% =====================================================================
% Fichier UNIQUE — mémoire PFE HAR (contenu fusionné)
% Bibliographie: references.bib | Figures: figures/
% =====================================================================

%% ============================================================
%%  Mémoire de Master — ESI Alger
%%  Structure calquée sur MSL_Walaa/main.tex (includes, ordre,
%%  bibliographie, annexes, listes). Contenu HAR + TikZ conservé.
%% ============================================================
\documentclass[a4paper,12pt]{report}

%% ── Packages (avant polyglossia / hyperref, comme MSL_Walaa) ─
\usepackage{amsmath,amssymb}
\usepackage[top=2.5cm,bottom=2.5cm,left=2.5cm,right=2.5cm]{geometry}
\usepackage{setspace}
\usepackage{indentfirst}
\usepackage{fancyhdr}
\usepackage{graphicx}
\graphicspath{{figures/}{../}{../har_project/results/figures/}}

\usepackage[nottoc,notlof,notlot]{tocbibind}
\usepackage{tocloft}
\usepackage{booktabs}
\usepackage{array}
\usepackage{tabularx}
\usepackage{xltabular}
\newcolumntype{Y}{>{\raggedright\arraybackslash}X}
\usepackage{longtable}
\usepackage{multirow}
\usepackage[table,xcdraw]{xcolor}
\usepackage{wrapfig}
\usepackage{float}
\usepackage[labelsep=space]{caption}
\usepackage{subcaption}
\usepackage{adjustbox}
\usepackage{placeins}
\captionsetup{list=true}
\captionsetup[subfigure]{list=false}
\usepackage{tikz}
\usepackage{pgfplots}
\usepackage{titlesec}
\usepackage{chngcntr}
\usepackage{listings}
\lstset{
  basicstyle=\small\ttfamily,
  breaklines=true,
  frame=single,
  framerule=0.4pt,
  columns=fullflexible,
  keepspaces=true,
}

\ifpdf
  \usepackage[pdftex]{hyperref}
\else
  \usepackage{hyperref}
\fi
\hypersetup{%
  colorlinks=true,
  linkcolor=black,
  citecolor=black,
  urlcolor=black}

\setcounter{tocdepth}{3}
\setcounter{secnumdepth}{3}

%% ── Couleurs HAR (contenu spécifique au PFE) ───────────────
\definecolor{harblue}{RGB}{52,120,196}
\definecolor{harlightblue}{RGB}{219,234,254}
\definecolor{hargreen}{RGB}{34,139,100}
\definecolor{harlightgreen}{RGB}{209,245,228}
\definecolor{hargray}{RGB}{248,249,250}
\definecolor{haraccent}{RGB}{234,88,12}
\definecolor{haryellow}{RGB}{254,243,199}
\definecolor{harpurple}{RGB}{109,40,217}
\definecolor{harlightpurple}{RGB}{237,233,254}
\definecolor{babyblue}{rgb}{0.79, 0.89, 0.95}

\pgfplotsset{compat=1.18}
\usetikzlibrary{shapes.geometric,arrows.meta,positioning,fit,
                backgrounds,calc,decorations.pathreplacing,babel}

%% ── Polyglossia + arabe (comme MSL_Walaa, fin de préambule) ─
\usepackage{fontspec}
\usepackage{polyglossia}
\setmainlanguage{french}
\setotherlanguage[calendar=gregorian]{arabic}
\IfFontExistsTF{Times New Roman}{%
  \setmainfont{Times New Roman}[%
    BoldFont = Times New Roman Bold,
    ItalicFont = Times New Roman Italic,
    BoldItalicFont = Times New Roman Bold Italic,
  ]%
}{%
  \setmainfont{TeX Gyre Termes}%
}
\newfontfamily\arabicfont[Script=Arabic,Scale=1.15]{Amiri}

%% ── Bibliographie (biblatex + biber — même bloc que MSL_Walaa) ─
\usepackage[citestyle=authoryear,bibstyle=authoryear,sorting=nyt,doi=false,isbn=false,eprint=false,maxcitenames=1,uniquelist=false,backend=bibtex]{biblatex}
\addbibresource{references.bib}

\AtEveryBibitem{
  \ifentrytype{misc}{
  }{
    \clearfield{url}
    \clearfield{urlyear}
  }
}

\renewcommand*{\nameyeardelim}{\addcomma\space}
\newrobustcmd*{\mkbibparensbold}[1]{\textbf{\mkbibparens{#1}}}

\DeclareCiteCommand{\parencite}[\mkbibparensbold]
  {\usebibmacro{prenote}}
  {%
    \usebibmacro{citeindex}%
    \bibhyperref{%
      \ifnameundef{labelname}
        {\printfield{label}}
        {\printnames{labelname}}%
      \nameyeardelim
      \printfield{year}%
    }%
  }
  {\multicitedelim\space}
  {\usebibmacro{postnote}}

\newcommand{\mychapter}[2]{%
  \chapter*{#2}
  \markboth{#2}{#2}
  \addcontentsline{toc}{chapter}{#2}
}

\usepackage[page,toc,titletoc,title]{appendix}
\addto\captionsfrench{%
  \renewcommand\appendixname{Annexe}
  \renewcommand\appendixpagename{Annexes}
  \renewcommand{\appendixtocname}{Annexes}
}
\appto\appendix{\addtocontents{toc}{\protect\setcounter{tocdepth}{0}}}
\appto\listoffigures{\addtocontents{lof}{\protect\setcounter{tocdepth}{1}}}
\appto\listoftables{\addtocontents{lot}{\protect\setcounter{tocdepth}{1}}}

\onehalfspacing
\setlength{\parskip}{6pt}
\setlength{\parindent}{1.25em}

\pagestyle{fancy}
\fancyhf{}
\fancyhead[L]{\small\leftmark}
\fancyfoot[R]{\thepage}
\renewcommand{\headrulewidth}{0.4pt}
\renewcommand{\chaptermark}[1]{\markboth{\chaptername\ \thechapter.\ #1}{}}

\fancypagestyle{plain}{%
  \fancyhf{}%
  \fancyfoot[R]{\thepage}%
  \renewcommand{\headrulewidth}{0pt}%
}

%% En-tête fixe pour la liste des sigles (évite « Liste des tableaux »)
\fancypagestyle{acronyms}{%
  \fancyhf{}%
  \fancyhead[L]{\small Liste des sigles et acronymes}%
  \fancyfoot[R]{\thepage}%
  \renewcommand{\headrulewidth}{0.4pt}%
}

\renewcommand{\figurename}{Figure}
\renewcommand{\tablename}{Tableau}
\DeclareCaptionFormat{chartefigure}{\textit{#1~#2}\quad #3}
\captionsetup[figure]{format=chartefigure,textfont=normalsize,list=true}
\captionsetup[table]{labelfont=bf,textfont=normalsize,labelsep=space,list=true}

\renewcommand{\cftfigleader}{\cftdotfill{\cftdotsep}}
\renewcommand{\cfttableader}{\cftdotfill{\cftdotsep}}
\renewcommand{\cftfigpresnum}{}
\renewcommand{\cftfigaftersnum}{}
\renewcommand{\cfttabpresnum}{}
\renewcommand{\cfttabaftersnum}{}

\titleformat{\chapter}[display]
  {\normalfont\large\bfseries}{\chaptertitlename\ \thechapter}{12pt}{\Large\bfseries}
\titleformat{\section}{\normalfont\bfseries}{\thesection}{0.8em}{}
\titleformat{\subsection}{\normalfont\bfseries}{\thesubsection}{0.8em}{}
\titlespacing*{\chapter}{0pt}{12pt}{12pt}
\titlespacing*{\section}{0pt}{12pt}{12pt}
\titlespacing*{\subsection}{0pt}{12pt}{12pt}
\titlespacing*{\subsubsection}{0pt}{12pt}{12pt}

\renewcommand{\thefigure}{\arabic{figure}}
\renewcommand{\thetable}{\Roman{table}}
\setlength{\extrarowheight}{3pt}
\renewcommand{\arraystretch}{1.15}

%% ── Figures : placeholders + images projet ─────────────────
\newcommand{\figslot}[4][2.6cm]{%
  \begin{center}%
  \fcolorbox{harblue!80!black}{harlightblue!50}{%
    \begin{minipage}{0.92\linewidth}%
      \centering\vspace{#1}%
      {\small\bfseries\textcolor{harblue}{Emplacement figure}}\\[0.55em]%
      {\footnotesize\path{figures/#2}}\\[0.45em]%
      {\footnotesize\itshape Source : #3}\\[0.25em]%
      {\footnotesize #4}%
      \vspace{#1}%
    \end{minipage}%
  }%
  \end{center}%
}

\newcommand{\insertfig}[4][]{%
  \IfFileExists{figures/#2}{%
    \adjustbox{max width=\linewidth,max height=0.82\textheight,center,#1}{%
      \includegraphics[keepaspectratio]{#2}}%
  }{%
    \IfFileExists{../har_project/results/figures/#2}{%
      \adjustbox{max width=\linewidth,max height=0.82\textheight,center,#1}{%
        \includegraphics[keepaspectratio]{#2}}%
    }{%
      \IfFileExists{../#2}{%
        \adjustbox{max width=\linewidth,max height=0.82\textheight,center,#1}{%
          \includegraphics[keepaspectratio]{#2}}%
      }{%
        \figslot{#2}{#3}{#4}%
      }%
    }%
  }%
}

\begin{document}

%% Page de garde (MSL_Walaa : \includepdf{00-Page-de-garde.pdf})

% ===== BEGIN input: 00-Page-de-garde.tex =====
\begin{titlepage}
  \enlargethispage{2\baselineskip}
  \setlength{\topskip}{0pt}
  \vspace*{-0.4cm}

  %% ── En-tête ──
  \begin{center}
    {\normalsize République Algérienne Démocratique et Populaire}\\[0.1cm]
    {\small\textarabic{الجمهورية الجزائرية الديمقراطية الشعبية}}\\[0.15cm]
    {\normalsize Ministère de l'Enseignement Supérieur et de la Recherche Scientifique}\\[0.1cm]
    {\small\textarabic{وزارة التعليم العالي و البحث العلمي}}
  \end{center}

  \vspace{0.15cm}
  \rule{\linewidth}{0.5pt}

  %% ── Logo + École ──
  \vspace{0.15cm}
  \noindent
  \begin{minipage}[t]{0.32\textwidth}
    \raggedright
    \includegraphics[width=3.0cm]{logo_esi.png}
  \end{minipage}%
  \hfill
  \begin{minipage}[t]{0.64\textwidth}
    \raggedleft
    {\small\textarabic{المدرسة الوطنية العليا للإعلام الآلي}}\\[0.05cm]
    {\small\textarabic{(المعهد الوطني للتكوين في الإعلام الآلي سابقا)}}\\[0.1cm]
    {\small École nationale Supérieure d'Informatique}\\[0.05cm]
    {\footnotesize ex. INI (Institut National de formation en Informatique)}
  \end{minipage}

  \vspace{0.55cm}

  %% ── Mémoire ──
  \begin{center}
    {\large\textbf{Mémoire de Master}}\\[0.2cm]
    {\normalsize Pour l'obtention du diplôme de Master en Informatique}\\[0.15cm]
    {\normalsize\textbf{Option : Systèmes d'information et technologies}}
  \end{center}

  \vspace{0.35cm}
  \rule{\linewidth}{0.5pt}

  %% ── Titre ──
  \vspace{0.3cm}
  \begin{center}
    {\Large\bfseries
      Une étude des approches d'IA\\[0.25cm]
      pour l'apprentissage adaptatif et continu\\[0.25cm]
      dans la reconnaissance d'activités humaines}
  \end{center}

  \vspace{0.3cm}
  \rule{\linewidth}{0.5pt}

  \vspace{0.55cm}

  %% ── Auteur / Encadrement ──
  \noindent
  \begin{minipage}[t]{0.48\textwidth}
    \raggedright
    \textit{Réalisée par :}\\[0.25cm]
    Mme. \textsc{Nedjam} Walaa
  \end{minipage}%
  \hfill
  \begin{minipage}[t]{0.48\textwidth}
    \raggedright
    \textit{Encadrée par :}\\[0.25cm]
    \textbf{Pr. Attal Ferhat} (LISSI, UPEC)\\[0.1cm]
    \textbf{Dr. Meziani Lila} (ESI)
  \end{minipage}

  \vspace{0.55cm}
  \noindent
  \textit{Organisme d'accueil :}\\[0.15cm]
  Laboratoire LISSI — Université Paris-Est Créteil (UPEC), France\\
  École Nationale Supérieure d'Informatique (ESI), Alger

  \vfill

  \begin{center}
    {\normalsize Promotion : \textbf{2025/2026}}
  \end{center}
  \vspace{0.2cm}
\end{titlepage}

% ===== END input: 00-Page-de-garde.tex =====


\pagenumbering{Roman}


% ===== BEGIN include: 01-Dedicace.tex =====
\clearpage
\mychapter{0}{Dédicace}

\begin{center}
\itshape
À mes parents, pour leur soutien et leurs encouragements tout au long de mon
parcours universitaire.\\[12pt]
À mes enseignants et encadrants, pour la qualité de leur accompagnement.\\[12pt]
À tous ceux qui m'ont aidée, de près ou de loin, à mener ce travail à bien.
\end{center}

\vfill
\begin{flushright}
\textit{--- Walaa Nedjam}
\end{flushright}

\clearpage

\clearpage
% ===== END include: 01-Dedicace.tex =====


% ===== BEGIN include: 02-Remerciements.tex =====
\clearpage
\mychapter{0}{Remerciements}

Tout d'abord, nous remercions Allah, le Tout-Puissant, de nous avoir accordé la
force, la patience et la persévérance nécessaires à l'accomplissement de ce travail.

Nous exprimons notre gratitude la plus sincère à notre encadrante Dr.\ Meziani Lila
de l'École Nationale Supérieure d'Informatique (ESI) d'Alger, et à notre promoteur
Pr.\ Attal Ferhat du Laboratoire LISSI de l'Université Paris-Est Créteil Val de Marne
(UPEC), pour la qualité de leur encadrement, la pertinence de leurs orientations et
leur accompagnement constant tout au long de ce projet.

Nous adressons également nos remerciements à l'ensemble du corps enseignant de
l'École Nationale Supérieure d'Informatique d'Alger pour la qualité de l'enseignement
dispensé durant notre formation.

À toutes celles et ceux qui ont, de près ou de loin, contribué à l'aboutissement
de ce projet, nous adressons notre plus sincère reconnaissance.

\clearpage
% ===== END include: 02-Remerciements.tex =====


% ===== BEGIN include: 03-Resume.tex =====
\clearpage
\mychapter{0}{Résumé}


Les smartphones et montres connectées embarquent des capteurs inertiels qui
permettent de reconnaître ce que fait une personne : marcher, monter des
escaliers, s'asseoir, etc. En pratique, un modèle entraîné une fois en
laboratoire se dégrade dès qu'un nouvel utilisateur arrive ou que le contexte
change. C'est l'oubli catastrophique. Peu de systèmes, en outre, cherchent à
anticiper une transition avant qu'elle ne soit complète.

Nous proposons ici une architecture HAR qui apprend de façon continue et qui
peut aussi prédire l'activité suivante à partir de fenêtres partielles. Le
cœur du système est un Transformer partagé, un tampon de rejeu guidé par
l'incertitude de prédiction (UWR), et une tête LSTM d'anticipation. Sur
HAPT et WISDM, UWR gagne environ 16,8 points de macro-F1 face au finetuning
naïf en scénario user-incrémental, et dépasse EWC et iCaRL sur nos protocoles.
L'anticipation multi-classe atteint un macro-F1 de validation d'environ
0,60--0,64 (contre $\approx$0,04 pour une baseline majorité). Un module
binaire de risque postural (fall-risk) monte à environ 0,88.

Le code (\texttt{har\_project}) contient aussi un pré-entraînement SSL léger,
un adaptateur de domaine affine, une calibration de confiance en ligne et un
petit bandit de seuil d'alerte. Le backbone fait environ 531\,K paramètres ;
le tampon de rejeu est limité à 2\,000 fenêtres. Le transfert HAPT vers WISDM
reste difficile (F1 $\approx$0,46 en zéro-shot). Nous n'avons pas de corpus
cuisine pour le scénario « préparation de repas » évoqué dans la fiche de
stage. Les annexes détaillent les données, les hyperparamètres et le code
pour reproduire les expériences.

\bigskip
\noindent\textbf{Mots clés :} Reconnaissance d'activités humaines, apprentissage
continu, capteurs inertiels, Transformer, Uncertainty-Weighted Replay, anticipation
d'activités.

\clearpage

\mychapter{0}{Abstract}


Inertial sensors in phones and watches make it possible to recognise daily
activities, but a model trained once on laboratory data usually degrades when
new users or contexts appear. Few systems try to anticipate a transition
before it is complete.

This thesis describes a continual HAR system built around a shared Transformer,
an uncertainty-weighted replay buffer (UWR), and an LSTM head that predicts the
next activity from partial windows. On HAPT and WISDM, UWR improves macro-F1 by
about 16.8 points over naive fine-tuning in the user-incremental setting and
outperforms EWC and iCaRL in our protocols. Next-activity anticipation reaches
validation macro-F1 around 0.60--0.64 (majority baseline $\approx$0.04). A
binary postural fall-risk head reaches about 0.88. Lighter modules cover SSL
pretraining, a domain adapter, online confidence calibration and a threshold
bandit. Cross-dataset transfer remains hard (HAPT$\rightarrow$WISDM F1
$\approx$0.46 zero-shot).

\bigskip
\noindent\textbf{Keywords:} Human activity recognition, continual learning, inertial
sensors, Transformer, Uncertainty-Weighted Replay, activity anticipation.

\clearpage

\chapter*{\hfill \textarabic{ملخص}}

\begin{otherlanguage}{arabic}
\noindent
\textarabic{%
يمكن لحساسات القصور الذاتي في الهواتف والساعات التعرف على الأنشطة اليومية،
لكن النموذج المدرّب مرة واحدة على بيانات مخبرية يتراجع أمام مستخدمين أو سياقات
جديدة، وقليل من الأنظمة يتوقع الانتقال قبل اكتماله.}

\bigskip
\noindent
\textarabic{%
يصف هذا المذكّر نظام HAR مستمراً يعتمد على Transformer مشترك، وإعادة تشغيل
مرجّحة بعدم اليقين (UWR)، ورأس LSTM لتوقع النشاط التالي. التجارب على HAPT
وWISDM تُظهر تحسناً واضحاً مقارنة بالضبط البسيط وEWC وiCaRL.}

\bigskip
\noindent\textbf{\textarabic{الكلمات المفتاحية:}}
\textarabic{%
التعرف على الأنشطة البشرية، التعلم المستمر، حساسات القصور الذاتي،
Transformer، UWR، توقع النشاط.}
\end{otherlanguage}

\addcontentsline{toc}{chapter}{ملخص}

% ===== END include: 03-Resume.tex =====


\renewcommand{\cftpartleader}{\cftdotfill{\cftdotsep}}
\renewcommand{\cftchapleader}{\cftdotfill{\cftdotsep}}
\tableofcontents
\clearpage
\listoffigures
\clearpage
\listoftables

\clearpage
\markboth{Liste des sigles et acronymes}{Liste des sigles et acronymes}

% ===== BEGIN include: 04-Acronymes.tex =====
\clearpage
\mychapter{0}{Liste des sigles et acronymes}

%% En-tête + pied de page dédiés (indépendants de \leftmark / list of tables)
\pagestyle{acronyms}
\thispagestyle{acronyms}

\vspace*{-0.6em}

\begingroup
\setlength{\tabcolsep}{5pt}
\setlength{\extrarowheight}{2pt}
\renewcommand{\arraystretch}{1.12}

\newcommand{\AcronymsTableHead}{%
    \hline
    \cellcolor{babyblue}\textbf{Sigle} &
    \cellcolor{babyblue}\textbf{Signification en anglais} &
    \cellcolor{babyblue}\textbf{Signification en français} \\
    \hline
}

%% Pleine largeur utile (éviter center + \textwidth qui rétrécit le tableau)

%% ── Page 1 ────────────────────────────────────────────────
\noindent
\begin{tabularx}{\textwidth}{|p{2.2cm}|>{\raggedright\arraybackslash}X|>{\raggedright\arraybackslash}X|}
    \AcronymsTableHead
%LIGNE =================================================
    \textbf{HAR} &
Human Activity Recognition &
Reconnaissance d'activités humaines \\
    \hline
%LIGNE =================================================
    \textbf{IMU} &
Inertial Measurement Unit &
Unité de mesure inertielle \\
    \hline
%LIGNE =================================================
    \textbf{CL} &
Continual Learning &
Apprentissage continu \\
    \hline
%LIGNE =================================================
    \textbf{AI} &
Artificial Intelligence (AI) &
Intelligence artificielle (IA) \\
    \hline
%LIGNE =================================================
    \textbf{ML} &
Machine Learning &
Apprentissage automatique \\
    \hline
%LIGNE =================================================
    \textbf{DL} &
Deep Learning &
Apprentissage profond \\
    \hline
%LIGNE =================================================
    \textbf{CNN} &
Convolutional Neural Network &
Réseau de neurones convolutionnel \\
    \hline
%LIGNE =================================================
    \textbf{LSTM} &
Long Short-Term Memory &
Mémoire à long et court terme \\
    \hline
%LIGNE =================================================
    \textbf{GRU} &
Gated Recurrent Unit &
Unité récurrente à portes \\
    \hline
%LIGNE =================================================
    \textbf{MLP} &
Multilayer Perceptron &
Perceptron multicouche \\
    \hline
%LIGNE =================================================
    \textbf{SVM} &
Support Vector Machine &
Machine à vecteurs de support \\
    \hline
%LIGNE =================================================
    \textbf{RF} &
Random Forest &
Forêt aléatoire \\
    \hline
%LIGNE =================================================
    \textbf{HMM} &
Hidden Markov Model &
Modèle de Markov caché \\
    \hline
%LIGNE =================================================
    \textbf{EWC} &
Elastic Weight Consolidation &
Consolidation élastique des poids \\
    \hline
%LIGNE =================================================
    \textbf{LwF} &
Learning without Forgetting &
Apprentissage sans oubli \\
    \hline
%LIGNE =================================================
    \textbf{ER} &
Experience Replay &
Rejeu d'expériences \\
    \hline
\end{tabularx}

\vspace{2em}

\clearpage
\thispagestyle{acronyms}

\vspace*{-0.6em}

%% ── Page 2 ────────────────────────────────────────────────
\noindent
\begin{tabularx}{\textwidth}{|p{2.2cm}|>{\raggedright\arraybackslash}X|>{\raggedright\arraybackslash}X|}
    \AcronymsTableHead
%LIGNE =================================================
    \textbf{iCaRL} &
Incremental Classifier and Representation Learning &
Classifieur et représentation incrémentaux \\
    \hline
%LIGNE =================================================
    \textbf{UWR} &
Uncertainty-Weighted Replay &
Rejeu pondéré par l'incertitude \\
    \hline
%LIGNE =================================================
    \textbf{MC} &
Monte Carlo &
Monte Carlo \\
    \hline
%LIGNE =================================================
    \textbf{SSL} &
Self-Supervised Learning &
Apprentissage auto-supervisé \\
    \hline
%LIGNE =================================================
    \textbf{CORAL} &
CORrelation ALignment &
Alignement de covariance \\
    \hline
%LIGNE =================================================
    \textbf{DA-BN} &
Domain Adaptive Batch Normalization &
Normalisation adaptative par lots \\
    \hline
%LIGNE =================================================
    \textbf{BWT} &
Backward Transfer &
Transfert rétroactif \\
    \hline
%LIGNE =================================================
    \textbf{FWT} &
Forward Transfer &
Transfert prospectif \\
    \hline
%LIGNE =================================================
    \textbf{AUC} &
Area Under the ROC Curve &
Aire sous la courbe ROC \\
    \hline
%LIGNE =================================================
    \textbf{ROC} &
Receiver Operating Characteristic &
Courbe caractéristique de fonctionnement du récepteur \\
    \hline
%LIGNE =================================================
    \textbf{KAN} &
Kolmogorov-Arnold Network &
Réseau de Kolmogorov-Arnold \\
    \hline
%LIGNE =================================================
    \textbf{OCL} &
Online Continual Learning &
Apprentissage continu en ligne \\
    \hline
%LIGNE =================================================
    \textbf{CLS} &
Classification Token &
Token de classification \\
    \hline
%LIGNE =================================================
    \textbf{FFN} &
Feed-Forward Network &
Réseau à propagation directe \\
    \hline
%LIGNE =================================================
    \textbf{IIR} &
Infinite Impulse Response &
Réponse impulsionnelle infinie \\
    \hline
%LIGNE =================================================
    \textbf{ADL} &
Activities of Daily Living &
Activités de la vie quotidienne \\
    \hline
%LIGNE =================================================
    \textbf{GAN} &
Generative Adversarial Network &
Réseau antagoniste génératif \\
    \hline
\end{tabularx}

\endgroup

\clearpage
\pagestyle{fancy}

\clearpage
% ===== END include: 04-Acronymes.tex =====


%%%%%%%%%%%%%%%%%%%%%%%%%%%%%%%%%%%%%%%%%%%%
%%% Content of the report and references %%%
%%%%%%%%%%%%%%%%%%%%%%%%%%%%%%%%%%%%%%%%%%%%

\cleardoublepage

\pagenumbering{arabic}


% ===== BEGIN include: 05-Introduction.tex =====
\clearpage
\chapter*{Introduction générale}
\addcontentsline{toc}{chapter}{Introduction générale}
\markboth{Introduction générale}{Introduction générale}
\label{chap:introduction}

Reconnaître une activité à partir d'un accéléromètre ou d'un gyroscope dans un
téléphone ou une montre est devenu courant, que ce soit pour le sport, la santé
ou l'aide à domicile. La manière habituelle de procéder reste simple : on
collecte des données en laboratoire, on entraîne un réseau une fois, puis on
déploie. Dès qu'un nouvel utilisateur arrive, que le capteur change de place
ou qu'une activité inédite apparaît, les performances chutent. Tout réapprendre
depuis zéro coûte cher ; un finetuning naïf fait souvent oublier ce qui
marchait. Et la plupart des systèmes se contentent de classer l'instant
présent, sans prévenir une transition qui s'annonce.

Ce mémoire part de ce constat. Nous construisons un système de HAR capable
d'apprendre au fil de l'eau quand de nouveaux utilisateurs ou de nouvelles
classes arrivent, et capable aussi de détecter tôt un changement d'activité.
L'implémentation vit dans le dépôt \texttt{har\_project} : un backbone
Transformer, un tampon de rejeu guidé par l'incertitude (UWR), une
anticipation multi-classe via une tête LSTM sur fenêtres partielles, un
adaptateur de domaine léger et un module de risque postural.

La première partie (chapitres~1 à~3) rappelle les bases de la HAR, les idées
de l'apprentissage continu et les travaux proches. La seconde (chapitres~4
à~6) décrit la conception, le code et les expériences, avec des comparaisons
à EWC, iCaRL et au finetuning séquentiel. La conclusion reprend les chiffres
et les limites.

\section{Problématique détaillée}

\subsection{Limites du paradigme batch}

Collecter un corpus, entraîner offline et figer les poids suppose une
distribution à peu près stable. En déploiement, ce n'est presque jamais le
cas. Les gens bougent différemment, le téléphone n'est pas toujours au même
endroit, de nouvelles activités apparaissent, et les signaux dérivent avec le
temps. Réentraîner périodiquement sur tout l'historique est coûteux, et les
données brutes ne peuvent pas toujours rester sur un serveur.

\subsection{Oubli catastrophique en pratique}

Quand on fine-tune sujet par sujet --- comme dans notre protocole à 44 tâches
--- l'accuracy tombe d'environ 91~\% (entraînement joint) à 62~\% sans
mécanisme anti-oubli (chapitre~6). EWC et iCaRL aident, mais ne ciblent pas
particulièrement les exemples près des frontières de décision. UWR oriente un
tampon de 2\,000 fenêtres vers celles où le modèle est le plus incertain.

\subsection{Anticipation vs classification instantanée}

Beaucoup de travaux HAR cherchent surtout une bonne accuracy sur la fenêtre
courante. Pour la sécurité (chute, sortie d'une zone sûre), il faudrait
parfois agir avant que le mouvement soit terminé. La littérature vidéo
\parencite{ryoo2011,gammulle2019} parle d'early recognition ; côté IMU, peu
d'articles mêlent anticipation et apprentissage continu. Dans
\texttt{har\_project}, on prédit la classe suivante à partir de $S=5$
fenêtres tronquées à $p\in\{25,50,75\}\%$, en ne gardant que les transitions.
On obtient un macro-F1 de validation d'environ 0,60--0,64 (baseline majorité
$\approx$0,04). Un score binaire de risque postural complète ce volet
(macro-F1 $\approx$0,88).

\section{Approche proposée}

Le modèle \texttt{HARContinualModel} s'appuie sur un encodeur Transformer
pré-entraîné (embeddings 128-D via un token CLS ; une initialisation SSL est
possible). L'apprentissage continu repose sur UWR, des prototypes de classes
et une perte contrastive, en scénarios user- et class-incremental. Autour de
ce noyau, on trouve une tête d'anticipation LSTM, un \texttt{DomainAdapter},
le fall-risk, une calibration de confiance en ligne et un petit bandit de
seuil. Chaque brique peut être testée seule ; les ablations du chapitre~6
s'en servent.

\section{Organisation du document}

Le chapitre~1 introduit le vocabulaire HAR et les pipelines classiques. Le
chapitre~2 formalise l'apprentissage continu. Le chapitre~3 compare les
travaux récents et pointe ce qui manque encore. Les chapitres~4 à~6
couvrent la conception, l'implémentation Python et l'évaluation. Les annexes
rassemblent datasets, hyperparamètres, extraits de code et fiches de lecture.

% ===== END include: 05-Introduction.tex =====



% ===== BEGIN include: part-1/06-Chapter1.tex =====
\clearpage
\part{État de l'art}

\chapter{Reconnaissance d'activités humaines, capteurs, pipelines et limites}
\label{ch:06-Chapter1}

La reconnaissance d'activités humaines (HAR) à partir de capteurs portés
accompagne aujourd'hui de nombreuses applications : santé mobile, aide à domicile,
sport, sécurité. Ce chapitre pose le vocabulaire de base pour la suite du mémoire :
définition de la HAR, capteurs utilisés, pipeline de traitement classique, passage
aux méthodes profondes, puis les limites qui nous poussent vers l'apprentissage
continu (chapitres~2 et~3) et notre contribution (partie~II).

\section{Définition et enjeux}

\subsection{Définition de la HAR}

En HAR « capteur », on cherche à inférer une activité physique (marche, assis,
escaliers, etc.) à partir de signaux mesurés sur le corps ou dans la poche
\parencite{bulling2014}. Les surveys récents rappellent que cette information sert
souvent de couche intermédiaire : une application en aval déclenche une alerte,
adapte une interface ou estime une charge d'effort \parencite{gu2021}.

On distingue classiquement la HAR vidéo de la HAR capteur. La seconde est
privilégiée quand la confidentialité et le déploiement embarqué comptent
\parencite{wang2019}. Un voisin proche est l'inférence précoce : reconnaître une
activité avant la fin de la séquence \parencite{ryoo2011}, ou anticiper la
transition suivante en vision \parencite{gammulle2019}. Ces travaux ne remplacent
pas la HAR IMU, mais expliquent pourquoi notre mémoire traite aussi
l'anticipation (chapitre~4).

\subsection{Applications}

Les activités couvrent la locomotion, les gestes du quotidien, la rééducation,
parfois la détection de chutes \parencite{gu2021,wang2019}. Selon l'usage, les
contraintes changent : latence, robustesse au bruit, respect de la vie privée.
Côté données, étiqueter des protocoles longs reste coûteux ; d'où l'intérêt du
pré-entraînement auto-supervisé ou contrastif
\parencite{yuan2024,tang2021}. Notre PFE vise explicitement l'adaptation
continue et l'anticipation, au-delà d'un classifieur statique entraîné une fois
for all.

\subsubsection{Santé et bien-être}

En santé mobile, la HAR sert à estimer le niveau d'activité physique (comptage
de pas, détection sédentarité), à suivre la rééducation post-opératoire ou à
détecter des chutes chez les personnes âgées \parencite{gu2021}. Les applications
cliniques exigent une faible latence et une robustesse aux variations de port du
capteur : un modèle entraîné sur des volontaires jeunes en laboratoire peut
dégrader fortement sur un patient âgé à domicile. D'où l'intérêt de
l'apprentissage continu utilisateur par utilisateur, thème central de notre
mémoire.

\subsubsection{Aide à domicile et sécurité}

Les systèmes d'aide à la vie autonome combinent souvent HAR et détection
d'événements critiques (chute, immobilité prolongée). La HAR fournit le contexte
(assis, debout, marche) ; un module spécialisé déclenche l'alerte. Notre
architecture prévoit explicitement un détecteur de chutes hybride (règles +
classifieur) alimenté par le même backbone IMU, même si l'évaluation actuelle
repose sur un proxy HAPT faute de MobiAct.

\subsubsection{Sport et coaching}

En sport, la HAR distingue types d'exercices, compte répétitions ou estime
l'intensité \parencite{wang2019}. Les contraintes diffèrent de la santé : signaux
plus intenses, mouvements non standardisés, besoin de personnalisation rapide
(nouvel athlète, nouvel exercice). Le class-incremental (nouvelles classes
d'activité) modélise partiellement ce cas — scénario évalué au chapitre~6.

\subsubsection{Industrie et ergonomie}

En environnement professionnel, la HAR peut repérer des postures à risque
(torsion, port de charge) ou valider le respect de gestes de sécurité
\parencite{gu2021}. Les capteurs portés remplacent parfois des caméras en zone
restreinte. L'adaptation à de nouveaux opérateurs ou postes de travail reprend
les mêmes enjeux d'oubli et de domain shift que la santé connectée.

\subsubsection{Synthèse des exigences transverses}

Quel que soit le domaine, trois exigences reviennent : (i)~généraliser à de
nouveaux utilisateurs sans réentraînement complet ; (ii)~intégrer de nouvelles
activités ou capteurs au fil du temps ; (iii)~anticiper les changements d'état
pour agir avant qu'un incident ne se produise. Les chapitres~2 à~6 structurent
notre réponse à ces trois points.

\section{Capteurs, modalités et déploiement}

\subsection{Les capteurs inertiels (IMU)}

Une IMU combine au minimum un accéléromètre et un gyroscope ; le magnétomètre
est souvent présent sur smartphone \parencite{bulling2014,gu2021}. L'accéléromètre
mesure l'accélération (souvent en g ou m/s²) ; le gyroscope apporte la dynamique
de rotation, utile pour discriminer des activités proches
\parencite{ordonyez2016}. En pratique, deux datasets HAPT et WISDM n'utilisent
pas les mêmes fréquences, unités ou conventions de gravité : toute chaîne
multi-sources doit homogénéiser ces détails \parencite{amrani2025}.

\begin{figure}[H]
  \centering
  \insertfig[max width=0.42\linewidth]{bulling_fig2_imu_setup.png}{Schéma IMU triaxial}{Capteur IMU}
  \caption[Capteur IMU triaxial]{Capteur IMU triaxial — axes accéléromètre et gyroscope.}
  \label{fig:imu_ch1}
\end{figure}

\subsection{Wearables versus smartphones}

Smartphones et montres restent les plateformes les plus répandues pour collecter
des IMU « au contact » de l'utilisateur \parencite{bulling2014,wang2019}. La HAR
capteur est en général mieux acceptée que la vidéo dans des contextes sensibles
\parencite{wang2019}. Sur téléphone, l'accéléromètre seul suffit parfois pour des
activités grossières ; le gyroscope améliore les cas ambigus
\parencite{ferrari2020}.

Les montres connectées offrent un placement poignet stable mais des amplitudes
de signal différentes de la ceinture (HAPT) ou de la poche (WISDM). Un modèle
entraîné sur ceinture peut perdre 20--30 points de F1 sur poignet sans
adaptation — phénomène documenté dans les comparaisons multi-datasets
\parencite{amrani2025}. C'est une motivation directe de nos expériences CORAL
et du protocole 44 sujets mélangeant placements.

\subsection{Fréquence d'échantillonnage et filtrage}

La fréquence typique varie de 20~Hz (WISDM) à 100~Hz (PAMAP2, MobiAct). Notre
pipeline unifie à 50~Hz par rééchantillonnage Fourier \parencite{amrani2025}.
La composante gravitationnelle ($\approx 9,81$~m/s²) biaise l'accéléromètre
selon l'orientation ; un filtre passe-bas 0,5~Hz estime et soustrait cette
composante avant fenêtrage (chapitre~5). Sans cette étape, « debout » et
« assis » peuvent se confondre sur la seule magnitude accélérométrique.

\subsection{Qualité des annotations}

En laboratoire (HAPT), les activités sont scriptées et étiquetées
échantillon-par-échantillon. En conditions libres (WISDM), les labels sont
déclaratifs et bruités : marche vs jogging flou, transitions non marquées.
Cette asymétrie explique pourquoi nous réservons WISDM surtout au test
cross-domain et pas au pré-train principal initial, avant la fusion étendue
du projet \texttt{har\_project}.

\subsection{Embarqué, cloud, et contraintes réelles}

Bulling et al.\ \parencite{bulling2014} opposent traitement hors ligne (acquisition
puis analyse) et traitement en ligne (décision pendant l'enregistrement). En
déploiement, on sépare souvent un entraînement lourd (serveur ou poste de travail)
d'une inférence ou d'un fine-tuning léger sur appareil
\parencite{amrani2025,mazankiewicz2020}. Les gros pré-entraînements SSL
\parencite{yuan2024} renforcent cette logique hybride : ressources importantes en
amont, adaptation ciblée ensuite.

\subsection{Multi-modalité}

Fusionner accéléromètre et gyroscope est courant lorsqu'une seule modalité prête
à confusion \parencite{ordonyez2016}. Les surveys élargissent parfois le spectre
(WiFi, RFID, etc.) \parencite{gu2021}. Côté modèle, les Transformers sur séries
\parencite{dirgova2022} offrent une alternative aux chaînes CNN/LSTM — piste retenue
dans notre backbone.

\section{Pipeline classique de la HAR}

\begin{figure}[H]
  \centering
  \insertfig[max width=\linewidth,max height=0.22\textheight]{bulling2014_arc_pipeline.png}{Bulling et al. (2014) — Fig. 1}{Pipeline ARC}
  \caption[Pipeline classique HAR — chaîne ARC \parencite{bulling2014}]{Pipeline classique de la HAR — chaîne ARC
  \parencite{bulling2014,wang2019,gu2021}.}
  \label{fig:pipeline_ch1}
\end{figure}

\subsection{Acquisition}

On enregistre des séries brutes : fréquence d'échantillonnage, placement du capteur,
protocole d'activités. En conditions réelles, la variabilité inter-sujets et la
dérive du capteur sont la norme, pas l'exception \parencite{bulling2014}. Fusionner
plusieurs bases publiques impose en plus d'aligner labels, unités et documentation
\parencite{amrani2025}.

\subsection{Segmentation}

Les activités quasi périodiques (marche, course) se découpent en fenêtres glissantes,
souvent avec recouvrement \parencite{bulling2014}. Les modèles profonds multimodaux
utilisent le même principe sur signaux bruts \parencite{ordonyez2016}. Dans notre
projet, une fenêtre de 3~s à 50~Hz (150 échantillons, stride 50~\%) suit cette
convention \parencite{amrani2025}.

\begin{figure}[H]
  \centering
  \insertfig[max width=0.75\linewidth]{bulling_fig5_windowing.png}{Fenêtrage glissant}{Fenêtre glissante}
  \caption[Fenêtrage glissant sur signaux IMU]{Principe de fenêtrage glissant avec
  chevauchement \parencite{bulling2014}.}
  \label{fig:windowing_ch1}
\end{figure}

\subsection{Features et classification}

Historiquement, chaque fenêtre était décrite par des statistiques ou des
descripteurs fréquentiels, puis classée (SVM, forêts aléatoires, k-NN, HMM)
\parencite{bulling2014,gu2021,wang2019}. Ces pipelines restent utiles comme baseline
légère ou quand les données sont rares. Leur limite commune en HAR : la qualité
des descripteurs fixe un plafond, surtout quand l'utilisateur ou le capteur change.

\subsection{Chaîne ARC}

Le tutorat de Bulling formalise la chaîne \emph{Acquisition -- Segmentation --
Features -- Classification -- Fusion -- Évaluation} \parencite{bulling2014}. Les
surveys deep learning \parencite{gu2021,wang2019} la retrouvent avec des blocs
neuronaux à la place de l'extraction manuelle.

\section{Méthodes profondes modernes}

\subsection{Du feature engineering à la représentation apprise}

Les réseaux profonds apprennent directement des représentations à partir des
fenêtres, ce qui réduit la dépendance aux descripteurs artisanaux
\parencite{gu2021,wang2019}. DeepConvLSTM \parencite{ordonyez2016} illustre le
couple CNN+LSTM sur signaux portés multimodaux.

\subsection{Attention et Transformers}

\begin{figure}[H]
  \centering
  \insertfig[max width=0.88\linewidth]{dirgova2022_transformer.png}{Dirgová et al. (2022)}{IMUTransformerEncoder}
  \caption[Architecture Transformer IMU \parencite{dirgova2022}]{Architecture IMUTransformerEncoder
  \parencite{dirgova2022,vaswani2017}.}
  \label{fig:transformer_ch1}
\end{figure}

Un Transformer calcule, pour chaque pas de temps, une combinaison pondérée de
toute la fenêtre via l'attention \parencite{vaswani2017,dirgova2022}. Avec
$Q, K, V$ les projections requête, clé et valeur :
\[
\mathrm{Attention}(Q,K,V) = \mathrm{softmax}\!\left(\frac{QK^\top}{\sqrt{d_k}}\right) V.
\]
En HAR, $Q_t$ encode « ce que cherche » l'instant $t$, $K_s$ mesure la
compatibilité avec l'instant $s$, et $V_s$ porte l'information agrégée. Les
encodages positionnels réintroduisent l'ordre temporel. Le coût quadratique en
longueur de séquence limite l'embarqué ; notre backbone reste volontairement
compact (4 blocs, $d=128$).

\subsection{SSL, intégration de données, personnalisation}

Yuan et al.\ \parencite{yuan2024} montrent l'intérêt d'un très grand
pré-entraînement non supervisé sur wearables, puis d'un fine-tuning aval.
\textcite{amrani2025} insistent plutôt sur l'homogénéisation multi-bases (fréquence,
unités, labels) avant fusion. Côté utilisateur, la personnalisation par sélection
de données \parencite{ferrari2020} ou par batch norm adaptive en ligne
\parencite{mazankiewicz2020} complète le tableau sans résoudre seul l'oubli
séquentiel.

\begin{figure}[H]
  \centering
  \insertfig[max width=0.82\linewidth]{gu_fig2_dl_har_overview.png}{Gu et al. (2021)}{Panorama deep learning HAR}
  \caption[Panorama des approches deep learning HAR \parencite{gu2021}]{Panorama
  des familles d'architectures profondes pour la HAR \parencite{gu2021,wang2019}.}
  \label{fig:dl_overview_ch1}
\end{figure}

\section{Méthodes traditionnelles (hors deep learning)}

Avant les réseaux profonds, la HAR capteur suivait presque toujours le même
schéma : fenêtre glissante, extraction de descripteurs, classifieur
\parencite{bulling2014}. Ces pipelines restent utiles comme baseline légère ou
sur microcontrôleur quand la mémoire est très contrainte \parencite{wang2019}.

\subsection{Machines à vecteurs de support (SVM)}

Une SVM cherche une frontière (éventuellement via noyau RBF) qui sépare les
vecteurs de features par classe. En multi-classe, on combine des classifieurs
binaires one-vs-one ou one-vs-rest. La performance dépend entièrement des
descripteurs : moyenne, écart-type, énergie, FFT, etc. Dès qu'un nouvel
utilisateur ou un capteur décalé modifie la distribution des features, la marge
s'effondre sans réentraînement complet.

\subsection{Forêts aléatoires}

Une forêt aléatoire vote entre des arbres entraînés sur des sous-échantillons
bootstrap. Elle tolère mieux le bruit et les variables hétérogènes qu'une SVM
mal calibrée, mais ne voit la dynamique fine à l'intérieur d'une fenêtre que si
les features la rendent explicite. En pratique, elle sert souvent de référence
rapide avant de passer au deep learning \parencite{ferrari2020}.

\subsection{Modèles de Markov cachés (HMM)}

Un HMM modélise une séquence comme une chaîne d'états latents (activités ou
sous-phases) avec des lois d'émission gaussiennes. Viterbi et Forward-Backward
permettent l'inférence. L'hypothèse markovienne limite la modélisation des
durées d'activité ; les HMM restent pédagogiques pour comprendre la dimension
temporelle, mais rarement compétitifs face aux LSTM ou Transformers sur IMU
brut \parencite{gu2021}.

\section{Protocoles d'évaluation en HAR}

\subsection{Split train / test}

Deux protocoles dominent. Le \emph{random split} mélange fenêtres de tous les
sujets : scores optimistes, peu réalistes en déploiement
\parencite{wang2019}. Le \emph{leave-one-subject-out} (LOSO) retient un sujet
entier pour le test : c'est le standard pour mesurer la généralisation inter-
individus. Notre pré-entraînement HAPT utilise 80/20 par sujet (24 train, 6 val).

\subsection{Métriques}

L'accuracy seule trompe sur jeux déséquilibrés (transitions rares en HAPT).
Le macro-F1 moyenne les F1 par classe ; le micro-F1 agrège les comptages
\parencite{schiemer2023}. En CL, on ajoute forgetting et backward transfer
(chapitre~2). Pour l'anticipation, précision et rappel sur la classe
« transition » comptent plus que l'accuracy globale.

\begin{figure}[H]
  \centering
  \insertfig[max width=0.75\linewidth]{bulling_fig8_sensor_placement.png}{Placements capteur}{Placement smartphone}
  \caption[Placements du capteur \parencite{bulling2014}]{Variations de placement
  du smartphone — ceinture, poignet, poche \parencite{bulling2014,ferrari2020}.}
  \label{fig:placement_ch1}
\end{figure}

\section{Jeux de données publics fréquents en HAR}

Au-delà de HAPT et WISDM utilisés dans ce mémoire, la littérature mobilise
régulièrement les jeux suivants \parencite{gu2021,wang2019,amrani2025} :

\begin{table}[H]
\centering
\small
\begin{tabularx}{\textwidth}{|Y|c|c|Y|}
\hline
\textbf{Dataset} & \textbf{Sujets} & \textbf{Capteurs} & \textbf{Usage typique} \\
\hline
UCI HAR (sans transitions) & 30 & Smartphone ceinture & Baseline smartphone \\
Opportunity & 4 & Corps + objets & HAR ambient, multimodal \\
PAMAP2 & 9 & 3 IMU corps & Activités + sport \\
DSADS & 8 & 5 IMU corps & Daily + sport \\
MobiAct & 61 & Smartphone & Chutes réelles \\
\hline
\end{tabularx}
\caption{Jeux HAR publics cités dans le corpus (chapitre~3)}
\label{tab:public_datasets_ch1}
\end{table}

OCL-HAR \parencite{schiemer2023} et LAPNet \parencite{adaimi2022} évaluent sur
PAMAP2 et Opportunity ; Amrani \parencite{amrani2025} fusionne huit bases. Notre
choix HAPT+WISDM+PAMAP2 couvre ceinture, poche et capteurs corps — triangle
représentatif des domain shifts rencontrés en déploiement consumer vs clinique.

\section{Limites et passage à l'apprentissage continu}

\subsection{Données et généralisation}

Les jeux publics restent des protocoles contrôlés ; généraliser à un nouvel
utilisateur ou un nouveau capteur est difficile \parencite{wang2019,gu2021}.
L'hétérogénéité entre datasets n'est pas un détail technique : elle reflète des
différences d'appareil et de protocole \parencite{amrani2025}. Amrani et al.\
montrent que fusionner huit bases sans homogénéisation préalable creuse l'écart
de F1 entre modèle global et modèles locaux ; d'où leur pipeline d'alignement
\parencite{amrani2025}.

\begin{figure}[H]
  \centering
  \insertfig[max width=0.78\linewidth]{amrani2025_label_ambiguity.png}{Amrani et al. (2025)}{Ambiguïté des labels}
  \caption[Homogénéisation multi-datasets \parencite{amrani2025}]{Ambiguïté et
  homogénéisation des labels entre jeux HAR \parencite{amrani2025}.}
  \label{fig:amrani_labels_ch1}
\end{figure}

\subsection{Variabilité inter-sujets et personnalisation}

Deux utilisateurs exécutant « marche » produisent des signaux différents :
morphologie, cadence, placement du téléphone \parencite{ferrari2020}. La
personnalisation par fine-tuning ou batch norm adaptive aide, mais ne traite
pas l'arrivée séquentielle de dizaines d'utilisateurs avec mémoire bornée
\parencite{mazankiewicz2020,amrani2025}.

\subsection{Oubli catastrophique}

Quand les tâches ou les utilisateurs arrivent en séquence, réentraîner naïvement
efface souvent les acquis précédents \parencite{kirkpatrick2017,parisi2019}. Le
cadre de \textcite{vdven2022} (task-, domain- et class-incremental) aide à
préciser \emph{ce qui change} dans le flux. EWC \parencite{kirkpatrick2017} et
iCaRL \parencite{rebuffi2017} sont des réponses fondateurs, reprises ensuite en
HAR (OCL-HAR \parencite{schiemer2023}, LAPNet \parencite{adaimi2022}, iKAN
\parencite{liu2024}).

\subsection{Anticipation}

La reconnaissance « instantanée » ne suffit pas si l'on veut alerter \emph{avant}
une chute ou une transition. Les références vidéo \parencite{ryoo2011,gammulle2019}
posent le problème ; notre mémoire l'implémente en anticipation multi-classe
(LSTM, fenêtres partielles).
sur IMU (partie~II).

\section{Conclusion}

Ce chapitre a posé la HAR capteur, ses pipelines et le passage aux modèles
profonds. Les verrous récurrents — données limitées, décalage de domaine, oubli
séquentiel, absence d'anticipation — motivent le chapitre~2 (concepts CL) et le
chapitre~3 (revue ciblée). La partie~II propose une réponse concrète avec
\texttt{har\_project}.

\clearpage
% ===== END include: part-1/06-Chapter1.tex =====


% ===== BEGIN include: part-1/07-Chapter2.tex =====
\clearpage
\chapter{Apprentissage continu et adaptation dynamique}
\label{ch:07-Chapter2}

Dans ce chapitre, nous fixons le vocabulaire et les outils d'évaluation de
l'apprentissage continu (CL), puis nous les relions au HAR sur capteurs portés.
L'anticipation d'activité est introduite comme extension naturelle de la
modélisation temporelle. Un fil rouge ressort de la littérature : le CL est
bien documenté en vision, mais beaucoup moins standardisé pour la HAR en
conditions réelles (nouveaux utilisateurs, nouvelles classes, flux continu).
C'est précisément ce décalage qui motive la revue du chapitre~3 et notre
contribution aux chapitres~4 à~6.

\subsection{Organisation du chapitre}

Nous commençons par les définitions (§2.1), l'oubli catastrophique et EWC (§2.2),
puis la taxonomie des méthodes (§2.3) et des scénarios (§2.4). Les métriques
CL (§2.5) sont essentielles pour lire le chapitre~6. La section~2.6 instancie
cette taxonomie pour la HAR ; la conclusion relie au chapitre~3.

\section{Apprentissage continu : concepts}

\subsection{Online learning}

En apprentissage en ligne, le modèle reçoit des observations au fil du temps
(mini-lots ou fenêtres) et met à jour ses paramètres sans repasser sur tout
l'historique. Le défi principal est double : suivre une distribution qui
dérive (\emph{concept drift}) et garder un coût calcul acceptable. Sur
smartphone ou montre, cela correspond à un utilisateur qui produit des signaux
IMU en continu : on préfère alors des mises à jour légères plutôt qu'un
réentraînement complet à chaque changement de contexte
\parencite{mazankiewicz2020,amrani2025}.

\subsection{Incremental learning}

L'apprentissage incrémental décrit les situations où l'information arrive par
\emph{blocs} : une nouvelle tâche, un nouvel utilisateur, une nouvelle classe
ou un nouveau capteur. \textcite{vdven2022} classent ces scénarios selon ce
qui change dans le flux et selon que la frontière de tâche est connue ou non au
test — ce qui conditionne à la fois le protocole d'évaluation et le choix de
méthode (replay, régularisation, isolation de paramètres, etc.). En HAR,
l'arrivée séquentielle d'utilisateurs ou de conditions de port est modélisée
ainsi dans les travaux d'online continual learning, par exemple
\parencite{schiemer2023}.

\subsection{Lifelong learning}

Le \emph{lifelong learning} insiste sur la durée : le système doit accumuler
des compétences utiles pour la suite, sans effacer les acquis précédents.
\textcite{parisi2019} résument le problème comme un équilibre plasticité /
stabilité, avec plusieurs familles de réponses (régularisation, replay,
architectures dynamiques).

\subsection{Continual learning}

Le terme \emph{continual learning} est souvent utilisé comme fourre-tout
couvrant l'incrémental et le lifelong, avec un focus explicite sur
l'oubli catastrophique et la rétention des performances passées
\parencite{delange2021,parisi2019}. Les surveys rappellent aussi que le CL en
classification n'est pas une seule boîte : comparer deux articles sans
préciser le scénario (task-incremental, class-incremental, domain-incremental)
peut conduire à des conclusions trompeuses \parencite{delange2021,vdven2022}.

\section{Problème clé : l'oubli catastrophique}

Quand on entraîne un réseau profond sur une nouvelle distribution, les poids
s'ajustent fortement : la performance sur les anciennes données chute
souvent brutalement. C'est l'oubli catastrophique
\parencite{kirkpatrick2017,delange2021} — une interférence entre objectifs
successifs dans le même espace de paramètres, distincte d'une simple perte de
généralisation.

\textcite{kirkpatrick2017} proposent l'Elastic Weight Consolidation (EWC) :
pénaliser le déplacement des poids jugés importants pour les tâches passées,
via une approximation diagonale de la matrice de Fisher. La
figure~\ref{fig:ewc_ch2} schématise l'idée : un finetuning naïf « efface » la
tâche~A ; une régularisation trop forte bloque la tâche~B ; EWC cherche un
compromis.

\begin{figure}[H]
  \centering
  \insertfig{kirkpatrick2017_ewc_paramspace.png}{Kirkpatrick et al. (2017), PNAS — Fig. 1}{EWC vs finetuning naïf}
  \caption[Oubli catastrophique et EWC \parencite{kirkpatrick2017}]{Oubli catastrophique et régularisation EWC dans l'espace
  des paramètres \parencite{kirkpatrick2017,delange2021}.}
  \label{fig:ewc_ch2}
\end{figure}

Du côté class-incremental, iCaRL \parencite{rebuffi2017} combine exemplaires
mémorisés, distillation et classificateur incrémental. Cette ligne reste une
référence fréquente, y compris dans les protocoles HAR continus
\parencite{delange2021,schiemer2023}.

\section{Taxonomie des approches}

\begin{figure}[H]
  \centering
  \insertfig{delange2021_taxonomy.png}{De Lange et al. (2021) — Fig. 1}{Taxonomie CL}
  \caption[Taxonomie de l'apprentissage continu \parencite{delange2021}]{Taxonomie de l'apprentissage continu
  \parencite{delange2021,parisi2019,vdven2022}.}
  \label{fig:cl_taxonomy_ch2}
\end{figure}

La figure~\ref{fig:cl_taxonomy_ch2} et le découpage ci-dessous reprennent les
grandes familles des surveys \parencite{delange2021,parisi2019,vdven2022}, en
les reliant aux travaux HAR retenus dans notre corpus.

\subsection{Régularisation (\emph{regularization-based})}

L'idée est simple : empêcher le réseau de s'éloigner trop de la solution
apprise avant. EWC \parencite{kirkpatrick2017} ajoute une pénalité quadratique
sur les poids, pondérée par l'importance Fisher. LwF (\emph{Learning without
Forgetting}) stabilise les sorties anciennes par distillation lors de
l'apprentissage d'un nouvel objectif — famille discutée dans
\parencite{delange2021} et reprise sous forme de distillation dans OCL-HAR
\parencite{schiemer2023}.

\subsection{Replay et mémoires}

Deux variantes dominent. Le \emph{rehearsal} garde un petit tampon d'exemples
passés et les mélange aux nouvelles données : c'est le cœur d'iCaRL
\parencite{rebuffi2017} et de nombreuses variantes d'experience replay. Le
\emph{generative replay} synthétise des pseudo-échantillons ; les surveys
\parencite{delange2021} rappellent le compromis fidélité / coût.

En HAR, \textcite{schiemer2023} combinent distillation et rehearsal dans un
cadre online avec décalage lié à l'arrivée d'utilisateurs.
\textcite{adaimi2022} proposent LAPNet-HAR : prototypes, replay et contraste
dans un cadre \emph{task-free} plus proche d'un déploiement long terme.

\subsubsection{Replay uniforme vs priorisé}

Le replay uniforme ($P(x)=1/|M|$) est la baseline la plus simple : chaque
exemple mémorisé a la même chance d'être rejoué. \emph{Prioritized Experience
Replay} (Schaul et al., 2016) pondère par la TD-error en RL ; en CL supervisé,
notre UWR transpose l'idée en pondérant par l'entropie de prédiction
\parencite{gal2016}. Intuition : les exemples près des frontières de décision
deviennent incertains avant d'être mal classifiés ; les rejouer en priorité
retarde l'oubli. Le tampon de 2\,000 fenêtres ($\approx$ 3,5~\% du corpus
fusionné) force cette sélection : on ne peut pas tout garder, il faut optimiser
l'information mémorisée.

\subsubsection{Prototypes vs exemplaires bruts}

iCaRL stocke des exemplaires bruts par classe ; LAPNet et notre système
maintiennent des prototypes (moyennes d'embeddings). Avantage : empreinte réduite
(un vecteur 128-D par classe vs fenêtres complètes 150$\times$6). Inconvénient :
perte d'information fine sur les transitions courtes — d'où le maintien du tampon
UWR \emph{en plus} des prototypes, pas à la place.


Plutôt que de tout partager, on isole : sous-réseaux par tâche, masques qui
figent certaines connexions, compression de poids. \textcite{delange2021}
mentionnent GEM, A-GEM, PackNet ou HAT comme réponses structurelles à
l'interférence. Pour des jeux HAR hétérogènes, iKAN \parencite{liu2024}
illustre une autre piste : branches par tâche et redistribution de
caractéristiques pour limiter à la fois l'oubli et l'incompatibilité entre
datasets.

\subsection{Approches prototypiques}

Les \emph{prototypical networks} représentent chaque classe par un centroïde
dans l'espace d'embeddings ; la décision se fait par distance au prototype le
plus proche. LAPNet-HAR \parencite{adaimi2022} s'inscrit dans cette logique,
avec mise à jour continue des prototypes et replay pour éviter qu'ils ne
« dérivent » quand le backbone évolue.

\subsection{Pré-entraînement et fine-tuning}

Une part croissante du corpus HAR passe par un gros pré-entraînement (souvent
SSL), puis une adaptation légère : Yuan et al.\ \parencite{yuan2024} sur
700\,000 person-days de wearables, contrastif santé \parencite{tang2021},
fine-tuning après homogénéisation multi-bases \parencite{amrani2025},
personnalisation par batch norm adaptive \parencite{mazankiewicz2020}. Les
surveys HAR \parencite{wang2019,gu2021} rappellent en parallèle que
l'hétérogénéité des capteurs reste un frein même avec de bonnes
représentations.

\section{Adaptation au contexte HAR}

\subsection{Adaptation de domaine}

Un même algorithme de marche ne se comporte pas pareil selon le capteur, sa
position (ceinture, poignet, poche) ou l'utilisateur. \textcite{gu2021,wang2019}
passent en revue les approches profondes et ces sources de décalage.
\textcite{mazankiewicz2020} proposent une personnalisation incrémentielle via
batch normalization \emph{domain-adaptive}, pensée pour un suivi en temps réel
sans réentraînement complet.

\subsection{Personnalisation}

La personnalisation vise à coller au geste individuel, au prix d'un risque de
sur-apprentissage et d'un coût de mise à jour. Dans notre corpus, elle se
croise souvent avec l'homogénéisation de données et le fine-tuning ciblé
\parencite{amrani2025,mazankiewicz2020}.

\subsection{Few-shot learning}

Quand les labels par utilisateur ou par classe sont rares, le few-shot ou le
métapprentissage devient pertinent. Les travaux retenus y arrivent surtout via
des représentations génériques (SSL, contrastif) plutôt que via un
meta-learner classique \parencite{yuan2024,tang2021}.

\subsection{Calibration en ligne}

On regroupe ici la mise à jour des statistiques de normalisation
\parencite{mazankiewicz2020}, le recalibrage des scores de confiance (température,
Platt scaling — détaillé au chapitre~5 si utilisé) et la stabilisation par
distillation \parencite{schiemer2023}.

\section{Anticipation d'activités}

Au-delà de « quelle activité maintenant ? », l'anticipation demande « que va-t-il
se passer ensuite ? » ou « peut-on reconnaître tôt une activité à partir d'une
fenêtre incomplète ? ». Ryoo \parencite{ryoo2011} pose des bases en vidéo ;
Gammulle et al.\ \parencite{gammulle2019} montrent des architectures jointes
en vision. Côté capteurs, Ordóñez et Roggen \parencite{ordonyez2016} ancrent
le séquentiel CNN+LSTM ; les transformeurs sur séries
\parencite{dirgova2022} ouvrent la même question avec d'autres blocs. Dans
notre projet, nous implémentons l'anticipation en prédiction multi-classe de
la classe suivante (tête LSTM, fenêtres partielles) ; les résultats du
chapitre~6 en précisent les performances.

Le lien avec l'apprentissage par renforcement reste surtout conceptuel pour le
HAR batch sur IMU ; les modèles séquentiels supervisés restent la voie la plus
fréquente dans la littérature capteur.

\subsection{Incertitude : MC Dropout (théorie) vs entropie (implémentation)}

Pour prioriser les exemples difficiles en replay (UWR, chapitre~4), nous
Gal et Ghahramani \parencite{gal2016} proposent le Monte Carlo Dropout pour
estimer l'incertitude épistémique. Dans \texttt{har\_project}, l'UWR utilise
une approximation plus légère : entropie de la softmax en une passe
(\texttt{uncertainty\_replay.py}), recalculée périodiquement sur le tampon.
Le MC Dropout reste la référence conceptuelle, pas le mécanisme exécuté.

\section{Métriques d'évaluation}

\subsection{Accuracy}

Pour $C$ classes et $N$ exemples de test :
\[
\mathrm{Acc} = \frac{1}{N}\sum_{i=1}^{N} \mathbf{1}\{\hat{y}_i = y_i\}.
\]
En CL, on trace souvent l'accuracy par tâche au fil des sessions
\parencite{delange2021}.

\subsection{F1-score}

Pour une classe $c$, avec $TP_c$, $FP_c$, $FN_c$ :
\[
P_c = \frac{TP_c}{TP_c + FP_c}, \quad
R_c = \frac{TP_c}{TP_c + FN_c}, \quad
F1_c = \frac{2 P_c R_c}{P_c + R_c}.
\]
Le \textbf{macro-F1} moyenne les $F1_c$ sur les classes ; le \textbf{micro-F1}
agrège d'abord les comptages globaux. Sur jeux déséquilibrés — fréquent en
HAR — \textcite{schiemer2023} argumentent pour les deux ; \textcite{amrani2025}
rapportent des F1 après fusion multi-datasets.

\subsection{Mesure d'oubli (\emph{forgetting})}

Soit $a_{i,j}$ l'accuracy sur la tâche $j$ après apprentissage jusqu'à la
tâche $i$ ($i \geq j$). L'oubli de la tâche $j$ après $T$ tâches :
\[
f_j = \max_{i \in \{j,\ldots,T-1\}} a_{i,j} - a_{T,j}, \qquad
\bar{f} = \frac{1}{T-1}\sum_{j=1}^{T-1} f_j.
\]
\textcite{delange2021} s'appuient sur des métriques de ce type. En online HAR,
\textcite{schiemer2023} proposent aussi un score d'oubli par classe à l'étape
$t$ :
\[
F_t = \frac{1}{|C|}\sum_{c \in C}
\left(\max_{l \in \{1,\ldots,t-1\}} a_l^{(c)} - a_t^{(c)}\right).
\]

\subsection{Backward transfer (BWT) et forward transfer (FWT)}

\textcite{vdven2022} distinguent le transfert vers le passé (BWT) et vers
l'avenir (FWT). Convention fréquente, avec $\bar{b}_j$ baseline isolée sur la
tâche $j$ :
\[
\mathrm{BWT} = \frac{1}{T-1}\sum_{j=1}^{T-1} (a_{T,j} - a_{j,j}), \qquad
\mathrm{FWT} = \frac{1}{T-1}\sum_{j=2}^{T} (a_{j-1,j} - \bar{b}_j).
\]
Dans nos expériences (chapitre~6), nous reportons BWT, forgetting et F1 macro
final avec la définition de $\bar{b}_j$ explicitée dans le protocole.

\subsection{Exemple numérique (3 tâches)}

Supposons trois sujets A, B, C et accuracy sur A après chaque phase :
$R[1,A]=0{,}80$, $R[2,A]=0{,}75$, $R[3,A]=0{,}70$. L'oubli de A est
$0{,}80 - 0{,}70 = 0{,}10$. Si B et C suivent le même schéma, le forgetting
moyen $\bar{f}$ agrège ces décroissances. Un algorithme CL efficace maintient
$R[3,A]$ proche de $R[1,A]$ ; le finetuning naïf typiquement fait chuter
$R[3,A]$ sous 0,50 — ce que nous observons sur 44 tâches sans rejeu (F1 0,483).

\subsection{Empreinte mémoire et latence embarquée}

L'empreinte agrège paramètres entraînables, taille du tampon de replay,
prototypes \parencite{adaimi2022} et structures auxiliaires (ex.\ termes Fisher
pour EWC). Sur mobile, on distingue latence d'inférence (une fenêtre de 3~s)
et latence d'adaptation (un pas de gradient). \textcite{amrani2025} montrent
que le fine-tuning peut rester léger pour de nouveaux utilisateurs — critère
retenu aussi dans notre cahier des charges (chapitre~4).

\section{Taxonomie des scénarios CL appliqués au HAR}

\textcite{delange2021} et \textcite{vdven2022} proposent des grilles de lecture
communes. En HAR capteur, on peut les instancier ainsi :

\begin{itemize}
  \item \textbf{Domain-incremental} : nouveau capteur ou placement (HAPT ceinture
  $\rightarrow$ WISDM poche $\rightarrow$ PAMAP2 corps). C'est le scénario le plus
  proche de nos tâches 37--44.
  \item \textbf{User-incremental} : même jeu, nouveaux sujets séquentiellement —
  protocole OCL-HAR \parencite{schiemer2023} et notre séquence 44 tâches.
  \item \textbf{Class-incremental} : nouvelles activités annotées au fil du temps
  (6 tâches $\times$ 2 classes dans nos expériences).
  \item \textbf{Task-incremental} : frontière de tâche connue au test — rare en
  déploiement réel ; LAPNet \parencite{adaimi2022} cherche justement à s'en affranchir.
\end{itemize}

\begin{figure}[H]
  \centering
  \insertfig[max width=0.85\linewidth]{delange2021_taxonomy.png}{De Lange et al. (2021)}{Taxonomie CL}
  \caption[Taxonomie continual learning \parencite{delange2021}]{Taxonomie des
  scénarios et méthodes CL \parencite{delange2021} — repère pour situer nos protocoles.}
  \label{fig:taxonomy_ch2}
\end{figure}

\begin{figure}[H]
  \centering
  \insertfig[max width=0.82\linewidth]{delange2021_eval_matrix.png}{De Lange et al. (2021)}{Matrice évaluation}
  \caption[Matrice d'évaluation CL]{Quelles métriques pour quel scénario
  \parencite{delange2021} : accuracy, forgetting, BWT, FWT.}
  \label{fig:eval_matrix_ch2}
\end{figure}

\subsection{Replay vs régularisation en HAR}

Le replay stocke un sous-ensemble de fenêtres passées et les remélange aux
nouvelles données — approche la plus robuste quand la mémoire le permet
\parencite{rebuffi2017,schiemer2023}. EWC \parencite{kirkpatrick2017} évite le
stockage explicite mais sous-estime les corrélations entre poids (Fisher diagonale)
; en HAR multi-sujets, nous l'avons gardé comme baseline intermédiaire (71,3~\%
accuracy, chapitre~6). Les prototypes \parencite{adaimi2022} compressent le replay
en représentants par classe : moins coûteux en RAM, mais sensibles au déséquilibre
quand une classe domine (class-incremental).

\subsection{Lien avec l'anticipation}

Anticiper une transition, c'est prédire une \emph{future} classe avant la fin de
la fenêtre courante. Le CL, lui, met à jour les poids pour ne pas oublier le passé.
Les deux problèmes partagent l'axe temporel mais optimisent des objectifs différents :
BCE/CE sur labels futurs vs distillation/replay sur labels passés. Notre
architecture (chapitre~4) les sépare en modules couplés seulement au niveau du
backbone gelé ou partagé — choix motivé par la lacune bibliographique (chapitre~3).

\section{Conclusion}

Régularisation, replay, isolation et prototypes donnent un vocabulaire commun
pour lire la littérature CL. Appliqué au HAR, ce vocabulaire se heurte encore
à des protocoles hétérogènes : peu d'articles combinent flux continu, nouveaux
utilisateurs, nouvelles classes \emph{et} métriques de rétention dans un même
benchmark \parencite{schiemer2023,adaimi2022,liu2024,amrani2025}. Le chapitre~3
détaille les travaux retenus ; les chapitres~4 à~6 montrent comment nous
comblons une partie de ce vide avec UWR, anticipation LSTM multi-classe et CORAL sur
notre pipeline \texttt{har\_project}.

\clearpage
% ===== END include: part-1/07-Chapter2.tex =====


% ===== BEGIN include: part-1/08-Chapter3.tex =====
\clearpage
\chapter{Travaux antérieurs sur la HAR Continue et l'Adaptation Dynamique}
\label{ch:08-Chapter3}

Ce chapitre complète le cadre théorique du chapitre~2 par une revue ciblée des
articles HAR + apprentissage continu (CL), adaptation et anticipation. Nous
expliquons d'abord comment le corpus a été construit (§3.1), puis nous regroupons
les travaux par familles (§3.2), nous les comparons dans des tableaux (§3.3) et
nous discutons ce qui manque encore (§3.4).

\section{Méthodologie de sélection des articles}

\subsection{Bases documentaires}

La recherche a surtout porté sur IEEE Xplore, ACM Digital Library (IMWUT, ISWC,
Ubicomp) et Springer (Sensors, IJMLC, etc.). Quelques prépublications arXiv ont
été gardées lorsqu'elles correspondent à la version citée dans un survey retenu.

\subsection{Requêtes et mots-clés}

Les requêtes combinaient HAR et continual learning, par exemple :
\emph{human activity recognition} AND (\emph{continual learning} OR
\emph{incremental learning} OR \emph{catastrophic forgetting}) ; variantes avec
\emph{online}, \emph{streaming}, \emph{personalization}, \emph{self-supervised},
\emph{activity anticipation}. Les chaînes de citations à partir de
\parencite{delange2021,gu2021,wang2019,parisi2019} ont permis d'élargir le corpus
au-delà du premier résultat de recherche.

\subsection{Critères d'inclusion et d'exclusion}

Nous avons retenu des articles revus par les pairs, explicitement liés à la HAR
capteur ou indispensables pour cadrer le CL (surveys, iCaRL, EWC). Sont exclus
les travaux purement vidéo sans transfert méthodologique clair (sauf pour
l'anticipation), les textes inaccessibles et les doublons arXiv/journal. Le corpus
final compte une vingtaine de références directement utiles au raisonnement du
mémoire.

\section{Catégorisation des travaux}

Le CL appliqué à la HAR progresse, mais reste « under-explored » comparé à la
vision : plusieurs articles récents le disent explicitement. Nous organisons le
corpus en cinq blocs.

\subsection{Continual learning et HAR (cœur du corpus)}

\textcite{amrani2025} homogénéisent huit bases HAR, comparent S-CNN/LSTM et
relient la démarche à une plateforme CLP. Le message principal : réduire l'écart de
F1 entre modèle fusionné et modèle entraîné jeu par jeu, puis fine-tuner pour de
nouveaux utilisateurs.

\textcite{schiemer2023} (OCL-HAR) formalisent un apprentissage en ligne avec
arrivée d'utilisateurs et de classes, distillation + rehearsal, métriques micro/macro
F1 et score d'oubli adapté au déséquilibre. Sur PAMAP2, DSADS, HAPT et WISDM,
OCL-HAR reste au-dessus des baselines tandis qu'un finetuning naïf s'effondre
rapidement — illustration concrète de l'oubli en HAR.

\textcite{adaimi2022} (LAPNet-HAR) visent un cadre \emph{task-free} :
prototypes, replay, perte contrastive. Les auteurs critiquent les hypothèses
« frontière de tâche connue » empruntées à la vision.

Ces trois articles forment le noyau CL+capteurs : intégration de données,
protocole OCL explicite, prototypes sans oracle de tâche.

\subsubsection{Amrani et al. (2025) — intégration multi-bases}

Le travail d'Amrani et al.\ \parencite{amrani2025} part d'un constat simple :
entraîner un modèle unique sur huit jeux HAR sans harmonisation produit un écart
de F1 de 24,3~\% par rapport à des modèles entraînés jeu par jeu. Leur pipeline
aligne d'abord les labels ambigus (marche vs jogging, assis vs debout), puis
entraîne des S-CNN/LSTM avec fine-tuning utilisateur. La plateforme CLP relie
l'expérience à un cadre continual explicite, même si le rejeu n'est pas le cœur
de l'article.

\begin{figure}[H]
  \centering
  \insertfig[max width=0.78\linewidth]{amrani2025_pretrain_finetune.png}{Amrani et al. (2025)}{Pré-train / fine-tune}
  \caption[Pipeline Amrani et al. \parencite{amrani2025}]{Schéma pré-entraînement global
  puis fine-tuning par utilisateur \parencite{amrani2025}.}
  \label{fig:amrani_pretrain_ch3}
\end{figure}

\begin{figure}[H]
  \centering
  \insertfig[max width=0.78\linewidth]{amrani_fig7_f1_incremental.png}{Amrani et al. (2025)}{F1 incrémental}
  \caption[F1 incrémental Amrani et al. \parencite{amrani2025}]{Évolution du F1 en
  apprentissage incrémental utilisateur \parencite{amrani2025} — repère externe.}
  \label{fig:amrani_f1_ch3}
\end{figure}

Pour notre mémoire, l'apport principal est méthodologique : montrer que
l'hétérogénéité inter-datasets n'est pas une fatalité si l'on investit dans
l'homogénéisation des labels et des fréquences. En revanche, l'anticipation et
la latence embarquée ne sont pas évaluées — d'où notre module d'anticipation séparé.

\subsubsection{Schiemer et al. (2023) — OCL-HAR}

OCL-HAR \parencite{schiemer2023} est l'article le plus proche de notre protocole
44 tâches. Les auteurs formalisent des scénarios \emph{within-subject},
\emph{between-subject} et \emph{domain shift}, avec distillation et rehearsal.
Sur PAMAP2, DSADS, HAPT et WISDM, ils rapportent des gains de +0,14 (micro-F1)
et +0,17 (macro-F1) vs la deuxième meilleure méthode, avec des scores d'oubli
modérés (0,24--0,28). Le finetuning naïf s'effondre en quelques tâches — preuve
directe de l'oubli en HAR.

Notre UWR reprend l'esprit du rehearsal mais ajoute des prototypes par classe,
une perte contrastive et un score d'incertitude (entropie softmax) pour pondérer les
mises à jour. OCL-HAR ne traite pas l'anticipation ; leurs métriques restent
notre référence pour comparer forgetting et F1 macro en conditions multi-jeux.

\subsubsection{Adaimi \& Thomaz (2022) — LAPNet-HAR}

LAPNet-HAR \parencite{adaimi2022} vise un cadre \emph{task-free} : pas de frontière
de tâche annoncée au test, ce qui correspond mieux au flux réel qu'un protocole
class-incremental artificiel. Réseaux prototypiques, rejeu et contraste limitent
l'oubli ; les auteurs mesurent explicitement Base/New/Intransigence. Sur
Opportunity, le forgetting LAPNet ($\sim$51~\%) bat le online finetuning
($\sim$60--70~\%), et l'adaptation de prototypes gagne +14 à +23 points.

La limite pour le déploiement wearable est le manque de chiffres temps réel et
empreinte mémoire. Notre choix Transformer + tampon fixe se situe entre la
richesse représentationnelle de LAPNet et la contrainte embarquée que nous
visons au chapitre~4.

\subsection{Incrémental de classes et représentation}

\subsubsection{iKAN \parencite{liu2024}}

iKAN combine des Kolmogorov-Arnold Networks (KAN) avec des branches par tâche
et une redistribution de capacité entre jeux hétérogènes. L'idée est de limiter
l'interférence en isolant partiellement les paramètres tout en partageant un
tronc commun. Sur des protocoles cross-dataset, iKAN se place au niveau des
Transformers classiques en accuracy, avec un coût d'ingénierie plus élevé.
Pour un déploiement mobile, la complexité des KAN reste un frein — nous
retenons plutôt un backbone unique + rejeu, plus simple à quantifier.

\begin{figure}[H]
  \centering
  \insertfig[max width=0.75\linewidth]{amrani_fig4_scnn_lstm.png}{Littérature HAR}{Architecture S-CNN/LSTM}
  \caption[Architecture CNN-LSTM de référence]{Exemple d'architecture S-CNN/LSTM
  utilisée dans les comparaisons multi-bases \parencite{amrani2025,gu2021}.}
  \label{fig:scnn_lstm_ch3}
\end{figure}

\subsubsection{iCaRL \parencite{rebuffi2017}}

iCaRL reste la baseline class-incremental de référence : exemplaires par classe,
distillation des logits, mise à jour incrémentale du classifieur NCM. Transposé
en HAR, iCaRL sert surtout de repère dans nos expériences (chapitre~6) plutôt
que comme architecture native capteur. Son hypothèse de frontières de tâche
connues limite la transposition directe au flux utilisateur continu.

\subsection{Self-supervised HAR}

\subsubsection{Apprentissage contrastif}

\textcite{tang2021} appliquent SimCLR-like à des signaux santé (ECG) : deux vues
augmentées d'un même segment doivent se rapprocher dans l'espace latent. Le
transfert vers IMU pur est discuté mais moins documenté que pour l'EEG. L'intérêt
pour le CL est indirect : un bon pré-train réduit le besoin de labels à l'arrivée
d'un nouvel utilisateur, sans garantir pour autant la rétention quand de nouvelles
classes apparaissent.

\subsubsection{Pré-entraînement à grande échelle}

Yuan et al.\ \parencite{yuan2024} pré-entraînent sur des centaines de milliers
d'heures de wearables avec masquage et contrastif. Les F1 downstream dépassent
93~\% sur certains benchmarks, mais le coût GPU et l'absence de protocole CL
explicite éloignent cette ligne de notre contrainte embarquée. Nous retenons
l'idée d'un pré-train supervisé plus modeste (HAPT + fusion) comme compromis.

\subsection{Personnalisation et adaptation de domaine}

\subsubsection{Batch norm adaptive}

\textcite{mazankiewicz2020} mettent à jour les statistiques de batch norm en
ligne pour suivre un utilisateur sans réentraîner tous les poids. Efficace quand
le drift est lent (fatigue, légère variation de posture), insuffisant quand
arrivent de nouvelles classes ou un capteur différent. CORAL \parencite{sun2016},
utilisé au chapitre~6, complète cette approche en alignant les covariances entre
domaines source et cible.

\subsubsection{Personnalisation de classifieurs}

Ferrari et al.\ \parencite{ferrari2020} discutent le risque de sur-apprentissage
utilisateur quand peu de fenêtres étiquetées sont disponibles. Leur message
recoupe le nôtre : la personnalisation doit être régularisée (rejeu, EWC, UWR)
pour ne pas effacer les connaissances communes apprises sur la population.

\begin{figure}[H]
  \centering
  \insertfig[max width=0.72\linewidth]{delange2021_eval_matrix.png}{De Lange et al. (2021)}{Matrice d'évaluation CL}
  \caption[Matrice d'évaluation CL \parencite{delange2021}]{Scénarios CL et métriques
  associées — cadre retenu pour nos protocoles \parencite{delange2021,vdven2022}.}
  \label{fig:delange_eval_ch3}
\end{figure}

\subsection{Anticipation d'activité}

Ryoo \parencite{ryoo2011} introduit la reconnaissance précoce en vidéo : prédire
l'activité avant la fin du geste. Gammulle et al.\ \parencite{gammulle2019}
étendent avec des modèles joints anticipation-classification. Côté capteurs,
Ordóñez \& Roggen \parencite{ordonyez2016} et Dirgová et al.\ \parencite{dirgova2022}
montrent que les architectures séquentielles (CNN-LSTM, Transformer) capturent déjà
partiellement la dynamique nécessaire à une prédiction lookahead.

\begin{figure}[H]
  \centering
  \insertfig[max width=0.78\linewidth]{dirgova2022_transformer.png}{Dirgová et al. (2022)}{Transformer HAR}
  \caption[Transformer pour HAR \parencite{dirgova2022}]{Architecture Transformer
  appliquée aux séries IMU \parencite{dirgova2022} — précurseur de notre backbone.}
  \label{fig:dirgova_ch3}
\end{figure}

Peu d'articles CL combinent anticipation et non-régression sur IMU : c'est la
lacune que notre tête LSTM multi-classe (chapitre~4) vise à combler, avec une
évaluation séparée du module CL principal.

\section{Analyse détaillée — synthèse transversale}

Au-delà du tableau comparatif, trois axes structurent la littérature récente.

\textbf{Alignement des données.} Amrani et al.\ \parencite{amrani2025} montrent
que sans homogénéisation, fusionner des jeux dégrade les performances ; avec
alignement, l'écart tombe de 24,3~\% à 7,8~\%. Notre pipeline de fusion
(HAPT+WISDM+PAMAP2) reprend cette leçon au chapitre~5.

\textbf{Protocoles CL explicites.} Schiemer \parencite{schiemer2023} et Adaimi
\parencite{adaimi2022} imposent des métriques d'oubli et des scénarios reproductibles.
Les surveys \parencite{delange2021,gu2021} rappellent que comparer deux papiers
sans préciser task- vs class-incremental est trompeur — d'où nos deux scénarios
distincts au chapitre~6.

\textbf{Anticipation vs adaptation.} La vidéo (Ryoo, Gammulle) traite l'anticipation
sans contrainte mémoire ; le CL capteur (OCL-HAR, LAPNet) traite l'adaptation sans
prédire la transition imminente. Notre contribution vise le couplage modulaire :
UWR pour l'adaptation, tête LSTM pour l'anticipation, évaluées séparément puis
discutées ensemble.

\section{Tableau comparatif (synthèse obligatoire)}

Le tableau~\ref{tab:comparatif_ch3} synthétise les travaux recensés selon six axes :
méthode, jeux de données, continual learning explicite, replay, anticipation et limites.

{\small
\begin{longtable}{|>{\raggedright\arraybackslash}p{2.0cm}|>{\raggedright\arraybackslash}p{2.6cm}|>{\raggedright\arraybackslash}p{2.1cm}|>{\centering\arraybackslash}p{1.1cm}|>{\centering\arraybackslash}p{1.1cm}|>{\centering\arraybackslash}p{1.0cm}|>{\raggedright\arraybackslash}p{2.5cm}|}
\hline
\textbf{Papier (année)} & \textbf{Méthode principale} & \textbf{Jeux / contexte} & \textbf{CL ?} & \textbf{Replay ?} & \textbf{Antic. ?} & \textbf{Limites principales} \\
\hline
\endfirsthead
\hline
\textbf{Papier (année)} & \textbf{Méthode principale} & \textbf{Jeux / contexte} & \textbf{CL ?} & \textbf{Replay ?} & \textbf{Antic. ?} & \textbf{Limites principales} \\
\hline
\endhead
Amrani et al. (2025) &
Homogénéisation multi-datasets + S-CNN/LSTM + fine-tuning ; plateforme CLP &
Huit bases HAR publiques &
Oui &
Non central &
Non &
Coût d'alignement ; peu d'évaluation embarquée \\
\hline
Schiemer et al. (2023) OCL-HAR &
OCL : distillation + rehearsal ; within/between sujet &
PAMAP2, WISDM, HAPT, DSADS &
Oui &
Oui &
Non &
Scénarios complexes ; généralisation capteur à discuter \\
\hline
Adaimi \& Thomaz (2022) LAPNet-HAR &
Prototypical networks + replay + contraste ; task-free &
Cinq bases HAR publiques &
Oui &
Oui &
Non &
Temps réel / embarqué peu détaillés \\
\hline
Liu et al. (2024) iKAN &
KAN + branches tâche + redistribution &
Jeux hétérogènes cross-dataset &
Oui &
Partiel &
Non &
Complexité ; validation terrain variable \\
\hline
Rebuffi et al. (2017) iCaRL &
Exemplars + distillation ; class-incremental &
ImageNet / CIFAR (vision) &
Paradigme CI &
Oui &
Non &
Pas natif wearable ; frontières de tâche \\
\hline
Mazankiewicz et al. (2020) &
Personnalisation incrémentielle ; DA-BN &
Signaux portés (Actitracer, etc.) &
Adaptation en ligne &
Non classique &
Non &
Souvent même utilisateur / drift léger \\
\hline
Yuan et al. (2024) &
SSL à grande échelle sur wearables &
Très grande cohorte wearables &
Non &
Non &
Non &
Calcul massif hors embarqué \\
\hline
Tang et al. (2021) &
Contrastif (SimCLR-like) santé &
PTB-XL, etc. &
Non &
Non &
Non &
Transfert vers IMU à nuancer \\
\hline
Dirgová Luptáková et al. (2022) &
Transformer sur séries HAR &
Bases type UCI &
Non &
Non &
Non &
Coût calcul ; pas CL \\
\hline
Ordóñez \& Roggen (2016) &
CNN + LSTM multimodal &
Opportunity, PAMAP2, etc. &
Non &
Non &
Non &
Pas adaptation longue durée / CL \\
\hline
Gammulle et al. (2019) &
Modèle joint anticipation &
THUMOS (vidéo) &
Non &
Non &
Oui &
Transfert direct capteur limité \\
\hline
Ryoo (2011) &
Early recognition / prédiction &
Vidéo &
Non &
Non &
Oui &
Hors IMU portable \\
\hline
De Lange et al. (2021) &
Survey CL classification &
N/A &
Cadre général &
N/A &
N/A &
Pas spécifique HAR capteur \\
\hline
Gu (2021) ; Wang (2019) &
Surveys deep learning HAR &
N/A &
Cadre général &
N/A &
N/A &
CL peu détaillé vs architectures \\
\hline
\end{longtable}
}

\captionof{table}{Synthèse comparative des travaux recensés (Chapitre~3)}
\label{tab:comparatif_ch3}

\subsection{Résultats quantitatifs rapportés (extraits des articles)}

Le tableau~\ref{tab:quantitatif_ch3} complète la synthèse méthodologique par les
indicateurs numériques publiés. «~N/R~» signifie non rapporté ou non comparable.

\begin{table}[H]
\centering
\small
\begin{tabularx}{\textwidth}{|>{\raggedright\arraybackslash}Y|>{\raggedright\arraybackslash}p{1.6cm}|>{\raggedright\arraybackslash}p{1.6cm}|>{\raggedright\arraybackslash}Y|>{\raggedright\arraybackslash}Y|>{\raggedright\arraybackslash}Y|}
\hline
\textbf{Papier} & \textbf{Micro-F1 (gain)} & \textbf{Macro-F1 (gain)} & \textbf{F1 / accuracy rapportés} & \textbf{Forgetting / stabilité} & \textbf{Note de comparabilité} \\
\hline
Amrani et al. (2025) &
N/R & N/R &
Écart F1 vs baseline sujet-dépendant : 24,3~\% $\rightarrow$ 7,8~\% (gain relatif 16,5~\%). Fine-tuning : +2,5~\% accuracy (nouveaux utilisateurs). &
N/R &
Protocole multi-datasets ; setup différent de OCL-HAR / LAPNet. \\
\hline
Schiemer et al. (2023) OCL-HAR &
+0,14 (within-subject vs 2\textsuperscript{e} meilleure méthode) &
+0,17 (idem) &
Jusqu'à +0,17 / +0,23 (micro/macro F1) selon configuration &
Scores d'oubli bas (ex.\ 0,24--0,28) &
Scénarios within/between/shift ; métriques par jeu. \\
\hline
Adaimi \& Thomaz (2022) LAPNet-HAR &
N/R & N/R &
Amélioration base/new vs baselines ; online finetuning abaisse les classes de base ($\sim$45~\% en moyenne). &
Forgetting mesuré ; ex.\ $\sim$51~\% (Opportunity) LAPNet vs $\sim$60--70~\% (online FT). Prototype adaptation : +14 à +23 pts. &
Indicateurs Base/New/intransigence ; jeux et ablations différents. \\
\hline
\end{tabularx}
\caption{Indicateurs quantitatifs extraits des articles de référence}
\label{tab:quantitatif_ch3}
\end{table}

\begin{figure}[H]
  \centering
  \insertfig{figure_positionnement_ch3.png}{Synthèse bibliographique}{Carte de positionnement}
  \caption[Positionnement des travaux recensés]{Carte de positionnement — apprentissage continu vs
  anticipation dans la littérature HAR.}
  \label{fig:positionnement_ch3}
\end{figure}

\begin{figure}[H]
  \centering
  \insertfig{figure_comparaison_quantitative_ch3.png}{Synthèse bibliographique}{Comparaison quantitative}
  \caption[Comparaison quantitative des approches]{Comparaison quantitative iKAN, LAPNet-HAR,
  OCL-HAR, Amrani et al. (2025) et baseline.}
  \label{fig:comparaison_ch3}
\end{figure}

\section{Analyse critique}

Quatre constats ressortent de cette revue.

\textbf{CL et anticipation restent découplés.} Amrani, Schiemer, Adaimi ou Liu
optimisent la classification stable dans le temps ; Gammulle ou Ryoo traitent la
prédiction sans oubli mesuré quand de nouveaux utilisateurs arrivent. Peu
d'articles combinent les deux — c'est l'espace que nous visons avec l'anticipation
LSTM + UWR.

\textbf{Peu de mesures « système ».} Même quand le vocabulaire dit « online » ou
« real-time », les expériences tournent souvent sur serveur, avec F1 et oubli, mais
sans latence bout-en-bout ni empreinte RAM mesurée sur la cible wearable.

\textbf{Personnalisation encore grossière.} Batch norm adaptive ou fine-tuning
utilisateur aident, mais peinent sur des variations fines (fatigue, pathologie
légère) avec peu de labels par personne.

\textbf{Datasets vs monde réel.} Fusionner HAPT, WISDM ou PAMAP2 réduit le biais
d'un seul protocole \parencite{amrani2025}, sans remplacer des collectes longues
en conditions libres. Le SSL massif \parencite{yuan2024} améliore les
représentations, pas automatiquement le CL avec nouvelles classes en ligne.

\subsection{Positionnement de notre contribution}

Le tableau~\ref{tab:positionnement_nous} situe explicitement notre travail par
rapport aux articles du corpus. Nous ne prétendons pas couvrir toute la grille :
l'objectif est de combiner, dans un même dépôt reproductible, pré-train Transformer,
UWR, anticipation LSTM, CORAL et détection de chutes proxy.

\begin{table}[H]
\centering
\small
\begin{tabularx}{\textwidth}{|>{\raggedright\arraybackslash}Y|c|c|c|c|}
\hline
\textbf{Critère} & \textbf{OCL-HAR} & \textbf{LAPNet} & \textbf{Amrani 2025} & \textbf{Notre système} \\
\hline
CL user-incrémental & Oui & Partiel & Oui (fine-tune) & Oui (44 tâches) \\
CL class-incrémental & Oui & Oui & Non central & Oui (6 tâches) \\
Replay / prototypes & Oui & Oui & Non & UWR + prototypes \\
Anticipation IMU & Non & Non & Non & LSTM multi-classe \\
Cross-dataset & Oui & Oui & Oui (8 bases) & CORAL HAPT$\rightarrow$WISDM \\
Code ouvert reproductible & Partiel & Partiel & Plateforme CLP & \texttt{har\_project} \\
\hline
\end{tabularx}
\caption{Positionnement qualitatif de notre contribution}
\label{tab:positionnement_nous}
\end{table}

\subsection{Limites de la revue}

Le corpus a été construit par recherche thématique et chaînes de citations ;
une revue systématique PRISMA n'a pas été menée. Les travaux très récents (2025)
postérieurs à la rédaction finale pourraient compléter iKAN ou le SSL wearables.
Enfin, la comparabilité numérique directe reste limitée : jeux, splits et métriques
diffèrent entre OCL-HAR, LAPNet et nos protocoles — d'où l'importance de fixer
nos propres baselines au chapitre~6 plutôt que de sur-interpréter un écart de F1
avec la littérature.

En résumé, le couple CL–HAR mûrit (OCL-HAR, LAPNet, iKAN, intégration
multi-bases), mais il manque encore une solution qui anticipe, s'adapte en flux,
et reste compatible embarqué — problème posé en partie~II.

\subsection{Scénarios types et couverture bibliographique}

Pour clarifier ce que la revue couvre — et ce qu'elle ne couvre pas — nous
proposons la lecture croisée suivante :

\begin{itemize}
  \item \textbf{Nouvel utilisateur, mêmes classes, même capteur} : couvert par
  OCL-HAR \parencite{schiemer2023}, Amrani fine-tuning \parencite{amrani2025},
  Mazankiewicz DA-BN \parencite{mazankiewicz2020}. Notre protocole 44 tâches
  étend ce cas au multi-dataset.
  \item \textbf{Nouvelles classes, même utilisateur} : couvert par iCaRL
  \parencite{rebuffi2017}, iKAN \parencite{liu2024}, LAPNet
  \parencite{adaimi2022}. Notre scénario 6 tâches $\times$ 2 classes suit
  cette ligne.
  \item \textbf{Nouveau capteur / placement} : partiellement couvert (domain
  shift dans OCL-HAR, CORAL en vision \parencite{sun2016}). Nos tâches 37--44
  et expérience CORAL comblent partiellement le gap côté IMU.
  \item \textbf{Anticipation + CL simultanés} : non couvert de façon satisfaisante
  — lacune centrale de ce mémoire.
\end{itemize}

Les fiches détaillées en annexe~F approfondissent chaque référence citée.

\subsection{Feuille de route implicite du mémoire}

La revue justifie les choix suivants, repris en partie~II :
\begin{enumerate}
  \item Backbone Transformer plutôt que KAN/iKAN (simplicité déploiement).
  \item UWR plutôt que replay uniforme seul (frontières incertaines).
  \item Anticipation LSTM sur fenêtres partielles.
  \item CORAL sur embeddings pré-entraînés (cross-dataset WISDM).
  \item Protocoles 10 et 44 sujets (comparabilité doc + robustesse multi-jeux).
\end{enumerate}

Chaque choix pourrait être contesté par une alternative de la littérature ; l'enjeu
n'est pas l'optimalité absolue sur un benchmark, mais une chaîne cohérente et
reproductible répondant simultanément à CL, anticipation et domain shift — combinaison
rare dans les articles recensés.

\section{Conclusion}

Ce chapitre a sélectionné et comparé les travaux les plus proches de notre sujet.
Les tableaux §3.3 fixent les références ; l'analyse §3.4 justifie nos choix
(UWR, anticipation LSTM, CORAL, protocole multi-scénarios). Le chapitre~4
formalise la réponse proposée.

\clearpage
% ===== END include: part-1/08-Chapter3.tex =====



% ===== BEGIN include: 09-Conception.tex =====
\part{Contribution et Évaluation}

\chapter{Conception de la solution}
\label{ch:09-Conception}

\section{Introduction}

Les chapitres précédents ont montré deux limites récurrentes. D'abord, un
modèle HAR classique oublie vite ce qu'il savait dès qu'on lui présente de
nouveaux utilisateurs ou de nouvelles activités. Ensuite, la plupart des
systèmes se contentent de reconnaître l'activité en cours : ils n'anticipent
pas une transition.

Nous répondons avec un Transformer partagé et deux modules autour. Le premier
fait de la reconnaissance continue grâce à un tampon de rejeu pondéré par
l'incertitude de prédiction (UWR). Le second anticipe l'activité suivante
avec une tête LSTM sur des fenêtres partielles. Les représentations apprises
par le backbone servent aux deux.

Le reste du chapitre formalise le problème et les scénarios, fixe les
objectifs, présente les données, décrit l'architecture, puis détaille UWR,
l'anticipation et le protocole d'évaluation.

\section{Définition et formulation du problème}

\subsection{Cadre formel et scénarios retenus}

\subsubsection{Formulation générale}

Soit $\mathcal{S} = \{(x_1, y_1), (x_2, y_2), \ldots\}$ un flux de données IMU non stationnaire, où chaque observation $x_t \in \mathbb{R}^{T \times C}$ est une fenêtre temporelle de $T = 150$ échantillons sur $C = 6$ canaux (accélération triaxiale $a_x, a_y, a_z$ et vitesse angulaire triaxiale $\omega_x, \omega_y, \omega_z$), et $y_t \in \mathcal{Y}$ est l'étiquette d'activité associée. Le flux est non stationnaire en ce sens que la distribution jointe $p(x, y)$ évolue au cours du temps, soit en raison de l'apparition de nouveaux utilisateurs aux habitudes motrices distinctes, soit en raison de l'introduction de nouvelles classes d'activités.

L'objectif général est d'entraîner un modèle $f_\theta : \mathbb{R}^{T \times C} \to \mathcal{Y}$ qui maintient des performances élevées sur l'ensemble des distributions rencontrées, sans accéder à l'intégralité des données historiques lors de chaque mise à jour, et sous contrainte mémoire bornée.

\subsubsection{Scénario 1 — Apprentissage incrémental par utilisateurs}

Dans ce scénario, les données arrivent par blocs correspondant à des utilisateurs distincts $\mathcal{U} = \{u_1, u_2, \ldots, u_K\}$. Les classes d'activités restent fixes ($\mathcal{Y}$ constant), mais la distribution conditionnelle $p(x \mid y, u)$ varie d'un utilisateur à l'autre en raison des différences morphologiques, de placement du capteur et de style de mouvement. Le modèle doit, après avoir appris sur l'utilisateur $u_k$, conserver ses performances sur tous les utilisateurs précédents $u_1, \ldots, u_{k-1}$ sans y avoir accès, tout en s'adaptant efficacement au nouvel utilisateur $u_k$.

Ce scénario correspond à la situation de déploiement la plus fréquente dans les applications de santé connectée : un même modèle installé sur un dispositif est progressivement personnalisé pour de nouveaux porteurs.

\subsubsection{Scénario 2 — Apprentissage incrémental par classes}

Dans ce scénario, les données arrivent par blocs correspondant à des tâches $\mathcal{T} = \{T_1, T_2, \ldots, T_N\}$, chaque tâche $T_k$ introduisant un sous-ensemble de nouvelles classes $\mathcal{Y}_k$ tel que $\mathcal{Y}_i \cap \mathcal{Y}_j = \emptyset$ pour $i \neq j$ et $\mathcal{Y} = \bigcup_{k=1}^{N} \mathcal{Y}_k$. Le modèle doit, à chaque tâche $k$, apprendre à reconnaître les nouvelles classes tout en maintenant la discrimination des classes des tâches précédentes. Ce scénario est formellement le plus difficile car il exige une expansion dynamique de l'espace de sortie.

Conformément à la taxonomie de \parencite{vdven2022}, l'identifiant de tâche n'est pas fourni au modèle lors de l'inférence (test sans oracle de tâche), ce qui constitue la configuration la plus réaliste.

\subsection{Contraintes de déploiement}

Le système est conçu pour un déploiement sur dispositifs embarqués à ressources contraintes, comme les smartphones et les montres connectées. Les contraintes suivantes sont imposées tout au long de la conception :

Contrainte mémoire. Le tampon de rejeu est borné à $M = 2000$ exemples au maximum, représentant moins de 4 \% du jeu d'entraînement total. Cette contrainte modélise la limitation de la mémoire vive disponible sur un dispositif mobile.

Contrainte temporelle. Le délai d'inférence doit rester compatible avec un usage
temps réel. Nous fixons comme \emph{objectif de conception} de traiter une
fenêtre de 3~secondes en moins de 150~ms sur CPU, ce qui exclut les
architectures trop profondes ou les inférences multi-passes coûteuses en
production (cette latence n'est pas une mesure expérimentale reportée ici).

Contrainte d'accès aux données. Lors de l'apprentissage d'une nouvelle tâche, le modèle n'a accès qu'aux données de la tâche courante et au contenu de son tampon de rejeu. Aucun accès aux données brutes des tâches passées n'est autorisé, conformément aux hypothèses standard de l'apprentissage continu.

Contrainte de stabilité. Les mises à jour du modèle doivent être incrémentales et ne pas dégrader significativement les performances sur les tâches précédentes. Le taux d'oubli toléré est défini comme une dégradation de moins de 5 points de F1 macro par rapport aux performances initiales sur chaque tâche.

\subsection{Positionnement par rapport à l'état de l'art}

Par rapport à iCaRL, qui choisit surtout des exemplaires proches des
prototypes, UWR privilégie les fenêtres où le modèle hésite. L'idée est
simple : ce sont souvent ces points, près des frontières, qu'il faut
réviser pour ne pas les perdre trop vite.

Sur l'anticipation, nous gardons une formulation multi-classe classique
(classe suivante depuis des fenêtres partielles, filtre
\texttt{transitions\_only}). Le déséquilibre des transitions rend la tâche
difficile ; le chapitre~6 montre néanmoins un macro-F1 nettement au-dessus
d'une baseline majorité.

Pour le cross-dataset, plutôt qu'un alignement de covariance type CORAL
dans le code livré, nous utilisons un adaptateur affine few-shot
(\texttt{DomainAdapter}) sur les embeddings du Transformer gelé.

\section{Objectifs du projet et contributions principales}

Ce projet poursuit cinq objectifs principaux :

Objectif 1 : Backbone universel pour la HAR. Concevoir et entraîner un encodeur Transformer léger, capable de produire des embeddings de haute qualité à partir de signaux IMU bruts, transférables à toutes les tâches en aval.

Objectif 2 : Apprentissage continu robuste. Développer un mécanisme de gestion de la mémoire et d'entraînement incrémental permettant de maintenir les performances sur les tâches passées tout en assimilant efficacement les nouvelles, dans les deux scénarios définis.

Objectif 3 : Anticipation temps réel. Proposer un module d'anticipation opérationnel en temps réel, capable de détecter les transitions imminentes d'activité avec une précision et un rappel suffisants pour des applications de sécurité (détection de chute imminente, surveillance de personnes âgées).

Objectif 4 : Adaptation cross-dataset. Évaluer la capacité du système à se transférer d'un dataset source à un dataset cible de distribution différente, avec et sans annotations cibles, via un adaptateur affine few-shot (\texttt{DomainAdapter}).

Objectif 5 : Détection de chutes hybride. Intégrer un module complémentaire de détection de chutes combinant règles physiques et classification apprise, démontrant la généricité de l'architecture proposée.

Les contributions scientifiques et techniques originales de ce travail sont :

(C1) Uncertainty-Weighted Replay (UWR) : sélection et échantillonnage du
tampon selon l'entropie softmax (une passe, paramètre $\alpha$), appliqué
au HAR continu.

(C2) Module d'anticipation multi-classe sur fenêtres partielles : tête LSTM
entraînée sur le backbone gelé, contexte $S=5$ fenêtres tronquées à
$p\in\{0{,}25,0{,}50,0{,}75\}$, filtrage \texttt{transitions\_only}, perte
cross-entropie.

(C3) Évaluation multi-dataset de l'apprentissage continu : protocole d'évaluation rigoureux couvrant deux scénarios, plusieurs baselines, et une mesure explicite de l'oubli (BWT) et du transfert (FWT).

(C4) Adaptation de domaine (\texttt{DomainAdapter}) : transformation affine few-shot sur embeddings d'un Transformer pré-entraîné (backbone gelé), pour réduire le domain shift cross-dataset (WISDM / PAMAP2).

\section{Données : homogénéisation et pipeline}

\subsection{Datasets retenus et leurs caractéristiques}

HAPT — UCI HAR with Postural Transitions

Le dataset HAPT (Reyes-Ortiz et al., 2016) est une extension du dataset UCI HAR qui inclut explicitement les transitions posturales entre activités statiques. Il comprend 30 sujets adultes (âge : 19–48 ans), équipés d'un smartphone Samsung Galaxy S2 fixé à la ceinture. Les données sont acquises à 50 Hz via l'accéléromètre et le gyroscope triaxiaux. Le dataset distingue 12 classes : 6 activités de base (marche, montée d'escaliers, descente d'escaliers, position debout, assis, allongé) et 6 transitions posturales (debout-vers-assis, assis-vers-debout, assis-vers-allongé, allongé-vers-assis, debout-vers-allongé, allongé-vers-debout).

Après application du pipeline de prétraitement décrit en §4.4.2, on obtient 56 861 fenêtres d'entraînement et de test, avec une distribution de classes déséquilibrée : les activités statiques représentent environ 45 \% des fenêtres, les activités dynamiques 42 \%, et les transitions posturales seulement 13 \%.

WISDM — Wireless Sensor Data Mining

Le dataset WISDM (Kwapisz et al., 2011) a été collecté auprès de 51 sujets portant leur smartphone en poche. La fréquence d'acquisition est de 20 Hz (rééchantillonnée à 50 Hz dans notre pipeline). Il couvre 6 activités : marche, jogging, montée d'escaliers, descente d'escaliers, assis, debout. Ce dataset est exclusivement utilisé dans les expériences d'adaptation de domaine, jouant le rôle du domaine cible dont on cherche à aligner la distribution sur HAPT (domaine source).

\begin{table}[H]
\centering
\small
\begin{tabularx}{0.92\textwidth}{|Y|c|c|}
\hline
\textbf{Propriété} & \textbf{HAPT} & \textbf{WISDM} \\
\hline
Sujets                    & 30  & 51 \\
Classes                   & 12  & 6 \\
Fréquence originale       & 50 Hz & 20 Hz (rééchant. 50 Hz) \\
Placement                 & Ceinture & Poche \\
Fenêtres (après prétrait.) & $\sim$56\,861 & $\sim$17\,000 \\
Déséquilibre max.         & $\times 3{,}4$ & $\times 2{,}1$ \\
\hline
\end{tabularx}
\caption{Comparaison des caractéristiques HAPT et WISDM}
\label{tab:hapt_wisdm}
\end{table}

\subsection{Pipeline de prétraitement}

Le pipeline de prétraitement suit les recommandations d'Amrani et al. (2025) pour l'homogénéisation de datasets IMU hétérogènes. Il comprend les étapes suivantes, implémentées dans le module preprocessing.py :

Étape 1 — Rééchantillonnage à 50 Hz. Pour les signaux dont la fréquence d'acquisition diffère de 50 Hz (cas de WISDM à 20 Hz), un rééchantillonnage par méthode de Fourier (scipy.signal.resample) est appliqué. Cette méthode préserve le contenu fréquentiel du signal tout en modifiant sa résolution temporelle.

Étape 2 — Suppression de la composante gravitationnelle. La gravité terrestre ($g \approx 9{,}81$ m/s²) est présente dans les mesures accéléromètre et constitue un biais qui varie selon l'orientation du capteur. Elle est estimée par un filtre passe-bas Butterworth d'ordre 5, de fréquence de coupure 0,5 Hz, appliqué sur chaque canal accéléromètre. La composante gravitationnelle ainsi estimée est soustraite du signal brut, isolant la composante dynamique pure liée au mouvement.

Étape 3 — Segmentation par fenêtre glissante. Le signal continu est découpé en fenêtres de $T = 150$ échantillons (équivalant à 3 secondes à 50 Hz) avec un chevauchement de 50 \% (stride $s = 75$ échantillons). Cette configuration est standard dans la littérature HAR et permet d'obtenir un bon compromis entre la résolution temporelle et la quantité de données. L'étiquette de chaque fenêtre est déterminée par vote majoritaire sur les étiquettes échantillon-par-échantillon de la fenêtre.

Étape 4 — Normalisation z-score par sujet. Pour chaque sujet, la moyenne et l'écart-type de chaque canal sont calculés sur l'ensemble de ses données d'entraînement. La normalisation est ensuite appliquée fenêtre par fenêtre, ramenant chaque canal à une distribution centrée-réduite. Cette normalisation par sujet, plutôt que globale, réduit les biais inter-sujets dus aux différences de calibration des capteurs.

\subsection{Unification des classes}

Pour les expériences d'adaptation de domaine cross-dataset (HAPT → WISDM), un mapping manuel est établi entre les classes des deux datasets en conservant uniquement les 5 activités communes : marche (HAPT: classe 1, WISDM: classe 1), montée d'escaliers (HAPT: 2, WISDM: 3), descente d'escaliers (HAPT: 3, WISDM: 4), assis (HAPT: 4, WISDM: 5), debout (HAPT: 5, WISDM: 6). Les transitions posturales spécifiques à HAPT et le jogging spécifique à WISDM sont exclus de ces expériences. Le mapping garantit une correspondance sémantique stricte entre les classes source et cible.

\section{Architecture générale du système}

\subsection{Vue d'ensemble}

L'architecture proposée est structurée autour d'un backbone partagé de type Transformer et de deux modules spécialisés branchés en aval. Cette organisation permet de mutualiser l'apprentissage des représentations générales du mouvement tout en permettant à chaque module d'optimiser sa tâche spécifique de manière indépendante.

Le flux de traitement global est le suivant : une fenêtre IMU brute $(T=150, C=6)$ est d'abord encodée par le backbone en un vecteur d'embedding de dimension $d=128$. Cet embedding est ensuite transmis soit au Module~1 (reconnaissance continue) pour la classification et la mise à jour incrémentale, soit au Module~2 (anticipation) pour la prédiction de l'activité suivante à partir de fenêtres partielles.

\input{includes/tikz-har-system}

\subsection{Backbone Transformer}

Le backbone IMUTransformerEncoder est un encodeur Transformer pour séries temporelles multivariées, conçu spécifiquement pour les signaux IMU. Son architecture s'inspire de Dirgová et al. (2022) et comprend quatre composants principaux :

Projection d'entrée. Les $C = 6$ canaux bruts sont projetés vers l'espace de représentation de dimension $d = 128$ par une couche linéaire suivie d'une normalisation de couche (LayerNorm). Cette projection est nécessaire car la dimension 6 est trop petite pour représenter des concepts abstraits de mouvement.

Token CLS. Un token de classification spécial, [CLS], appris par optimisation, est prépendu à la séquence. Après passage dans tous les blocs Transformer, la représentation de ce token agrège l'information de toute la fenêtre par le mécanisme d'attention.

Encodage positionnel. Des codes positionnels sinusoïdaux fixes (sin/cos à différentes fréquences) sont additionnés aux représentations après projection, permettant au modèle de connaître la position temporelle de chaque échantillon. Un facteur d'échelle appris est appliqué aux codes positionnels.

Blocs Transformer. Quatre blocs Transformer encodeurs sont empilés. Chaque bloc applique, dans l'ordre : normalisation pré-couche → attention multi-têtes → connexion résiduelle → normalisation → FFN (Feed-Forward Network) → connexion résiduelle. L'attention multi-têtes utilise $H = 4$ têtes de dimension 32 chacune ($4 \times 32 = 128$). La FFN est composée de deux couches linéaires avec activation GELU et une dimension cachée de 256.

La sortie du backbone est la représentation du token CLS après le dernier bloc et une normalisation finale : $\mathbf{e} = f_\theta(x) \in \mathbb{R}^{128}$.

\begin{figure}[H]
  \centering
  \insertfig[max width=0.88\linewidth]{dirgova2022_transformer.png}{Architecture Transformer IMU \parencite{dirgova2022}}{Encodeur Transformer pour séries temporelles IMU.}
  \caption[Architecture du backbone Transformer]{Architecture du backbone \emph{IMUTransformerEncoder}
  inspirée de \parencite{dirgova2022} — projection, token CLS, encodage positionnel et blocs d'attention.}
  \label{fig:backbone_transformer}
\end{figure}

Nombre total de paramètres : 531\,457, dont 118\,784 pour les blocs Transformer et 412\,673 pour les couches de projection et de classification.

\subsection{Module 1 — Reconnaissance continue}

Le Module 1 gère l'ensemble du cycle de vie de la reconnaissance d'activité en mode continu. Il est composé de trois sous-composants :

Tête de classification : une séquence linéaire (128 → 64 → dropout(0,1) → n\_classes) utilisée lors du pré-entraînement supervisé.

Mémoire de prototypes : un vecteur moyen par classe maintenu et mis à jour en ligne.

Tampon de rejeu UWR : stockage borné des exemples passés, avec priorisation par incertitude épistémique.

Ce module est décrit en détail en §4.6.

\subsection{Module 2 — Anticipation d'activité}

Le Module 2 opère sur le backbone gelé (mode standard). Il prend en entrée une
séquence de $S=5$ fenêtres partielles, les encode via le Transformer, puis les
agrège avec un LSTM unidirectionnel dont le dernier état alimente un
classifieur linéaire multi-classe (\texttt{AnticipationHead}). La cible est le
label de la fenêtre suivante ; l'entraînement filtre les transitions
(\texttt{transitions\_only=True}).

Ce module est décrit en détail en §4.7.

\section{Conception du module de reconnaissance continue}

\subsection{Mémoire de prototypes}

Un prototype de classe $\boldsymbol{\mu}_c \in \mathbb{R}^{128}$ est associé à chaque classe $c \in \mathcal{Y}$. Il est défini comme la moyenne exponentielle mobile des embeddings des exemples de la classe $c$ observés jusqu'à l'instant courant :

$$\boldsymbol{\mu}_c \leftarrow \alpha \cdot \boldsymbol{\mu}_c + (1 - \alpha) \cdot \frac{1}{|\mathcal{B}_c|} \sum_{x \in \mathcal{B}_c} f_\theta(x)$$

où $\mathcal{B}_c$ est l'ensemble des exemples de classe $c$ dans le lot courant et $\alpha = 0{,}9$ est le coefficient de momentum. Cette mise à jour s'effectue après chaque pas de gradient, garantissant que les prototypes reflètent toujours les représentations du modèle courant.

Lors de l'inférence continue, la classe prédite est celle dont le prototype est le plus proche en distance euclidienne de l'embedding de l'exemple test :

$$\hat{y} = \arg\min_{c \in \mathcal{Y}} |\boldsymbol{e} - \boldsymbol{\mu}_c|_2$$

Cette stratégie de classification par plus proche prototype présente l'avantage de ne pas nécessiter de recalibration de la tête de classification lors de l'ajout de nouvelles classes, contrairement aux approches softmax classiques.

\subsection{Tampon de rejeu standard}

Un tampon de capacité bornée $M = 2000$ exemples est maintenu. Sa structure est une liste circulaire qui associe à chaque exemple $(x, y)$ un score de priorité $s \in [0, 1]$ initialisé à 0,5. L'ajout d'un nouvel exemple se fait selon la stratégie de réservoir : si le tampon est plein, l'exemple est inséré avec probabilité $M / n_{total}$ en remplaçant un exemple existant tiré aléatoirement. Lors de l'entraînement, un lot de rejeu de taille $B_{replay} = 32$ est échantillonné uniformément depuis le tampon et la perte de cross-entropie est calculée sur ces exemples :

$$\mathcal{L}_{\mathrm{replay}} = \mathcal{L}_{\mathrm{CE}}(f_\theta(x_{\mathrm{replay}}), y_{\mathrm{replay}})$$

Cette perte de rejeu est additionnée à la perte sur le lot courant, forçant le modèle à maintenir ses prédictions sur les données passées.

\subsection{Uncertainty-Weighted Replay — contribution originale}

Motivation et intuition

La stratégie de rejeu uniforme présente une limitation fondamentale : elle ne distingue pas les exemples faciles (bien classifiés, loin des frontières de décision) des exemples difficiles (près des frontières, source d'erreurs). Or, les frontières de décision sont précisément les régions de l'espace de représentation qui risquent de se déplacer lors de l'apprentissage de nouvelles tâches, causant ainsi l'oubli. En concentrant la mémoire sur les exemples à forte incertitude, l'UWR maximise la couverture des frontières de décision avec un budget mémoire fixe.

Estimation de l'incertitude par entropie de prédiction

Dans l'implémentation livrée (\texttt{uncertainty\_replay.py}), l'incertitude d'un exemple $x$ est l'entropie de Shannon de la distribution softmax \emph{en une seule passe} (mode \texttt{eval}) :

$$p(y \mid x) = \mathrm{softmax}(f_\theta(x)), \qquad
\mathcal{U}(x) = -\sum_{c=1}^{|\mathcal{Y}|} p(c \mid x) \log p(c \mid x).$$

Les scores sont ensuite élevés à la puissance $\alpha$ (défaut $\alpha=1$) et normalisés pour former une distribution d'échantillonnage. Cette mesure est moins coûteuse que le Monte Carlo Dropout \parencite{gal2016} ($T$ passes) : nous citons Gal \& Ghahramani comme fondement théorique de l'incertitude, mais le code priorise l'entropie predictive single-pass, recalculée périodiquement sur le tampon (\texttt{update\_freq}).

Algorithme UWR

Lors de l'ajout d'un exemple au tampon : le score de priorité de l'exemple est fixé à $s(x) = \mathcal{U}(x) / \log |\mathcal{Y}|$, normalisant l'incertitude dans $[0, 1]$. Si le tampon est plein, l'exemple est inséré en remplaçant un exemple existant tiré avec probabilité inversement proportionnelle à son score de priorité (les exemples les moins incertains sont préférentiellement évincés).

Lors de l'échantillonnage pour le rejeu : les exemples du tampon sont sélectionnés avec une probabilité proportionnelle à leur score de priorité, légèrement lissée par une constante $\epsilon = 0{,}1$ pour éviter la famine des exemples à faible incertitude :

$$P(\text{sélectionner } x_i) \propto s(x_i) + \epsilon$$

Propriétés théoriques

L'UWR converge vers un tampon dont la distribution marginale est concentrée sur les exemples à forte incertitude, c'est-à-dire près des frontières de décision. Lorsque le modèle évolue et que certaines frontières se déplacent, les exemples affectés voient leur incertitude augmenter, déclenchant leur priorisation dans les futurs lots de rejeu. Ce mécanisme crée une boucle de rétroaction qui guide activement le tampon vers les zones à risque d'oubli.

\subsection{Calibration adaptative en ligne (prototypes)}

Après chaque pas de gradient, les prototypes de toutes les classes présentes dans le lot courant sont recalibrés. Pour une classe $c$ présente dans le lot, l'embedding moyen du lot $\bar{\boldsymbol{e}}_c = \frac{1}{|\mathcal{B}_c|}\sum_{x \in \mathcal{B}_c} f_\theta(x)$ est calculé avec les nouveaux paramètres $\theta$ mis à jour. Le prototype est ensuite mis à jour selon la formule présentée en §4.6.1. Pour les classes absentes du lot courant, les prototypes sont conservés inchangés.

Cette recalibration garantit que les prototypes restent cohérents avec les représentations courantes du backbone, évitant un décalage progressif (prototype drift) qui dégraderait les performances de l'inférence par plus proche voisin.

\subsection{Fonction de perte totale}

La fonction de perte globale lors d'un pas d'apprentissage continu combine trois termes :

$$\mathcal{L}_{\mathrm{total}} = \mathcal{L}_{\mathrm{CE}}(\mathcal{B}_{\mathrm{curr}}) + \mathcal{L}_{\mathrm{replay}}(\mathcal{B}_{\mathrm{replay}}) + \lambda_c \cdot \mathcal{L}_{\mathrm{contrast}}(\mathcal{B}_{\mathrm{curr}})$$

Terme 1 — Perte de classification courante $\mathcal{L}_{\mathrm{CE}}(\mathcal{B}_{\mathrm{curr}})$ : cross-entropie sur le lot courant, assurant l'apprentissage de la tâche courante.

Terme 2 — Perte de rejeu $\mathcal{L}_{\mathrm{replay}}(\mathcal{B}_{\mathrm{replay}})$ : cross-entropie sur le lot échantillonné depuis le tampon UWR, assurant la rétention des connaissances passées.

Terme 3 — Perte contrastive $\mathcal{L}_{\mathrm{contrast}}(\mathcal{B}_{\mathrm{curr}})$ : perte de séparation inter-classes calculée par rapport à la mémoire de prototypes. Pour chaque exemple $(x, y)$ du lot courant, elle maximise la similarité cosinus avec le prototype de sa classe et minimise celle avec les prototypes des autres classes, avec une température $\tau = 0{,}07$ :

$$\mathcal{L}_{\mathrm{contrast}}(x, y) = -\log \frac{\exp(\mathrm{sim}(\boldsymbol{e}, \boldsymbol{\mu}_y) / \tau)}{\sum_{c \neq y} \exp(\mathrm{sim}(\boldsymbol{e}, \boldsymbol{\mu}_c) / \tau)}$$

Le coefficient $\lambda_c = 0{,}1$ est fixé empiriquement pour équilibrer la contribution de la perte contrastive sans dominer la dynamique d'apprentissage.

\section{Conception du module d'anticipation}

\subsection{Architecture LSTM sur fenêtres partielles}

Le module d'anticipation prédit la \emph{classe suivante} à partir d'un
contexte partiel. Pour un indice $i$ :

\begin{itemize}
  \item le contexte est la séquence de $S=5$ fenêtres
  $X[i-S],\ldots,X[i-1]$, chacune tronquée aux premiers $p\,\%$
  ($p\in\{0{,}25,0{,}50,0{,}75\}$) ;
  \item la cible est le label $y[i]$ ;
  \item avec \texttt{transitions\_only=True}, on ne garde l'exemple que si
  $y[i]\neq y[i-1]$ ;
  \item chaque fenêtre partielle est encodée par le backbone (gelé) ;
  \item la séquence d'embeddings $(B,S,128)$ passe dans un LSTM
  unidirectionnel (2 couches, hidden 128, dropout 0,2) ;
  \item le dernier état alimente \texttt{LayerNorm + Linear} (logits
  multi-classe) ;
  \item la perte est la cross-entropie (\texttt{F.cross\_entropy}).
\end{itemize}

L'entraînement standard gèle le backbone et la tête HAR
(\texttt{train\_anticipation.py}) ; une variante conjoint
(\texttt{train\_anticipation\_joint.py}) adapte aussi le backbone avec un
taux d'apprentissage réduit.

\subsection{Limites de la formulation multi-classe}

Avec truncature en fin de fenêtre, pondération des classes et restauration
du meilleur checkpoint, la prédiction multi-classe atteint un macro-F1 de
validation d'environ 0,59 ($p=25\%$), 0,63 ($p=50\%$) et 0,64 ($p=75\%$),
contre une baseline majorité d'environ 0,04. La tâche reste non triviale
(transitions rares, déséquilibre, split par sujet), mais nettement au-dessus
du niveau hasard/majorité.

\section{Modules complémentaires}

La fiche mentionnait aussi l'auto-supervision, un prétrain de type fondation,
la calibration en ligne, un peu de renforcement, l'anticipation d'un déséquilibre
et des activités domestiques comme la préparation d'un repas. Voici ce que
nous avons réellement ajouté, avec les limites qui vont avec.

\subsection{Pré-entraînement auto-supervisé (SSL)}

Avant le fine-tuning supervisé, un pré-entraînement SimCLR-style sur fenêtres
IMU (sans labels) produit deux vues augmentées (bruit gaussien, mise à
l'échelle canal, masquage temporel) et minimise une perte NT-Xent sur une
tête de projection. Le backbone obtenu (\texttt{ssl\_backbone.pt}) sert
d'initialisation optionnelle. Ce n'est pas un SSL massif type industrie
\parencite{yuan2024}, mais une étape pragmatique compatible avec nos volumes
de données.

\subsection{Pipeline foundation-style}

Par \emph{foundation-style}, nous entendons le motif
prétrain générique $\rightarrow$ adaptation ciblée, et non l'intégration d'un
modèle de fondation public (TimesNet, UniTS, MOMENT). Le script
\texttt{train\_foundation\_style.py} enchaîne SSL puis fine-tune supervisé
sur HAPT/WISDM et produit \texttt{foundation\_style.pt}. Ce checkpoint
démontre le pattern ; le checkpoint de production reste
\texttt{pretrained.pt} (supervisé direct).

\subsection{Calibration de confiance en ligne}

Outre la mise à jour des prototypes, un module
\texttt{OnlineCalibrator} combine :
\begin{itemize}
  \item \texttt{OnlineTemperatureScaler} : un scalaire $T$ appris en ligne
  (NLL sur un tampon étiqueté) pour recalibrer $\mathrm{softmax}(\ell/T)$ ;
  \item \texttt{AdaptiveConfidenceThreshold} : seuil d'acceptation / rejet
  ajusté dynamiquement selon la fiabilité observée.
\end{itemize}
Ce mécanisme répond à l'exigence fiche de \emph{calibration adaptative en ligne}
des scores d'alerte et de classification.

\subsection{Bandit de seuil (RL léger)}

La fiche évoque l'apprentissage séquentiel par renforcement. Nous
implémentons un bandit $\varepsilon$-greedy (\texttt{ThresholdBandit}) sur une
grille de seuils de confiance d'alerte. La récompense pénalise les fausses
alertes et récompense les vraies détections. Ce n'est pas du deep RL ; c'est
une adaptation en ligne du seuil d'alerte, suffisante pour la démonstration
et le couplage avec la calibration.

\subsection{Anticipation du risque de chute / équilibre}

Au-delà de la détection d'impact a posteriori, une tête binaire
\texttt{FallRiskHead} (LSTM sur séquence d'embeddings) prédit si la
\emph{prochaine} activité est une transition posturale à risque
(labels 9--12 : stand$\leftrightarrow$sit, sit$\leftrightarrow$lie). Un
détecteur physique heuristique (\texttt{PhysicalFallDetector}) complète la
couche sécurité. La classe unifiée \texttt{fall} étant absente du merge
HAPT+WISDM, nous ne prétendons pas détecter des chutes cliniques réelles ;
MobiAct reste une perspective.

\subsection{Scénario « préparation de repas »}

La fiche cite la préparation d'un repas comme exemple d'activité. Aucun
label cuisine n'existe dans HAPT/WISDM. Le module
\texttt{adl\_scenarios.py} documente le scénario comme \emph{unavailable}
et expose un proxy narratif UI (transitions domestiques) uniquement pour
la démo Streamlit — à ne pas présenter comme un classifieur de repas.

\section{Scénarios d'utilisation et flux de données}

Cette section illustre comment les modules interagissent dans trois cas d'usage
concrets. Les diagrammes reprennent l'architecture figure TikZ du §4.5.

\subsection{Cas 1 — Personnalisation progressive (user-incrémental)}

Un service de coaching sportif déploie l'application sur le smartphone de nouveaux
utilisateurs chaque semaine. Au jour~0, le backbone pré-entraîné sur HAPT classifie
correctement les activités génériques (marche, assis). À l'arrivée de l'utilisateur
$u_k$, l'application collecte une vingtaine de minutes de données étiquetées (ou
pseudo-étiquetées par interaction). Le module UWR fine-tune le modèle sur $u_k$
pendant quelques dizaines d'époques (typiquement 10 à 50 selon la configuration
expérimentale), remplit le tampon avec les fenêtres les plus incertaines, et
met à jour les prototypes. L'inférence pour $u_k$ utilise le NCM sur prototypes ;
les utilisateurs $u_1..u_{k-1}$ restent couverts grâce au rejeu. Métrique suivie :
macro-F1 par sujet et forgetting (chapitre~6).

\subsection{Cas 2 — Extension du répertoire d'activités (class-incrémental)}

Un programme de rééducation ajoute de nouveaux exercices au fil des semaines
(ex.\ équilibre sur une jambe, montée d'escaliers guidée). Chaque semaine introduit
deux nouvelles classes. Le classifieur doit reconnaître l'ensemble des exercices
cumulés sans oublier les premiers. Le tampon stratifié privilégie les exemplaires
par classe ; la taille (2\,000 vs 5\,000) conditionne directement le forgetting
(22~\% de réduction observée). Les transitions posturales HAPT modélisent ce
déséquilibre extrême entre classes massives et classes rares.

\subsection{Cas 3 — Déploiement cross-appareil (DomainAdapter + anticipation)}

Un patient quitte l'hôpital (capteur ceinture, protocole HAPT-like) et continue
le suivi à domicile avec son téléphone en poche (distribution WISDM-like). Un
\texttt{DomainAdapter} few-shot sur quelques fenêtres étiquetées à domicile
calibre les embeddings avant inférence. En parallèle, le module d'anticipation LSTM peut signaler
une activité suivante probable lorsque le contexte partiel suggère une
transition ; le détecteur hybride de chutes fusionne règles physiques et MLP.
Ce scénario combine adaptation domaine, CL utilisateur si le patient est suivi
longitudinalement, et anticipation / sécurité.

\subsection{Contraintes temps réel rappelées}

Dans les trois cas, une fenêtre de 3~s doit idéalement être inférée en moins de
150~ms sur CPU (\emph{objectif} de conception rappelé au §4.2, non mesuré comme
benchmark dans nos expériences). Le backbone $\approx$531~K paramètres et la
tête d'anticipation légère sont dimensionnés pour rester compatibles avec cette
cible ; le recalcul d'entropie UWR est périodique (\texttt{update\_freq}), pas à chaque inférence utilisateur en production.

\section{Protocole d'évaluation}

\subsection{Baselines comparées}

\begin{table}[H]
\centering
\small
\begin{tabularx}{\textwidth}{|Y|Y|Y|}
\hline
\textbf{Baseline} & \textbf{Description} & \textbf{Référence} \\
\hline
Naïf (Finetuning) & Entraînement séquentiel sans mécanisme anti-oubli & — \\
EWC & Régularisation élastique des poids critiques & \parencite{kirkpatrick2017} \\
iCaRL & Replay + exemplaires + classification NCM & \parencite{rebuffi2017} \\
\textbf{UWR (notre méthode)} & Replay pondéré par incertitude + prototypes + contrastif & Ce travail \\
\hline
\end{tabularx}
\caption{Baselines — apprentissage continu}
\label{tab:baselines_cl}
\end{table}

Toutes les méthodes utilisent le même backbone pré-entraîné pour garantir une comparaison équitable. EWC et iCaRL sont réimplémentés en PyTorch et adaptés au contexte HAR.

\begin{table}[H]
\centering
\small
\begin{tabularx}{0.88\textwidth}{|Y|c|}
\hline
\textbf{Méthode d'anticipation (livrée)} & \textbf{Macro-F1 (ordre de grandeur)} \\
\hline
LSTM multi-classe, fenêtres partielles,
\texttt{transitions\_only} (notre méthode) & $\approx$ 0,59--0,64 \\
\hline
\end{tabularx}
\caption{Baseline anticipation — formulation implémentée dans \texttt{har\_project}}
\label{tab:baselines_anticipation}
\end{table}

\begin{table}[H]
\centering
\small
\begin{tabularx}{\textwidth}{|Y|Y|}
\hline
\textbf{Méthode CORAL} & \textbf{Description} \\
\hline
Sans adaptation & Évaluation directe sur WISDM du modèle entraîné sur HAPT \\
CORAL non supervisé & Alignement de covariance sans étiquettes cibles \\
CORAL supervisé (20~\%) & Alignement + fine-tuning avec 20~\% d'étiquettes cibles \\
\hline
\end{tabularx}
\caption{Baselines — adaptation cross-dataset HAPT $\rightarrow$ WISDM}
\label{tab:baselines_coral}
\end{table}

\subsection{Métriques retenues}

Les métriques d'évaluation sont choisies pour couvrir à la fois la performance instantanée et la dynamique d'apprentissage continu :

Accuracy (ACC) : taux de classification correcte global, calculé sur l'ensemble des tâches vues après la dernière mise à jour.

Macro-F1 : moyenne non pondérée des F1-scores par classe, robuste au déséquilibre de classes.

Forgetting (F) : mesure de la dégradation des performances sur les tâches passées (Chaudhry et al., 2018) :

$$F = \frac{1}{T-1} \sum_{i=1}^{T-1} \max_{t \in \{i,\ldots,T\}} a_{t,i} - a_{T,i}$$

où $a_{t,i}$ est l'accuracy sur la tâche $i$ après l'entraînement sur la tâche $t$.

Backward Transfer (BWT) : version signée du forgetting, permettant de détecter également les cas de transfert positif rétroactif ($BWT > 0$).

AUC-ROC : aire sous la courbe ROC, utilisée pour la détection de chutes et l'anticipation, insensible au seuil de décision.

\section{Module auxiliaire — risque de chute et alerte}

Pour l'alerte préventive, on combine trois éléments. Une tête LSTM binaire
(\texttt{FallRiskHead}) indique si la prochaine activité est une transition
posturale à risque (labels 9--12). Une heuristique physique
(\texttt{PhysicalFallDetector}) repère un pic d'accélération suivi d'une
phase calme. Enfin, un bandit de seuil et une calibration de température
ajustent l'alerte en ligne.

\subsection{Métriques cibles}

Macro-F1 binaire et accuracy pour le fall-risk ; objectif clinique sur chutes
réelles (MobiAct) reporté. Résultat actuel : macro-F1 binaire $\approx$ 0,88
sur transitions HAPT ($p\in\{50\%,75\%\}$) — preuve de concept, pas validation
clinique. L'AUC-ROC 0,762 du détecteur hybride historique (proxy impact)
reste documenté comme baseline complémentaire.

\section{Conclusion}

En résumé, le système repose sur un Transformer partagé, UWR pour l'apprentissage
continu, et une anticipation LSTM sur fenêtres partielles. Autour, SSL léger,
adaptateur de domaine, calibration, bandit de seuil et fall-risk complètent
le dispositif sans le surcharger. Le chapitre suivant décrit comment tout
cela est codé et entraîné.

% ===== END include: 09-Conception.tex =====


% ===== BEGIN include: 10-Realisation.tex =====
\clearpage
\chapter{Réalisation de la solution}
\label{ch:10-Realisation}

\section{Introduction}

Ce chapitre décrit l'implémentation concrète du système conçu au chapitre précédent. Nous passons de la conception théorique à la réalisation logicielle, en détaillant les choix techniques effectués à chaque étape : environnement de développement, préparation des données, implémentation du backbone Transformer, des modules de reconnaissance continue et d'anticipation, et des contributions supplémentaires (détection de chutes et adaptation CORAL). Chaque section s'appuie sur le code source produit, en en explicitant la logique, les paramètres et les justifications des choix réalisés. L'ensemble du projet est organisé en modules Python réutilisables, documentés et versionnés.

\section{Environnement de développement}

\subsection{Plateformes utilisées}

Le développement et les expériences ont été conduits sur une machine Apple MacBook Pro équipée d'un processeur Apple M-series (architecture ARM), disposant de 16 Go de mémoire unifiée. L'accélération matérielle est assurée par le backend MPS (Metal Performance Shaders) de PyTorch, permettant d'exploiter le GPU Apple Silicon sans recourir à CUDA. Pour les expériences plus longues, l'environnement Google Colab (GPU Tesla T4) a également été utilisé.

Le système d'exploitation est macOS Sequoia. L'environnement Python est géré par Homebrew avec Python 3.14, et les dépendances sont isolées dans un environnement virtuel dédié au projet.

\subsection{Langages et frameworks}

Python 3.14 est le langage principal du projet. PyTorch 2.x constitue le framework de deep learning, choisi pour sa flexibilité dans la définition d'architectures personnalisées et sa compatibilité native avec le backend MPS. Le calcul scientifique repose sur NumPy pour les opérations matricielles et SciPy pour le prétraitement du signal (filtrage Butterworth, rééchantillonnage).

L'évaluation et la comparaison des modèles utilisent scikit-learn pour le calcul des métriques (F1-score, accuracy, AUC-ROC, classification\_report). La visualisation des résultats et des embeddings fait appel à Matplotlib et Seaborn.

\subsection{Bibliothèques et outils}

\begin{table}[H]
\centering
\small
\begin{tabularx}{0.88\textwidth}{|Y|c|Y|}
\hline
\textbf{Bibliothèque} & \textbf{Version} & \textbf{Rôle} \\
\hline
torch          & 2.x   & Framework deep learning \\
numpy          & 1.26+ & Calcul numérique \\
scipy          & 1.12+ & Traitement du signal \\
scikit-learn   & 1.4+  & Métriques et évaluation \\
matplotlib     & 3.8+  & Visualisation \\
seaborn        & 0.13+ & Visualisation avancée \\
tqdm           & 4.x   & Barres de progression \\
\hline
\end{tabularx}
\caption{Principales dépendances Python (\texttt{requirements.txt})}
\label{tab:dependencies}
\end{table}

\subsection{Structure du projet et suivi de version}

\begin{lstlisting}[language=bash,caption={Structure du projet \texttt{har\_project/}},label=lst:project_tree]
har_project/
|-- src/
|   |-- data/
|   |   |-- preprocessing.py       # pipeline pretraitement
|   |   |-- har_dataset.py         # Dataset PyTorch
|   |   `-- anticipation_dataset.py
|   |-- models/
|   |   |-- backbone.py            # IMUTransformerEncoder
|   |   |-- har_model.py           # modele complet deux tetes
|   |   |-- replay_buffer.py       # tampon de rejeu standard
|   |   |-- uncertainty_replay.py  # UWR (entropie softmax)
|   |   |-- prototype_memory.py    # memoire de prototypes
|   |   |-- fall_risk.py           # fall-risk + heuristique physique
|   |   |-- online_calibration.py  # temperature + seuil
|   |   `-- domain_adapter.py      # DomainAdapter (affine few-shot)
|   |-- training/
|   |   |-- trainer.py             # boucle pre-entrainement
|   |   |-- anticipation_trainer.py
|   |   |-- ssl_pretrain.py        # SimCLR-style IMU
|   |   |-- fall_risk_trainer.py
|   |   `-- rl_threshold_agent.py  # bandit de seuil
|   `-- scenarios/
|       `-- adl_scenarios.py       # repas (indisponible) + proxy UI
|-- scripts/                       # scripts executables
|-- checkpoints/                   # modeles sauvegardes
|-- data/
|   |-- raw/                       # donnees brutes
|   `-- processed/                 # donnees pretraitees
`-- docs/                          # memoire LaTeX
\end{lstlisting}

Le suivi de version est géré par Git. Chaque composant fonctionnel (backbone, UWR, DomainAdapter, anticipation) fait l'objet d'une implémentation et d'une validation indépendantes avant intégration dans le modèle complet.

\section{Préparation des données}

\subsection{Chargement et homogénéisation des datasets}

Le chargement des données brutes est implémenté dans le module src/data/dataset\_loaders.py. Pour chaque dataset, un chargeur spécifique lit les fichiers sources, applique les conversions d'unités nécessaires et produit un tableau NumPy brut $(N_{sujets}, T_{total}, C)$.

Chargement HAPT. Les données HAPT sont distribuées sous forme de fichiers texte séparés pour les signaux d'accélération et de gyroscope, avec des annotations temporelles d'activité. Le chargeur lit les 30 fichiers sujets, concatène les canaux accéléromètre (3 canaux, unité g) et gyroscope (3 canaux, unité rad/s), et charge les étiquettes échantillon-par-échantillon depuis les fichiers de labels.

Chargement WISDM. Le dataset WISDM est distribué sous forme d'un fichier CSV unique. Chaque ligne correspond à un échantillon avec les colonnes : identifiant sujet, étiquette d'activité, timestamp, $a_x$, $a_y$, $a_z$. Le chargeur regroupe les lignes par sujet et activité, et comble les éventuelles discontinuités temporelles détectées par des sauts dans les timestamps.

Homogénéisation. Les deux datasets sont ramenés à une représentation commune $(N, T=150, C=6)$ par application du pipeline décrit en §5.3.2. Les étiquettes sont remappées vers un espace commun entier ${0, 1, \ldots, n_classes-1}$. Dans les expériences multi-dataset, un dictionnaire de correspondance de classes est appliqué manuellement.

\subsection{Implémentation du pipeline de prétraitement}

\begin{lstlisting}[language=Python,caption={Constantes globales du prétraitement (\texttt{preprocessing.py})},label=lst:preproc_const]
TARGET_HZ   = 50      # fréquence cible (Hz)
WINDOW_SIZE = 150     # 3 secondes × 50 Hz
STEP_SIZE   = 75      # chevauchement 50 %
OVERLAP     = 0.5
G           = 9.80665 # m/s²
\end{lstlisting}

Rééchantillonnage. La fonction \texttt{resample\_signal(signal, original\_hz)} utilise \texttt{scipy.signal.resample} pour interpoler ou sous-échantillonner le signal vers \texttt{TARGET\_HZ = 50} Hz.

\begin{lstlisting}[language=Python,caption={Rééchantillonnage à 50 Hz},label=lst:resample]
def resample_signal(signal, original_hz):
    if original_hz == TARGET_HZ:
        return signal
    n_new = int(signal.shape[0] * TARGET_HZ / original_hz)
    return resample(signal, n_new, axis=0)
\end{lstlisting}

Suppression de la gravité. La fonction \texttt{remove\_gravity(signal, fs, cutoff=0.5)} implémente un filtre passe-bas Butterworth d'ordre 5 à fréquence de coupure 0,5 Hz.

\begin{lstlisting}[language=Python,caption={Suppression de la composante gravitationnelle},label=lst:gravity]
def remove_gravity(signal, fs=TARGET_HZ, cutoff=0.5):
    b, a = butter(5, cutoff / (fs / 2.0), btype='low', analog=False)
    gravity = filtfilt(b, a, signal, axis=0)
    return signal - gravity
\end{lstlisting}

Segmentation par fenêtre glissante. La fonction sliding\_windows(signal) extrait des fenêtres de WINDOW\_SIZE = 150 échantillons avec un stride de STEP\_SIZE = 75 (chevauchement 50 \%). Pour chaque fenêtre, l'étiquette majoritaire est assignée via numpy.bincount. La fonction sliding\_windows\_with\_labels(signal, labels) est utilisée quand les annotations sont disponibles au niveau échantillon.

Normalisation z-score. La fonction normalize\_signal(signal, mean, std) normalise chaque canal indépendamment. Les statistiques (moyenne et écart-type) sont calculées sur les données d'entraînement du sujet courant et réutilisées pour normaliser ses données de test, évitant toute fuite d'information.

Pipeline complet. La fonction preprocess\_signal(signal, original\_hz, unit) enchaîne les quatre étapes ci-dessus et produit directement les fenêtres normalisées prêtes à être injectées dans le modèle.

\subsection{Construction du flux incrémental simulé}

Pour simuler un flux de données réaliste, les données prétraitées sont réorganisées selon les scénarios d'apprentissage définis en §4.2.1. Cette logique est implémentée dans src/data/har\_dataset.py.

Scénario user-incrémental. Les fenêtres sont partitionnées par sujet via
\texttt{subjects.npy}. L'ordre des tâches suit les identifiants sujets
triés (\texttt{user\_incremental\_tasks} : un sujet = une tâche), pas une
permutation aléatoire.

Scénario class-incrémental. Les classes sont regroupées par paquets de
\texttt{classes\_per\_task=2} (défaut de \texttt{train.py}), avec mélange de
l'ordre des classes sous graine fixe (\texttt{seed=42}).

Dataset. La classe \texttt{HARDataset} encapsule les tableaux NumPy
$(X, y, \mathrm{subjects}, \mathrm{origins})$. L'augmentation n'est pas un
flag du Dataset : elle est appliquée dans la boucle d'entraînement / module
les scripts d'entraînement / SSL (\texttt{imu\_augment} dans \texttt{ssl\_pretrain.py}) lorsque l'augmentation est activée.

\section{Implémentation du backbone Transformer}

\subsection{Architecture détaillée}

Le backbone est implémenté dans src/models/backbone.py. La classe principale IMUTransformerEncoder hérite de nn.Module et est construite par composition de trois sous-modules : PositionalEncoding, une liste de TransformerBlock, et une couche de normalisation finale.

Encodage positionnel. La classe PositionalEncoding calcule à l'initialisation les codes sinusoïdaux fixes selon la formule standard :

$$PE_{pos, 2i} = \sin\left(\frac{pos}{10000^{2i/d}}\right), \quad PE_{pos, 2i+1} = \cos\left(\frac{pos}{10000^{2i/d}}\right)$$

Ces codes sont stockés comme buffer non-paramètre (non mis à jour par le gradient). Un paramètre scalaire self.scale appris, initialisé à 1, module l'amplitude des codes positionnels. Le dropout est appliqué après addition pour régularisation.

Blocs Transformer. Chaque TransformerBlock implémente l'architecture pré-norme (pre-LN), plus stable que la post-norme pour les petits datasets. L'ordre des opérations est : LayerNorm → MultiheadAttention → connexion résiduelle → LayerNorm → FFN → connexion résiduelle. L'attention multi-têtes est implémentée via nn.MultiheadAttention(d\_model=128, num\_heads=4, batch\_first=True). La FFN est composée de Linear(128, 256) → GELU → Dropout(0.1) → Linear(256, 128) → Dropout(0.1).

Le passage forward du backbone effectue les opérations suivantes :

\begin{lstlisting}[language=Python,caption={Forward pass du backbone},label=lst:backbone_forward]
def forward(self, x):          # x: (B, T, 6)
    B = x.size(0)
    x = self.input_proj(x)     # (B, T, 128)
    cls = self.cls_token.expand(B, -1, -1)
    x = torch.cat([cls, x], dim=1)  # (B, T+1, 128)
    x = self.pos_enc(x)
    for block in self.blocks:
        x = block(x)
    x = self.norm(x)
    return x[:, 0]             # (B, 128) — token CLS
\end{lstlisting}

Le token CLS est initialisé avec des valeurs proches de zéro (trunc\_normal\_(std=0.02)) et évolue librement pendant l'entraînement pour apprendre à agréger l'information globale de la fenêtre par le mécanisme d'attention.

Initialisation des poids. Toutes les couches linéaires sont initialisées par la méthode trunc\_normal\_(std=0.02), standard pour les Transformers, avec les biais initialisés à zéro.

Fonction utilitaire. La fonction build\_backbone(cfg=None) permet de créer un backbone avec des hyperparamètres personnalisés ou les valeurs par défaut :

defaults = dict(in\_channels=6, d\_model=128, n\_heads=4,

n\_blocks=4, ffn\_dim=256, dropout=0.1)

\subsection{Pré-entraînement initial}

Le pré-entraînement du backbone est réalisé par le script scripts/train.py via la classe BaselineTrainer du module src/training/baseline\_trainer.py. Le modèle complet HARContinualModel (backbone + tête de classification HAR) est entraîné en mode supervisé standard sur les données HAPT.

Optimiseur. AdamW avec taux d'apprentissage $lr = 10^{-3}$ et pénalité de régularisation $\lambda_{wd} = 10^{-4}$. Un scheduler cosinus (CosineAnnealingLR) réduit le taux d'apprentissage de $lr$ à $lr \times 0{,}01$ sur 100 epochs.

Perte. Cross-entropie standard sur les 12 classes HAPT, sans pondération de classes. Un mixup léger (alpha=0.2) est appliqué pour régularisation.

Split train/validation. 80 \% des sujets pour l'entraînement (24 sujets), 20 \% pour la validation (6 sujets), avec séparation stricte par sujet pour éviter la fuite de données inter-sujets.

Augmentation. Lors du pré-entraînement, l'augmentation de fenêtres IMU est activée avec probabilité $p_{\mathrm{apply}} = 0{,}5$ par transformation : bruit gaussien ($\sigma = 0{,}01$), mise à l'échelle ($s \in [0{,}85, 1{,}15]$), déformation temporelle ($\sigma_{\mathrm{warp}} = 0{,}05$), dropout de canaux ($p_{\mathrm{drop}} = 0{,}1$), et découpage de fenêtre ($r_{\mathrm{crop}} = 0{,}9$). L'augmentation permet d'améliorer l'accuracy de validation de 81~\% (sans) à 91,3~\% (avec), soit un gain de +10 points.

\begin{figure}[H]
  \centering
  \insertfig[max width=0.88\linewidth]{pretrain_history.png}{har\_project/results/figures/}{Courbes d'apprentissage}
  \caption[Courbes de pré-entraînement]{Courbes de loss et macro-F1 lors du pré-entraînement
  supervisé sur HAPT.}
  \label{fig:pretrain_impl}
\end{figure}

Sauvegarde. Le meilleur checkpoint (selon le F1 macro de validation) est sauvegardé dans checkpoints/pretrained\_backbone.pt avec la structure :

torch.save(\{

"model\_state":      model.state\_dict(),

"prototype\_state":  model.har\_head.prototype\_memory.state(),

"val\_acc":          best\_acc,

"val\_f1":           best\_f1,

\}, path)

\subsection{Visualisation des représentations apprises}

Pour vérifier la qualité des représentations apprises, une projection t-SNE est calculée sur les embeddings du backbone sur 1 000 fenêtres test tirées aléatoirement. Les embeddings 128-dimensionnels sont réduits à 2 dimensions par t-SNE (perplexité=30, itérations=1000). La visualisation confirme une séparation nette des clusters par classe d'activité, avec les activités dynamiques (marche, jogging) bien séparées des activités statiques (assis, debout, allongé). Les transitions posturales se placent naturellement entre les paires d'activités statiques correspondantes, témoignant d'une représentation géométriquement cohérente.

\begin{figure}[H]
  \centering
  \insertfig[max width=0.82\linewidth]{tsne_embeddings.png}{har\_project/results/figures/}{Projection t-SNE}
  \caption[Projection t-SNE des embeddings]{Projection t-SNE des embeddings CLS — séparation
  par classe d'activité.}
  \label{fig:tsne_impl}
\end{figure}

\section{Implémentation du module de reconnaissance continue}

\subsection{Gestion des prototypes}

La mémoire de prototypes est implémentée dans src/models/prototype\_memory.py via la classe PrototypeMemory. Elle maintient un dictionnaire \{class\_id → prototype\_vector\} et un compteur de mises à jour par classe.

Mise à jour. La méthode update(embeddings, labels) calcule l'embedding moyen par classe dans le lot courant et met à jour les prototypes par moyenne exponentielle mobile :

def update(self, embeddings, labels):

for c in labels.unique():

mask = (labels == c)

batch\_mean = embeddings[mask].mean(dim=0)

if c.item() in self.prototypes:

self.prototypes[c.item()] = (

self.momentum * self.prototypes[c.item()]

+ (1 - self.momentum) * batch\_mean.detach()

)

else:

self.prototypes[c.item()] = batch\_mean.detach().clone()

Inférence. La méthode predict(embeddings) calcule la distance euclidienne de chaque embedding à tous les prototypes connus et retourne la classe du prototype le plus proche. Si aucun prototype n'est encore enregistré pour une classe, celle-ci est exclue de la prédiction.

Perte contrastive. La méthode contrastive\_loss(embeddings, labels, temperature=0.07) implémente une perte de séparation inter-classes en espace de prototypes. Pour chaque exemple, la perte est calculée comme le log-softmax négatif de la similarité cosinus normalisée par la température, maximisant la proximité au prototype de la classe correcte et minimisant celle aux autres prototypes.

\subsection{Tampon de rejeu et Uncertainty-Weighted Replay}

Deux implémentations du tampon de rejeu sont disponibles, selon une interface commune définie dans src/models/replay\_buffer.py.

Tampon standard (ReplayBuffer). La classe ReplayBuffer maintient deux tableaux NumPy circulaires buffer\_X et buffer\_y de capacité capacity=2000. La méthode add\_batch(X, y) insère les nouveaux exemples selon la stratégie de réservoir. La méthode sample(batch\_size, device) retourne un lot aléatoire sous forme de tenseurs PyTorch sur le dispositif spécifié.

\begin{lstlisting}[language=Python,caption={Échantillonnage UWR depuis le tampon},label=lst:uwr_sample]
def sample_uncertain(self, model, batch_size, device):
    scores = self.priority_scores[:self.size] + self.epsilon
    probs  = scores / scores.sum()
    indices = np.random.choice(self.size, size=batch_size,
                               replace=False, p=probs)
    X = torch.from_numpy(self.buffer_X[indices]).to(device)
    y = torch.from_numpy(self.buffer_y[indices]).to(device)
    return X, y
\end{lstlisting}

Le paramètre alpha contrôle la force de la priorisation : $\alpha = 1{,}0$ correspond à un échantillonnage purement proportionnel aux incertitudes, tandis que $\alpha \to 0$ converge vers l'échantillonnage uniforme.

\subsection{Calcul d'incertitude UWR (entropie softmax)}

Le module \texttt{UncertaintyWeightedReplayBuffer} (\texttt{src/models/uncertainty\_replay.py}) calcule l'entropie sur le tampon en mode \texttt{eval} (une passe forward), puis échantillonne avec $P(x)\propto \mathcal{U}(x)^{\alpha}$.

\begin{lstlisting}[language=Python,caption={Entropie de prédiction pour UWR},label=lst:uwr_entropy]
@torch.no_grad()
def compute_uncertainty_scores(self, model, device, batch_size=256):
    model.eval()
    entropies = []
    for start in range(0, len(self._X), batch_size):
        batch = torch.from_numpy(self._X[start:start+batch_size]).to(device)
        probs = F.softmax(model(batch), dim=-1)
        H = -(probs * (probs + 1e-10).log()).sum(dim=-1)
        entropies.append(H.cpu().numpy())
    scores = np.concatenate(entropies) ** self.alpha
    return scores / scores.sum()
\end{lstlisting}

Il n'existe pas de fichier \texttt{mc\_dropout.py} dans le dépôt livré : le MC Dropout reste une référence bibliographique, pas l'implémentation utilisée pour pondérer le rejeu.

\subsection{Entraînement incrémental complet}

Le modèle complet HARContinualModel est implémenté dans src/models/har\_model.py. Il encapsule le backbone, la tête HAR, la mémoire de prototypes et le tampon de rejeu dans une interface unifiée.

\begin{lstlisting}[language=Python,caption={Pas d'apprentissage continu (\texttt{continual\_step})},label=lst:continual_step]
def continual_step(self, x, y, optimizer, replay_batch_size=32, device=None):
    self.train()
    optimizer.zero_grad()
    emb_curr    = self.backbone(x)
    logits_curr = self.har_head(emb_curr)
    loss_ce     = F.cross_entropy(logits_curr, y)
    loss_replay = torch.tensor(0.0, device=device)
    if len(buffer) >= replay_batch_size:
        rx, ry        = buffer.sample_uncertain(self, replay_batch_size, device)
        emb_replay    = self.backbone(rx)
        logits_replay = self.har_head(emb_replay)
        loss_replay   = F.cross_entropy(logits_replay, ry)
    loss_contrast = self.har_head.contrastive_loss(emb_curr, y)
    total_loss = loss_ce + loss_replay + 0.1 * loss_contrast
    total_loss.backward()
    optimizer.step()
    self.har_head.update_prototypes(emb_curr.detach(), y)
    buffer.add_batch(x.cpu().numpy(), y.cpu().numpy())
    return {"loss_ce": ..., "loss_replay": ..., "loss_contrast": ...}
\end{lstlisting}

Ce design garantit que chaque pas d'apprentissage met simultanément à jour les paramètres du modèle, les prototypes et le contenu du tampon de rejeu, maintenant une cohérence globale entre tous les composants.

\subsection{Orchestration des scripts d'expérience}

Chaque protocole du chapitre~6 correspond à un script exécutable dans
\texttt{har\_project/scripts/}. Le tableau~\ref{tab:scripts_exp} résume la
chaîne de dépendances.

\begin{table}[H]
\centering
\small
\begin{tabularx}{\textwidth}{|>{\raggedright\arraybackslash}p{4.5cm}|Y|>{\raggedright\arraybackslash}p{3.2cm}|}
\hline
\textbf{Script} & \textbf{Rôle} & \textbf{Sortie principale} \\
\hline
\texttt{preprocess.py} & Charge / fenêtre / normalise les jeux & \texttt{data/processed/*.npz} \\
\texttt{train.py --mode pretrain} & Pré-train backbone + tête CE & \texttt{checkpoints/*.pt} \\
\texttt{train.py --mode continual --scenario user} & Boucle user-incrémental UWR & Matrice $R[i,j]$, métriques CL \\
\texttt{train.py --mode continual --scenario class} & Class-incrémental & Figures forgetting class \\
\texttt{train\_anticipation.py} & Tête LSTM multi-classe (backbone gelé) & F1 / acc. par ratio $p$ \\
\texttt{train\_anticipation\_joint.py} & Variante joint backbone+LSTM & Checkpoint anticipation \\
\texttt{cross\_dataset\_eval.py} & Adaptation / eval HAPT$\rightarrow$WISDM & F1 cross-dataset \\
\texttt{regenerate\_figures.sh} & Regénère les figures & \texttt{results/figures/*.png} \\
\hline
\end{tabularx}
\caption{Scripts d'expérience et artefacts produits}
\label{tab:scripts_exp}
\end{table}

Le script user-incrémental charge le checkpoint pré-entraîné, initialise un
\texttt{UncertaintyWeightedReplayBuffer(2000)} et une \texttt{PrototypeMemory},
puis itère sur la liste ordonnée des sujets. Après chaque tâche, il sérialise
la matrice $R$ et les courbes d'oubli au format NumPy/PNG. La reprise sur
erreur est possible en passant \texttt{--start\_task k} pour éviter de
réentraîner les $k-1$ premières tâches — utile lors du développement.

Le logging console affiche, par époque : loss CE, loss replay, loss contrastive,
macro-F1 sur le sujet courant et sur un sous-échantillon de sujets passés (max
5 pour limiter le temps). Les checkpoints intermédiaires sont sauvegardés tous
les 10 sujets.

\subsection{Validation unitaire et tests de non-régression}

Avant l'intégration, chaque module a été testé isolément : \texttt{backbone.py}
(forward shape $(B,128)$), \texttt{replay\_buffer.py} (capacité circulaire,
stratégie réservoir), \texttt{uncertainty\_replay.py} (somme des probabilités
d'échantillonnage = 1), \texttt{domain\_adapter.py} (symétrie CORAL source/cible).
Un jeu synthétique de 100 fenêtres aléatoires permet de vérifier en $<$30~s que
la chaîne complète \texttt{continual\_step} ne diverge pas numériquement (NaN).

\section{Implémentation du module d'anticipation}

\subsection{Architecture LSTM multi-classe}

Le module d'anticipation est la classe \texttt{AnticipationHead}
(\texttt{src/models/har\_model.py}), entraînée via
\texttt{scripts/train\_anticipation.py} (backbone gelé) et
\texttt{scripts/train\_anticipation\_joint.py} (joint optionnel).

\begin{verbatim}
self.lstm = nn.LSTM(input_size=d_model, hidden_size=128,
                    num_layers=2, batch_first=True, dropout=0.2)
self.classifier = nn.Sequential(
    nn.LayerNorm(128),
    nn.Linear(128, n_classes),
)
\end{verbatim}

Le LSTM est unidirectionnel. \texttt{anticipate()} aplatit
$(B,S,T_{\mathrm{partial}},C)$, encode chaque fenêtre, reforme $(B,S,d)$
puis appelle le LSTM.

\subsection{Dataset et entraînement}

\texttt{AnticipationDataset} : contexte $S=5$ fenêtres tronquées à $p\,\%$,
cible = label de la fenêtre suivante, filtre \texttt{transitions\_only}.
Le split train/val (80/20) est obtenu après mélange des fenêtres
(\texttt{build\_anticipation\_datasets}).

Dans \texttt{anticipation\_trainer.py} : backbone + HAR gelés ; AdamW sur la
tête LSTM ; perte \texttt{cross\_entropy} ; métriques accuracy et macro-F1 ;
un entraînement par ratio $p$.

\begin{table}[H]
\centering
\small
\begin{tabularx}{0.88\textwidth}{|Y|c|c|}
\hline
\textbf{Configuration} & \textbf{Macro-F1} & \textbf{Commentaire} \\
\hline
LSTM multi-classe, backbone gelé, $p=25\%$ & $\approx$ 0,59 & Fin de fenêtre + poids de classes \\
LSTM multi-classe, backbone gelé, $p=50\%$ & $\approx$ 0,63 & Idem \\
LSTM multi-classe, backbone gelé, $p=75\%$ & $\approx$ 0,64 & Checkpoint \texttt{anticipation.pt} \\
Baseline majorité & $\approx$ 0,04 & Toujours la classe la plus fréquente \\
\hline
\end{tabularx}
\caption{Résultats du module d'anticipation}
\label{tab:anticipation_impl}
\end{table}
\section{Contributions supplémentaires}

\subsection{Détection de chute — approche hybride / fall-risk}

La composante alerte / équilibre est réalisée par \texttt{FallRiskHead}
(\texttt{src/models/fall\_risk.py}, script \texttt{train\_fall\_risk.py}) :
prédiction \emph{binaire} indiquant si la prochaine activité est une transition
posturale à risque (labels 9--12 : stand$\leftrightarrow$sit, sit$\leftrightarrow$lie).
Un détecteur physique heuristique (\texttt{PhysicalFallDetector}) complète la
couche sécurité (pic d'accélération puis phase calme). La classe unifiée
\texttt{fall} est absente du merge HAPT+WISDM actuel.

Résultats. Sur validation, le macro-F1 binaire de la tête fall-risk atteint
environ 0,88 ($p\in\{50\%,75\%\}$). L'interface Streamlit et le bandit de
seuil (\texttt{ThresholdBandit}) illustrent l'adaptation en ligne du seuil
d'alerte (RL léger), couplée à une calibration de température
(\texttt{OnlineCalibrator}).

\subsection{SSL, foundation-style, calibration et RL léger}

Pour coller aux orientations de la fiche PFE, des modules légers ont été
ajoutés :
\begin{itemize}
  \item pré-entraînement auto-supervisé SimCLR-style
  (\texttt{scripts/train\_ssl.py} $\rightarrow$ \texttt{ssl\_backbone.pt}) :
  augmentations bruit / scale / mask + perte NT-Xent ;
  \item pipeline foundation-style SSL puis fine-tune
  (\texttt{train\_foundation\_style.py} $\rightarrow$ \texttt{foundation\_style.pt},
  run démonstratif, distinct du \texttt{pretrained.pt} de production) ;
  \item calibration online
  (\texttt{online\_calibration.py} : température + seuil adaptatif) ;
  \item bandit $\varepsilon$-greedy de seuil d'alerte
  (\texttt{rl\_threshold\_agent.py}).
\end{itemize}
Le scénario « préparation de repas » reste hors portée faute de corpus cuisine
(\texttt{adl\_scenarios.py}, statut \texttt{available=False}). L'interface
Streamlit expose six scénarios (reconnaissance, nouvel utilisateur,
anticipation, fall-risk, calibration+RL, SSL/foundation).

\subsection{Adaptateur de domaine (\texttt{DomainAdapter})}

L'adaptation de domaine livrée est dans \texttt{src/models/domain\_adapter.py} :
une transformation affine légère $x' = x \odot s + b$ (\texttt{DomainAdapter})
entraînée en few-shot pendant que le backbone reste gelé
(\texttt{adapt\_domain}), contre l'objectif plus proche prototype (même règle
qu'à l'inférence). Le script \texttt{cross\_dataset\_eval.py} mesure le
domain shift et applique cet adaptateur surtout sur PAMAP2 ($n$ fenêtres
cibles étiquetées).

CORAL \parencite{sun2016} (alignement de covariance) reste une référence
bibliographique pertinente, mais \emph{n'est pas} la fonction de perte
exécutée dans le dépôt actuel : il n'existe pas de \texttt{coral\_loss} ni de
classe \texttt{CORALAdapter} dans le code.

\begin{table}[H]
\centering
\small
\begin{tabularx}{0.88\textwidth}{|Y|c|}
\hline
\textbf{Transfert (checkpoint \texttt{cross\_dataset\_results.npy})} & \textbf{Macro-F1} \\
\hline
HAPT $\rightarrow$ HAPT (référence intra) & 0,838 \\
HAPT $\rightarrow$ WISDM (zéro-shot) & 0,464 \\
HAPT $\rightarrow$ PAMAP2 (zéro-shot) & 0,124 \\
\hline
\end{tabularx}
\caption{Transfert cross-dataset mesuré (sans réécrire le backbone)}
\label{tab:coral_impl}
\end{table}

Ces chiffres confirment un fort domain shift (surtout PAMAP2). L'adaptateur
affine few-shot vise à récupérer une partie de ce gap sans réentraînement
complet ; les anciens tableaux « CORAL 0,629 / $\times 15$ » ne correspondent
pas au code livré et ne sont plus rapportés comme résultat reproductible.

\section{Conclusion}

Le chapitre~4 est maintenant incarné dans des modules Python séparés : backbone,
UWR, anticipation, fall-risk, SSL, calibration, bandit et
\texttt{DomainAdapter}. Les runs confirment surtout UWR et le diagnostic
cross-dataset ; l'anticipation multi-classe se situe autour de 0,60--0,64 de
macro-F1. Le chapitre suivant regarde les chiffres de plus près.

\subsection{Retour d'expérience}

Quelques choix ont simplifié le travail. Une interface commune
\texttt{ReplayBuffer} / \texttt{UncertaintyWeightedReplayBuffer} a permis de
comparer les baselines sans tout réécrire. Geler le backbone au début des
expériences d'anticipation a isolé la tête LSTM. Un script de régénération
des figures évite les captures obsolètes. Sur Apple Silicon (MPS), le
pré-entraînement est nettement plus rapide qu'en CPU ; les runs à 44 tâches
restent longs (de l'ordre d'une nuit), d'où l'intérêt des checkpoints
intermédiaires.

% ===== END include: 10-Realisation.tex =====


% ===== BEGIN include: 11-Evaluation.tex =====
\clearpage
\chapter{Évaluation de la solution}
\label{ch:11-Evaluation}

\section{Introduction}

Ce chapitre présente l'évaluation expérimentale complète du système
\texttt{HARContinualModel} (\texttt{har\_project/}). Les protocoles reprennent
ceux définis au chapitre~4 : pré-entraînement supervisé, apprentissage
user- et class-incrémental, comparaison aux baselines (EWC, iCaRL), ablation,
adaptation cross-dataset (CORAL), anticipation et détection de chutes.

Nous structurons la lecture en trois blocs : (1)~qualité du backbone seul
(classification HAPT et comparaison SOTA) ; (2)~apprentissage continu avec
métriques forgetting/BWT sur protocoles 10 et 44 sujets ; (3)~modules auxiliaires
(anticipation LSTM multi-classe, CORAL, chutes). Les chiffres 10 sujets proviennent du
cahier des charges initial ; les runs 44 sujets reflètent l'état du dépôt
\texttt{har\_project} fusion multi-jeux — les deux coexistent volontairement avec
explicitation des protocoles dans chaque tableau.

\section{Résultats du pré-entraînement}

\subsection{Courbes d'apprentissage}

Le backbone \emph{IMUTransformerEncoder} ($d=128$, 4 blocs, 531\,457 paramètres)
est pré-entraîné sur HAPT avec AdamW ($\mathrm{lr}=10^{-3}$, scheduler cosinus,
30--100 époques). La figure~\ref{fig:pretrain} montre la convergence ; l'augmentation
de données apporte \textbf{+10,3 points} d'accuracy (81,0~\% $\rightarrow$ 91,3~\%).

\begin{figure}[H]
  \centering
  \insertfig[max width=0.88\linewidth]{pretrain_history.png}{har\_project/results/figures/}{Courbes loss et macro-F1}
  \caption[Courbes d'apprentissage du backbone]{Courbes d'apprentissage — loss et macro-F1
  sur la validation HAPT.}
  \label{fig:pretrain}
\end{figure}

\begin{table}[H]
\centering
\small
\begin{tabularx}{\textwidth}{|Y|c|c|c|}
\hline
\textbf{Configuration} & \textbf{Époques} & \textbf{Val. Accuracy} & \textbf{Val. F1 macro} \\
\hline
4 blocs, sans augmentation & 100 & 81,0~\% & — \\
4 blocs, avec augmentation  & 30  & 91,3~\% & 0,752 \\
6 blocs, avec augmentation  & 30  & 90,7~\% & 0,716 \\
\hline
\end{tabularx}
\caption{Impact de l'augmentation et de la profondeur du backbone}
\label{tab:pretrain_configs}
\end{table}

\subsection{Performances par classe}

\begin{table}[H]
\centering
\small
\begin{tabularx}{\textwidth}{|Y|c|c|c|}
\hline
\textbf{Classe} & \textbf{Précision} & \textbf{Rappel} & \textbf{F1} \\
\hline
Marche (WALKING)           & 0,96 & 0,97 & 0,97 \\
Montée escaliers           & 0,93 & 0,94 & 0,93 \\
Descente escaliers         & 0,91 & 0,90 & 0,90 \\
Debout / Assis / Allongé   & 0,92--0,99 & 0,91--0,99 & 0,92--0,99 \\
Transitions (6 classes)    & 0,72--0,88 & 0,65--0,84 & 0,68--0,86 \\
\hline
\textbf{Macro moyenne}     & 0,91 & 0,90 & 0,752 \\
\hline
\end{tabularx}
\caption{Performances par classe — validation HAPT}
\label{tab:per_class_pretrain}
\end{table}

\subsection{Visualisation des embeddings (t-SNE)}

\begin{figure}[H]
  \centering
  \insertfig[max width=0.82\linewidth]{tsne_embeddings.png}{har\_project/results/figures/}{Projection t-SNE}
  \caption[Projection t-SNE des embeddings CLS]{Projection t-SNE des embeddings CLS
  (1\,000 fenêtres de validation). Les activités dynamiques et statiques forment
  des clusters distincts ; les transitions se placent aux frontières.}
  \label{fig:tsne}
\end{figure}

\section{Comparaison avec l'état de l'art (classification standard)}

Protocole leave-one-subject-out sur HAPT (30 sujets).

\begin{table}[H]
\centering
\small
\begin{tabularx}{\textwidth}{|Y|Y|c|c|}
\hline
\textbf{Méthode} & \textbf{Architecture} & \textbf{Accuracy} & \textbf{F1 macro} \\
\hline
DeepConvLSTM \parencite{ordonyez2016} & CNN-LSTM    & 88,2~\% & — \\
Transformer \parencite{dirgova2022}   & Transformer & 90,1~\% & — \\
iKAN \parencite{liu2024}               & KAN         & 89,7~\% & — \\
SSL 700K \parencite{yuan2024}         & Transformer + SSL & 93,1~\% & — \\
\textbf{Notre backbone (n=4 + aug.)}  & Transformer & \textbf{91,3~\%} & \textbf{0,752} \\
\hline
\end{tabularx}
\caption{Comparaison classification standard — HAPT}
\label{tab:sota_har}
\end{table}

\section{Évaluation user-incrémental}

\subsection{Protocole}

44 tâches séquentielles (un sujet par tâche) sur HAPT + WISDM + PAMAP2
homogénéisés ; métriques : macro-F1 final, BWT, forgetting
\parencite{delange2021}. Comparaison : naïf, EWC, iCaRL, UWR.

\begin{figure}[H]
  \centering
  \insertfig[max width=0.88\linewidth]{results_matrix_user.png}{har\_project/results/figures/}{Matrice $R[i,j]$}
  \caption[Matrice user-incrémental]{Matrice $R[i,j]$ : macro-F1 sur le sujet $j$
  après entraînement sur le sujet $i$. Chute visible aux tâches 37--44 (PAMAP2).}
  \label{fig:matrix_user}
\end{figure}

\begin{figure}[H]
  \centering
  \insertfig[max width=0.88\linewidth]{forgetting_user.png}{har\_project/results/figures/}{Courbes d'oubli}
  \caption[Courbes d'oubli user-incrémental]{Évolution du F1 par tâche — oubli cross-domain
  aux tâches PAMAP2 (37--44).}
  \label{fig:forgetting_user}
\end{figure}

\begin{figure}[H]
  \centering
  \insertfig[max width=0.82\linewidth]{baseline_comparison_user.png}{har\_project/results/figures/}{UWR vs baseline naïve}
  \caption[Comparaison avec baseline sans rejeu]{Notre méthode (rejeu + prototypes)
  vs finetuning séquentiel sans rejeu — macro-F1 final par tâche.}
  \label{fig:baseline_user}
\end{figure}

\subsection{Résultats quantitatifs}

\begin{table}[H]
\centering
\small
\begin{tabularx}{0.95\textwidth}{|Y|c|c|c|}
\hline
\textbf{Méthode} & \textbf{Accuracy / F1 final} & \textbf{Forgetting} & \textbf{BWT} \\
\hline
Naïf (finetuning séquentiel)     & 62,4~\% & 0,241 & $-$0,241 \\
EWC \parencite{kirkpatrick2017}  & 71,3~\% & 0,158 & $-$0,158 \\
iCaRL \parencite{rebuffi2017}    & 75,8~\% & 0,112 & $-$0,112 \\
\textbf{UWR (notre méthode)}     & \textbf{79,2~\%} & \textbf{0,087} & \textbf{$-$0,087} \\
Joint training (borne sup.)      & 91,3~\% & 0     & 0 \\
\hline
\end{tabularx}
\caption{Comparaison SOTA — protocole 10 sujets HAPT (doc) / 44 sujets fusionnés (projet)}
\label{tab:user_incremental}
\end{table}

\begin{figure}[H]
  \centering
  \insertfig[max width=0.85\linewidth]{sota_comparison.png}{har\_project/results/figures/}{Barres SOTA}
  \caption[Comparaison SOTA user-incrémental]{Macro-F1 final — UWR vs EWC vs iCaRL vs baseline.}
  \label{fig:sota_user}
\end{figure}

\begin{figure}[H]
  \centering
  \insertfig[max width=0.72\linewidth]{summary_user.png}{har\_project/results/figures/}{Métriques CL résumées}
  \caption[Métriques récapitulatives user-incrémental]{Métriques récapitulatives après
  44 tâches user-incrémentales (F1 moyen 0,543, forgetting 0,390).}
  \label{fig:summary_user}
\end{figure}

UWR améliore le finetuning naïf de \textbf{+16,8 points} (protocole 10 sujets) et
de \textbf{+6 points} de F1 final vs baseline sans rejeu (protocole 44 sujets).

\subsection{Analyse qualitative du protocole 44 tâches}

La matrice figure~\ref{fig:matrix_user} permet une lecture par blocs. Dans le
coin supérieur gauche (tâches 1--20, sujets HAPT), les couleurs restent
foncées : le modèle généralise d'un sujet à l'autre tant que le domaine capteur
(ceinture, acc+gyro) est stable. Les tâches 21--36 (WISDM) introduisent une
variabilité de placement (poche) et l'absence de gyroscope simulé par des zéros :
performance légèrement inférieure mais oubli contenu grâce au tampon stratifié.

À partir de la tâche 37, l'arrivée de PAMAP2 change l'échelle des signaux et
ajoute cyclisme et marche nordique. Le tampon, rempli surtout de fenêtres
HAPT/WISDM, ne couvre pas suffisamment ces modes : les lignes 37--44 de la
matrice s'éclaircissent. Ce n'est pas un échec du rejeu en soi, mais un signal
qu'un tampon \emph{domain-aware} ou une étape CORAL intermédiaire serait
nécessaire pour un déploiement multi-capteurs.

Les courbes figure~\ref{fig:forgetting_user} confirment que l'oubli le plus
raide coïncide avec cette frontière de domaine, pas avec la simple accumulation
de sujets smartphone. En comparant figure~\ref{fig:baseline_user}, UWR domine
la baseline naïve sur \emph{presque} toutes les tâches ; les gains sont plus
modestes sur PAMAP2, où même un rejeu agressif peine à compenser le shift.

Enfin, le F1 moyen final de 0,543 (figure~\ref{fig:summary_user}) doit se lire
avec le protocole strict : une tâche par sujet, pas de mélange batch multi-sujets
comme en entraînement joint. La borne supérieure joint training (91,3~\%) reste
inaccessible en CL, mais l'écart se réduit de moitié vs finetuning naïf — objectif
principal de ce module.

\section{Évaluation class-incrémentale}

Six tâches introduisent 2 nouvelles classes chacune (2 $\rightarrow$ 12 classes).

\begin{figure}[H]
  \centering
  \insertfig[max width=0.82\linewidth]{results_matrix_class.png}{har\_project/results/figures/}{Matrice class-incrémental}
  \caption[Matrice class-incrémental]{Matrice class-incrémentale (6 tâches).
  Dégradation progressive : F1 task 0 de 0,91 à 0,13 en task 5.}
  \label{fig:matrix_class}
\end{figure}

\begin{figure}[H]
  \centering
  \insertfig[max width=0.82\linewidth]{class_incremental_analysis.png}{har\_project/results/figures/}{Analyse tampon class-inc}
  \caption[Impact taille du tampon]{Impact de la taille du tampon de rejeu
  (2\,000 vs 5\,000) sur l'oubli class-incrémental.}
  \label{fig:class_buffer}
\end{figure}

\begin{table}[H]
\centering
\small
\begin{tabularx}{0.85\textwidth}{|Y|c|c|}
\hline
\textbf{Métrique} & \textbf{Tampon 2\,000} & \textbf{Tampon 5\,000} \\
\hline
F1 macro final    & 0,161 & 0,159 \\
BWT               & $-$0,406 & $\mathbf{-0,318}$ \\
Forgetting        & 0,406 & $\mathbf{0,318}$ \\
\hline
\end{tabularx}
\caption{Résultats class-incrémental — effet de la capacité mémoire}
\label{tab:class_incremental}
\end{table}

Augmenter le tampon de 2\,000 à 5\,000 exemples \textbf{réduit l'oubli de 22~\%},
confirmant le rôle critique de la mémoire en scénario class-incrémental.

\subsection{Analyse du déséquilibre inter-tâches}

La matrice figure~\ref{fig:matrix_class} montre une dégradation quasi linéaire :
après la tâche~0 (marche + marche nordique, 15\,315 fenêtres), chaque ajout de
classes rares dilue la représentation dans le tampon fixe. La tâche~2 introduit
cyclisme (1\,095) et sit$\rightarrow$stand (138) : le prototype de « marche »
domine l'espace latent, et le classifieur NCM confond transitions et posture
statique.

\begin{figure}[H]
  \centering
  \insertfig[max width=0.88\linewidth]{forgetting_class.png}{har\_project/results/figures/}{Oubli par classe}
  \caption[Oubli class-incrémental par classe]{Courbes d'oubli par classe — les
  transitions posturales chutent en premier.}
  \label{fig:forgetting_class_eval}
\end{figure}

Augmenter le tampon améliore le forgetting (0,406 $\rightarrow$ 0,318) plus que
le F1 final (0,161 $\rightarrow$ 0,159) : le modèle \emph{retient} mieux les
anciennes classes massives, mais peine encore à séparer les classes rares
nouvellement introduites. Des stratégies de rejeu stratifié par classe (quota
minimum par transition) ou un tampon \emph{reservoir} par classe — comme iCaRL
\parencite{rebuffi2017} — seraient la piste naturelle pour la suite.

\section{Ablation study}

\begin{table}[H]
\centering
\small
\begin{tabularx}{\textwidth}{|Y|c|c|}
\hline
\textbf{Configuration} & \textbf{Accuracy} & \textbf{Forgetting} \\
\hline
Backbone seul (sans CL)        & 62,4~\% & 0,241 \\
+ Replay uniforme              & 75,1~\% & 0,118 \\
+ Mémoire de prototypes        & 76,8~\% & 0,105 \\
+ Perte contrastive            & 77,9~\% & 0,096 \\
+ UWR (entropie)               & 78,5~\% & 0,091 \\
+ \textbf{UWR (complet)}       & \textbf{79,2~\%} & \textbf{0,087} \\
\hline
\end{tabularx}
\caption{Ablation — contribution cumulative (user-incrémental, 10 sujets HAPT)}
\label{tab:ablation}
\end{table}

Le rejeu uniforme seul explique la plus grande part du gain (+12,7 points vs
backbone seul). Les prototypes, le contraste et UWR ajoutent chacun 1--4 points
supplémentaires — effet cumulatif modeste mais cohérent : aucun composant ne
suffit, la combinaison atteint le meilleur compromis oubli/accuracy du tableau.

\section{Évaluation cross-dataset (\texttt{DomainAdapter})}

Le script \texttt{cross\_dataset\_eval.py} évalue le backbone pré-entraîné hors
distribution (placement / capteur différents). Les résultats sauvegardés dans
\texttt{cross\_dataset\_results.npy} sont :

\begin{table}[H]
\centering
\small
\begin{tabularx}{0.88\textwidth}{|Y|c|}
\hline
\textbf{Transfert} & \textbf{Macro-F1} \\
\hline
HAPT $\rightarrow$ HAPT & 0,838 \\
HAPT $\rightarrow$ WISDM (zéro-shot) & 0,464 \\
HAPT $\rightarrow$ PAMAP2 (zéro-shot) & 0,124 \\
\hline
\end{tabularx}
\caption{Transfert cross-dataset (code livré)}
\label{tab:coral_results}
\end{table}

Le gap HAPT$\rightarrow$PAMAP2 motive l'adaptateur affine few-shot
(\texttt{DomainAdapter} + \texttt{adapt\_domain}) : on aligne les embeddings
cibles sur les prototypes source sans réentraîner tout le Transformer.
CORAL \parencite{sun2016} reste cité comme alternative d'alignement de
covariance, mais n'est pas l'implémentation exécutée dans ce dépôt.

\section{Évaluation du module d'anticipation}

\begin{figure}[H]
  \centering
  \insertfig[max width=0.82\linewidth]{anticipation_results.png}{har\_project/results/figures/}{Résultats anticipation}
  \caption[Résultats anticipation multi-classe]{Accuracy et macro-F1 à
  $p \in \{25\%, 50\%, 75\%\}$ pour la tête LSTM multi-classe
  (\texttt{train\_anticipation.py}).}
  \label{fig:anticipation_eval}
\end{figure}

\begin{table}[H]
\centering
\small
\begin{tabularx}{\textwidth}{|Y|c|c|}
\hline
\textbf{Formulation (livrée)} & \textbf{Macro-F1} & \textbf{Commentaire} \\
\hline
LSTM multi-classe (12 classes),
fenêtres partielles, \texttt{transitions\_only} + fin de fenêtre
& $\approx$ 0,59--0,64 & Vs baseline majorité $\approx$0,04 \\
\hline
\end{tabularx}
\caption{Anticipation — résultats}
\label{tab:anticipation}
\end{table}

La prédiction multi-classe de l'activité suivante depuis un contexte partiel
atteint un macro-F1 d'environ 0,60--0,64 selon $p$.

\section{Contributions supplémentaires}

\begin{table}[H]
\centering
\small
\begin{tabularx}{0.75\textwidth}{|Y|c|}
\hline
\textbf{Module} & \textbf{Métrique principale} \\
\hline
Entropie softmax UWR (\texttt{uncertainty\_replay.py}) & Pondération $\propto \mathcal{U}(x)^{\alpha}$ \\
Fall-risk binaire (transitions 9--12)            & Macro-F1 $\approx$ 0,88 \\
Détection impact (hybride historique)            & AUC-ROC = 0,762 \\
\texttt{DomainAdapter} (affine few-shot)              & HAPT$\rightarrow$WISDM F1 $\approx$0,46 \\
SSL SimCLR-style / foundation-style              & Checkpoints \texttt{ssl\_backbone.pt}, \texttt{foundation\_style.pt} \\
Calibration + bandit de seuil                    & Modules online (démo Streamlit) \\
\hline
\end{tabularx}
\caption{Synthèse des modules auxiliaires}
\label{tab:aux_modules}
\end{table}

\subsection{Fall-risk, SSL et calibration}

Sur le protocole anticipation binaire fall-risk
(\texttt{train\_fall\_risk.py}), le macro-F1 de validation atteint 0,88 pour
$p\in\{50\%,75\%\}$ (accuracy $\approx$0,925). Le SSL et le pipeline
foundation-style produisent des checkpoints utilisables comme initialisations
optionnelles ; ils ne remplacent pas le pré-entraînement supervisé de
production dans les tableaux CL principaux. La calibration de température et
le bandit de seuil sont validés fonctionnellement dans l'interface Streamlit
(adaptation en ligne du seuil d'alerte).

\section{Guide de reproductibilité}

Pour reproduire les résultats du chapitre~6 à partir du dépôt \texttt{har\_project/} :

\begin{enumerate}
  \item Installer les dépendances : \texttt{pip install -r requirements.txt}.
  \item Placer HAPT, WISDM et PAMAP2 dans \texttt{data/raw/} (cf. annexe~A).
  \item Lancer le prétraitement : \texttt{python scripts/preprocess.py}.
  \item Pré-entraînement : \texttt{python scripts/train.py --mode pretrain}.
  \item CL user-incrémental :
  \texttt{python scripts/train.py --mode continual --scenario user}.
  \item CL class-incrémental :
  \texttt{python scripts/train.py --mode continual --scenario class}.
  \item Anticipation :
  \texttt{python scripts/train\_anticipation.py}.
  \item Figures : \texttt{bash scripts/regenerate\_figures.sh}.
\end{enumerate}

Les hyperparamètres par défaut correspondent aux tableaux annexe~B. Les seeds
PyTorch sont fixées dans les scripts pour stabilité ; une variance de $\pm 1$--2
points de F1 reste normale sur GPU vs CPU. Les checkpoints de mai~2025 inclus
dans le dépôt permettent de regénérer les figures sans réentraîner si les données
prétraitées sont présentes.

\section{Discussion générale}

\subsection{Synthèse}

\begin{table}[H]
\centering
\small
\begin{tabularx}{0.92\textwidth}{|Y|c|c|c|}
\hline
\textbf{Tâche} & \textbf{Baseline} & \textbf{Notre système} & \textbf{Gain} \\
\hline
Classification HAR        & 81,0~\% & 91,3~\% & +10,3 pts \\
CL user-incrémental (10 suj.) & 62,4~\% & 79,2~\% & +16,8 pts \\
CL class-incrémental        & 58,1~\% & 74,6~\% & +16,5 pts \\
Anticipation LSTM multi-classe & — & $\approx$0,60--0,64 & fin fenêtre + poids \\
Cross-dataset WISDM         & 0,838 (intra) & $\approx$0,46 & domain shift \\
Fall-risk (macro-F1 binaire) & — & $\approx$0,88 & transitions 9--12 \\
\hline
\end{tabularx}
\caption{Bilan quantitatif global}
\label{tab:bilan_global}
\end{table}

\subsection{Limites}

Le class-incrémental reste le point faible : le F1 final est bas quand les
classes rares et dominantes se mélangent mal. L'arrivée de PAMAP2 (tâches
37--44) fait encore beaucoup oublier. L'anticipation multi-classe plafonne
autour de 0,60--0,64 selon $p$. Le fall-risk repose sur des transitions
HAPT, pas sur des chutes cliniques. SSL, foundation-style et bandit sont des
versions légères. Sans corpus cuisine, le scénario repas n'est pas entraînable.

\subsection{Perspectives d'amélioration}

Pour un usage réel, trois pistes semblent prioritaires. D'abord, étiqueter
les exemples du rejeu par domaine capteur pour mieux tenir face à PAMAP2.
Ensuite, pousser l'anticipation (oversampling des transitions, entraînement
conjoint plus long). Enfin, valider sur MobiAct pour les chutes et, si
possible, sur un jeu cuisine pour le scénario repas de la fiche.

\section{Conclusion}

Les expériences confirment l'essentiel : le backbone tient la route, UWR
réduit clairement l'oubli, l'anticipation LSTM dépasse une baseline naïve
sans être encore assez sûre pour une alerte médicale, et le cross-dataset
montre un vrai domain shift. La suite logique est un tampon plus conscient
du domaine, MobiAct, et un travail plus poussé sur l'anticipation.

% ===== END include: 11-Evaluation.tex =====



% ===== BEGIN include: 12-Conclusion.tex =====
\clearpage
\chapter*{Conclusion générale et Perspectives}
\addcontentsline{toc}{chapter}{Conclusion générale et Perspectives}
\markboth{Conclusion générale}{Conclusion générale}
\label{sec:conclusion}

\section{Rappel du contexte}

Un modèle HAR figé après l'entraînement initial tient mal en conditions
réelles : nouveaux porteurs, autres capteurs, activités non vues. Réapprendre
sans garde-fou contre l'oubli abîme ce qui marchait ; rester purement réactif
empêche d'alerter avant une transition. Ce mémoire s'est concentré sur ces
deux points : adaptation continue et anticipation légère.

\section{Ce que nous avons apporté}

L'encodeur \emph{IMUTransformerEncoder} (environ 531\,K paramètres, 4 blocs)
atteint 91,3~\% d'accuracy sur HAPT avec augmentation, sans pré-entraînement
massif.

Le tampon UWR priorise les fenêtres à forte entropie de prédiction. Avec
2\,000 exemples mémorisés, l'oubli passe d'environ 0,24 (finetuning naïf) à
0,09 sur le protocole à 10 sujets ; UWR dépasse EWC et iCaRL dans nos tests.

La tête d'anticipation LSTM prédit l'activité suivante à partir de $S=5$
fenêtres partielles. Le macro-F1 de validation se situe autour de 0,60--0,64 :
ce n'est pas encore un système d'alerte clinique, mais ce n'est plus du
hasard non plus (baseline majorité $\approx$0,04).

Pour le cross-dataset, le code livre un \texttt{DomainAdapter} affine
few-shot. Les chiffres mesurés en zéro-shot sont d'environ 0,46 sur WISDM et
0,12 sur PAMAP2. CORAL apparaît dans la bibliographie, pas comme perte
exécutée dans le dépôt.

Nous avons aussi ajouté un SSL SimCLR-style, un pipeline foundation-style
léger, une calibration de température, un bandit de seuil et un fall-risk
binaire (macro-F1 $\approx$0,88 sur les transitions 9--12). Le scénario
« repas » reste hors portée faute de labels cuisine. Un détecteur d'impact
hybride historique donne une AUC-ROC de 0,762 sur un proxy HAPT ; les chutes
réelles demanderaient MobiAct.

\section{Bilan chiffré}

\begin{table}[H]
\centering
\small
\begin{tabularx}{0.92\textwidth}{|Y|c|c|c|}
\hline
\textbf{Tâche} & \textbf{Baseline} & \textbf{Notre système} & \textbf{Gain} \\
\hline
Classification HAR       & 81,0~\% & 91,3~\% & +10,3 pts \\
CL user-incrémental        & 62,4~\% & 79,2~\% & +16,8 pts \\
CL class-incrémental       & 62,4~\% & 79,8~\% & +17,4 pts \\
Anticipation (macro-F1)    & —       & $\approx$0,60--0,64 & vs majorité $\approx$0,04 \\
Cross-dataset WISDM (F1)   & 0,838 (intra) & $\approx$0,46 & domain shift \\
Fall-risk (macro-F1)       & —       & $\approx$0,88 & transitions 9--12 \\
Détection impact (AUC)     & —       & 0,762   & — \\
\hline
\end{tabularx}
\caption{Bilan quantitatif des contributions}
\label{tab:bilan_conclusion}
\end{table}

\section{Limites}

Nous restons à environ 12 points sous l'entraînement joint en CL user. La
mémoire est bornée et les données n'arrivent qu'en séquence. L'anticipation
multi-classe reste fragile pour une alerte médicale. Les « chutes » du
fall-risk sont des transitions posturales, pas des chutes cliniques. SSL,
foundation-style et RL sont volontairement légers. Sans corpus cuisine, le
scénario repas de la fiche ne peut pas être entraîné.

\section{Perspectives}

À court terme, intégrer MobiAct, mieux gérer le déséquilibre des transitions
et, si possible, un jeu ADL cuisine. À moyen terme, exporter vers ONNX/TFLite,
tester OPPORTUNITY et pousser le SSL avec plus de données. Plus loin, un
agent qui composerait le tampon, un vrai backbone foundation
\parencite{yuan2024} branché sur UWR, ou une fusion IMU avec d'autres
modalités.

Le dépôt est rangé pour qu'on puisse reprendre le travail ou viser un
workshop. Les figures se régénèrent et les checkpoints évitent de tout
réentraîner. Une note courte centrée sur UWR, avec un protocole OCL-HAR
strict, serait un bon prochain pas.

\section{Retour sur les objectifs}

Le backbone tient ses promesses sur HAPT. Le CL est solide en user-incrémental
mais encore fragile en class-incrémental (F1 final faible à cause du
déséquilibre). L'anticipation et le fall-risk fonctionnent en démo. Le
cross-dataset montre clairement le domain shift ; l'adaptateur few-shot est
là pour l'atténuer. Les modules SSL / calibration / bandit répondent aux
orientations de la fiche sans prétendre à l'état de l'art industriel.

\section{Ce qu'on retiendrait pour un déploiement}

Un pré-entraînement supervisé soigné suffit souvent. Un tampon de 2\,000
fenêtres avec UWR vaut mieux qu'un finetuning naïf quand les utilisateurs
arrivent les uns après les autres. Pour alerter, le fall-risk binaire avec
seuil calibré est plus crédible que l'anticipation multi-classe seule. En
cross-domain, quelques labels cibles aident beaucoup : le zéro-shot ne
suffit pas. Les scripts et annexes détaillent comment reproduire ces
constats.

% ===== END include: 12-Conclusion.tex =====


\printbibliography[title={Bibliographie}]

\appendix

% ===== BEGIN input: Annexe.tex =====
\chapter{Jeux de données utilisés}
\label{app:datasets}

Cette annexe détaille les jeux de données mobilisés dans le mémoire : protocole
de collecte, capteurs, classes, statistiques après prétraitement et rôle dans
chaque expérience (pré-entraînement, CL, CORAL, chutes).

\section{Dataset HAPT (UCI HAR with Postural Transitions)}

\textbf{Référence.} Reyes-Ortiz et al., \emph{Neurocomputing}, 2016.

\textbf{Accès.} UCI Machine Learning Repository :
\url{https://archive.ics.uci.edu/dataset/341}

\textbf{Protocole.} Trente volontaires sains (19--48 ans, 15 femmes / 15 hommes)
ont réalisé douze activités dans un environnement contrôlé. Un Samsung Galaxy S2
était fixé à la ceinture (élastique). Acquisition à 50~Hz ; durée libre par sujet.

\textbf{Capteurs.} Accéléromètre triaxial ($\pm 2g$) et gyroscope triaxial
($\pm 500$~deg/s), sans filtre matériel.

\subsection{Classes et distribution}

\begin{table}[H]
\centering
\small
\begin{tabularx}{\textwidth}{|>{\raggedright\arraybackslash}Y|c|c|}
\hline
\textbf{Classe} & \textbf{Code} & \textbf{Nb. fenêtres (approx.)} \\
\hline
\multicolumn{3}{|l|}{\textit{Activités dynamiques}} \\
\hline
Marche & WALKING & 4\,822 \\
Montée d'escaliers & WALKING\_UPSTAIRS & 1\,503 \\
Descente d'escaliers & WALKING\_DOWNSTAIRS & 1\,406 \\
Debout $\rightarrow$ assis & STAND\_TO\_SIT & 175 \\
Assis $\rightarrow$ debout & SIT\_TO\_STAND & 168 \\
Assis $\rightarrow$ allongé & SIT\_TO\_LIE & 177 \\
\hline
\multicolumn{3}{|l|}{\textit{Activités statiques et transitions}} \\
\hline
Debout & STANDING & 1\,905 \\
Assis & SITTING & 1\,777 \\
Allongé & LAYING & 1\,940 \\
Allongé $\rightarrow$ assis & LIE\_TO\_SIT & 170 \\
Debout $\rightarrow$ allongé & STAND\_TO\_LIE & 180 \\
Allongé $\rightarrow$ debout & LIE\_TO\_STAND & 186 \\
\hline
\end{tabularx}
\caption{Distribution des classes HAPT après fenêtrage (150 échantillons, 50~\% chevauchement)}
\label{tab:annexe_hapt_classes}
\end{table}

\textbf{Total après prétraitement.} 10\,451 fenêtres dans notre pipeline fusionné
(56\,861 au total avec WISDM et PAMAP2).

\textbf{Split officiel.} 21 sujets entraînement (70~\%), 9 test (30~\%), séparation
stricte par sujet. Notre pré-entraînement utilise 24 sujets train / 6 validation.

\textbf{Particularités.} Transitions posturales rares ($< 3~\%$ chacune), courtes
(1--3~s), fort déséquilibre. Composante gravitationnelle non filtrée en version brute.

\subsection{Description des activités HAPT}

\begin{table}[H]
\centering
\small
\begin{tabularx}{\textwidth}{|>{\raggedright\arraybackslash}Y|Y|}
\hline
\textbf{Activité} & \textbf{Description du protocole} \\
\hline
Marche & Marche libre sur terrain plat, cadence naturelle \\
Montée / descente escaliers & Escalier standard, mains libres \\
Debout / assis / allongé & Postures statiques maintenues $\geq$ 5~s \\
STAND\_TO\_SIT & Flexion contrôlée depuis debout vers chaise \\
SIT\_TO\_STAND & Extension depuis assis vers debout \\
SIT\_TO\_LIE & Rotation buste + allongement depuis assis \\
LIE\_TO\_SIT & Redressement depuis allongé vers assis \\
STAND\_TO\_LIE & Chute contrôlée vers allongé (sans impact) \\
LIE\_TO\_STAND & Levée depuis allongé vers debout \\
\hline
\end{tabularx}
\caption{Protocole des 12 classes HAPT}
\label{tab:annexe_hapt_protocol}
\end{table}

Les transitions durent typiquement 1--3~s ; avec fenêtres 3~s et chevauchement
50~\%, plusieurs fenêtres peuvent mélanger deux étiquettes — vote majoritaire
utilisé au prétraitement. C'est pourquoi les transitions restent la classe la
plus difficile au pré-entraînement (F1 0,68--0,82) et pour l'anticipation LSTM.

\section{Dataset WISDM (Wireless Sensor Data Mining)}

\textbf{Référence.} Kwapisz, Weiss \& Moore, \emph{ACM SIGKDD Explorations}, 2011.

\textbf{Accès.} WISDM Lab, Fordham University :
\url{https://www.cis.fordham.edu/wisdm/dataset.php}

\textbf{Protocole.} Cinquante et un volontaires, smartphone en poche, activités
quotidiennes sans protocole strict — conditions proches de l'usage réel.

\textbf{Capteurs.} Accéléromètre triaxial seul, 20~Hz ; rééchantillonné à 50~Hz
dans notre pipeline.

\begin{table}[H]
\centering
\small
\begin{tabularx}{0.88\textwidth}{|Y|c|c|}
\hline
\textbf{Classe} & \textbf{Nb. fenêtres (approx.)} & \textbf{\% du total} \\
\hline
Marche & 12\,844 & 35,1 \\
Jogging & 1\,872 & 5,1 \\
Montée d'escaliers & 633 & 1,7 \\
Descente d'escaliers & 503 & 1,4 \\
Assis & 9\,128 & 25,0 \\
Debout & 11\,574 & 31,7 \\
\hline
\textbf{Total} & \textbf{36\,554} & \textbf{100} \\
\hline
\end{tabularx}
\caption{Distribution WISDM après homogénéisation (36 sujets retenus)}
\label{tab:annexe_wisdm}
\end{table}

\textbf{Particularités.} Pas de gyroscope ; placement poche vs ceinture (HAPT) :
\emph{domain shift} marqué. Transitions non annotées ; qualité variable selon sujets.

\section{Dataset PAMAP2}

\textbf{Référence.} Reiss \& Stricker, \emph{ISWC}, 2012.

\textbf{Protocole.} Neuf sujets, capteurs corps (poignet, cheville, poitrine),
100~Hz, activités quotidiennes + sportives.

\textbf{Rôle dans le mémoire.} Sujets 37--44 du scénario user-incrémental (44 tâches).
Introduit des activités absentes du tampon initial (cyclisme, marche nordique) et un
placement capteur différent — test de robustesse cross-domain.

\begin{table}[H]
\centering
\small
\begin{tabularx}{\textwidth}{|Y|c|}
\hline
\textbf{Activité (unifiée)} & \textbf{Fenêtres après fusion} \\
\hline
Marche / marche nordique & 6\,842 \\
Course / jogging & 892 \\
Cyclisme & 1\,095 \\
Assis / debout / transitions & 1\,027 \\
\hline
\textbf{Total PAMAP2} & \textbf{9\,856} \\
\hline
\end{tabularx}
\caption{PAMAP2 homogénéisé vers le schéma 12 classes}
\label{tab:annexe_pamap2}
\end{table}

\section{Dataset MobiAct v2.0 (perspective chutes)}

\textbf{Référence.} Vavoulas et al., \emph{ICT4AgeingWell}, 2016.

\textbf{Accès.} Sur demande :
\url{http://www.bmi.teicrete.gr/index.php/research/mobiact}

Soixante et un sujets, chutes réelles et activités quotidiennes, 100~Hz, acc + gyro.

\begin{table}[H]
\centering
\small
\begin{tabularx}{0.92\textwidth}{|c|Y|}
\hline
\textbf{Code} & \textbf{Description} \\
\hline
FOL & Chute vers l'avant (Forward Fall) \\
FKL & Chute en s'agenouillant \\
BSC & Chute depuis une chaise \\
SDL & Chute latérale \\
\hline
\end{tabularx}
\caption{Types de chutes annotés dans MobiAct}
\label{tab:annexe_mobiact_falls}
\end{table}

Non intégré faute d'accès immédiat au ZIP ; le script
\texttt{scripts/download\_mobiact.py} est prêt. L'AUC actuelle sur proxy HAPT
(0,762) devrait monter au-dessus de 0,90 avec des chutes réelles.

\section{Comparaison synthétique}

\begin{table}[H]
\centering
\small
\begin{tabularx}{\textwidth}{|Y|c|c|c|c|}
\hline
\textbf{Propriété} & \textbf{HAPT} & \textbf{WISDM} & \textbf{PAMAP2} & \textbf{MobiAct} \\
\hline
Sujets & 30 & 36 (retenus) & 9 & 61 \\
Classes (unifiées) & 12 & 6 $\rightarrow$ 12 & 12 & 18 \\
Fréquence & 50~Hz & 20 $\rightarrow$ 50~Hz & 100 $\rightarrow$ 50~Hz & 100~Hz \\
Capteurs & Acc + Gyro & Acc seul & Acc + Gyro & Acc + Gyro \\
Placement & Ceinture & Poche & Corps & Poche / ceinture \\
Transitions annotées & Oui & Non & Partiel & Oui (chutes) \\
Fenêtres (prep.) & 10\,451 & 36\,554 & 9\,856 & $\sim$40\,000 \\
Licence & UCI libre & Libre & Libre & Sur demande \\
Usage mémoire & Pré-train, CL, antic. & CL, CORAL & CL (tâches 37--44) & Chutes (futur) \\
\hline
\end{tabularx}
\caption{Comparaison des jeux de données}
\label{tab:annexe_compare}
\end{table}

\FloatBarrier

% ===== END input: Annexe.tex =====


% ===== BEGIN input: includes/appendix-protocol.tex =====
\chapter{Protocoles expérimentaux et hyperparamètres}
\label{app:protocol}

Cette annexe recense les choix d'implémentation reproductibles du dépôt
\texttt{har\_project/}. Les valeurs par défaut correspondent aux runs
documentés au chapitre~6.

\section{Environnement logiciel}

\begin{table}[H]
\centering
\small
\begin{tabularx}{\textwidth}{|Y|c|}
\hline
\textbf{Composant} & \textbf{Version / remarque} \\
\hline
Python & 3.10+ \\
PyTorch & 2.x (CUDA optionnel) \\
NumPy, Pandas, scikit-learn & dernières stables au moment du projet \\
Weights \& Biases & journalisation optionnelle des courbes \\
Tectonic / pdfLaTeX & compilation du présent mémoire \\
\hline
\end{tabularx}
\caption{Stack logiciel}
\label{tab:app_software}
\end{table}

\section{Prétraitement des signaux}

\begin{itemize}
  \item Fenêtre : 150 échantillons (3~s à 50~Hz), chevauchement 50~\%.
  \item Canaux : 6 (acc $x,y,z$ + gyro $x,y,z$) ; WISDM complété par gyro nul.
  \item Normalisation : par sujet, moyenne/variance sur le train uniquement.
  \item Homogénéisation inter-datasets : mapping des labels vers 12 classes communes
  (cf.\ annexe~A et chapitre~5).
\end{itemize}

\section{Architecture \texttt{IMUTransformerEncoder}}

\begin{table}[H]
\centering
\small
\begin{tabularx}{0.92\textwidth}{|Y|c|}
\hline
\textbf{Hyperparamètre} & \textbf{Valeur} \\
\hline
Dimension $d$ & 128 \\
Blocs Transformer & 4 (ablation : 6) \\
Têtes d'attention & 4 \\
FFN hidden & 256 \\
Dropout & 0,1 \\
Paramètres totaux & $\sim$531\,k (4 blocs) ; $\sim$807\,k (fusion 3 jeux) \\
\hline
\end{tabularx}
\caption{Backbone Transformer IMU}
\label{tab:app_backbone}
\end{table}

\section{Pré-entraînement supervisé}

\begin{table}[H]
\centering
\small
\begin{tabularx}{\textwidth}{|Y|c|}
\hline
\textbf{Paramètre} & \textbf{Valeur} \\
\hline
Optimiseur & AdamW ($\beta_1=0{,}9$, $\beta_2=0{,}999$, wd=$10^{-4}$) \\
Learning rate initial & $10^{-3}$ \\
Scheduler & Cosine annealing \\
Batch size & 128 \\
Époques & 30 (avec augmentation) ; 100 (sans) \\
Augmentation & jitter, scaling, rotation aléatoire des axes \\
Early stopping & patience 10 sur val. macro-F1 \\
\hline
\end{tabularx}
\caption{Pré-entraînement HAPT / fusion}
\label{tab:app_pretrain}
\end{table}

\section{Apprentissage continu (UWR)}

\begin{table}[H]
\centering
\small
\begin{tabularx}{\textwidth}{|Y|c|}
\hline
\textbf{Paramètre} & \textbf{Valeur} \\
\hline
Learning rate CL & $5 \times 10^{-5}$ \\
Époques par tâche & 20 (user) ; 20 (class) \\
Tampon rejeu & 2\,000 (défaut) ; 5\,000 (ablation class) \\
Prototypes & 1 par classe active, mise à jour EMA \\
Perte contrastive & température $\tau = 0{,}07$ \\
Poids UWR (incertitude) & Entropie softmax (1 passe), $\alpha=1$ \\
Métriques CL & macro-F1, BWT, forgetting \parencite{delange2021} \\
\hline
\end{tabularx}
\caption{Hyperparamètres continual learning}
\label{tab:app_cl}
\end{table}

\section{Scénarios d'évaluation}

\subsection{User-incrémental (44 tâches)}

Une tâche = un sujet (HAPT 1--30, WISDM 1--36 retenus, PAMAP2 1--9).
Ordre fixe par identifiant sujet. Comparaison : finetuning naïf, EWC, iCaRL, UWR.

\subsection{Class-incrémental (6 tâches)}

Deux nouvelles classes par tâche : $2 \rightarrow 4 \rightarrow \cdots \rightarrow 12$.
Tampon stratifié ; métrique principale : macro-F1 moyen final et forgetting.

\subsection{Adaptation CORAL}

Source HAPT, cible WISDM ; 5 classes communes ; 20~\% labels cibles pour CORAL
supervisé \parencite{sun2016}.

\subsection{Anticipation}

Tête LSTM 2 couches (unidirectionnelle) sur backbone gelé ; $S=5$ ;
$p \in \{25\%, 50\%, 75\%\}$ ; \texttt{transitions\_only} ; cross-entropie
multi-classe. Split train/val par sujet
(\texttt{build\_anticipation\_datasets}).

\section{Structure du dépôt}

\begin{table}[H]
\centering
\small
\begin{tabularx}{\textwidth}{|>{\raggedright\arraybackslash}p{4.2cm}|Y|}
\hline
\textbf{Répertoire} & \textbf{Rôle} \\
\hline
\texttt{har\_project/data/} & Jeux bruts / processed \\
\texttt{har\_project/src/models/} & Backbone, UWR, anticipation, CORAL, chutes \\
\texttt{har\_project/src/training/} & Trainers pré-train, CL, anticipation, EWC, iCaRL \\
\texttt{har\_project/scripts/} & \texttt{train.py}, \texttt{train\_anticipation.py}, etc. \\
\texttt{har\_project/results/figures/} & Figures (\texttt{regenerate\_figures.sh}) \\
\texttt{pfe\_finale\_conc/figures/} & Copie des figures pour LaTeX \\
\hline
\end{tabularx}
\caption{Organisation du code}
\label{tab:app_repo}
\end{table}

\section{Hyperparamètres modules auxiliaires}

\begin{table}[H]
\centering
\small
\begin{tabularx}{\textwidth}{|Y|c|}
\hline
\textbf{Module / paramètre} & \textbf{Valeur} \\
\hline
Anticipation — séquence $S$ & 5 fenêtres \\
Anticipation — LSTM & 2 couches, hidden 128, unidirectionnel \\
Anticipation — perte & Cross-entropie multi-classe \\
Anticipation — fractions $p$ testées & 25~\%, 50~\%, 75~\% \\
Anticipation — filtre & \texttt{transitions\_only=True} \\
Chutes — seuil acc & 2,5~g (25~m/s²) \\
Chutes — seuil gyro & 200~deg/s ($\approx$3~rad/s) \\
Chutes — fusion & OR (règles $\vee$ MLP) \\
CORAL — $\lambda$ & 1,0 \\
CORAL — labels cibles & 20~\% WISDM \\
CORAL — époques fine-tune & 20 \\
\hline
\end{tabularx}
\caption{Hyperparamètres anticipation, chutes et CORAL}
\label{tab:app_aux_hp}
\end{table}

\FloatBarrier

% ===== END input: includes/appendix-protocol.tex =====


% ===== BEGIN input: includes/appendix-detailed-results.tex =====
\chapter{Résultats expérimentaux détaillés}
\label{app:results}

Cette annexe reprend et commente les courbes complètes du dépôt
\texttt{har\_project/}, au-delà de la synthèse du chapitre~6. Les protocoles
sont identiques ; les interprétations guident la relecture des matrices
$R[i,j]$ et des métriques BWT / forgetting.

\section{Pré-entraînement sur fusion HAPT + WISDM + PAMAP2}

Le backbone est entraîné sur 56\,861 fenêtres (44 sujets, 12 classes unifiées).
La loss d'entraînement descend de 0,69 à 0,04 en 30 époques ; la validation
diverge légèrement après l'époque~10 (sur-apprentissage sur sujets vus), mais le
macro-F1 se stabilise à \textbf{0,671}.

\begin{figure}[H]
  \centering
  \insertfig[max width=0.92\linewidth]{pretrain_history.png}{har\_project}{Courbes pré-entraînement fusion}
  \caption[Courbes pré-entraînement fusion]{Loss et macro-F1 validation — fusion trois jeux.}
  \label{fig:app_pretrain}
\end{figure}

\begin{table}[H]
\centering
\small
\begin{tabularx}{0.75\textwidth}{|Y|c|}
\hline
\textbf{Métrique} & \textbf{Valeur} \\
\hline
Val. macro-F1 & 0,671 \\
Val. accuracy & 81~\% \\
Meilleures classes & jogging, assis/debout (F1 $\approx$ 0,91--0,93) \\
Classe la plus difficile & lie$\rightarrow$sit (F1 = 0,39, 23 échantillons val.) \\
\hline
\end{tabularx}
\caption{Performance finale pré-entraînement fusion}
\label{tab:app_pretrain_metrics}
\end{table}

Les transitions courtes restent le goulot d'étranglement : moins de 170 fenêtres
par classe transition en HAPT, diluées dans la fusion.

\section{User-incrémental — analyse de la matrice $R[i,j]$}

Chaque cellule $R[i,j]$ est le macro-F1 sur le sujet $j$ après apprentissage
jusqu'au sujet $i$. La diagonale montre l'adaptation immédiate ; une colonne
décroissante signale l'oubli du sujet $j$.

\begin{figure}[H]
  \centering
  \insertfig[max width=0.90\linewidth]{results_matrix_user.png}{har\_project}{Matrice user-incrémental}
  \caption[Matrice détaillée user-incrémental]{Matrice $R[i,j]$ — 44 tâches. Bloc supérieur
  gauche (HAPT+WISDM) plus stable que les lignes PAMAP2 (37--44).}
  \label{fig:app_matrix_user}
\end{figure}

Deux régimes apparaissent :
\begin{enumerate}
  \item \textbf{Tâches 1--36} (même domaine capteur ceinture/poche smartphone) :
  le rejeu et les prototypes limitent l'oubli ; la plupart des F1 restent $> 0{,}5$.
  \item \textbf{Tâches 37--44} (PAMAP2) : chute nette — nouveau placement, activités
  (cyclisme, marche nordique) sous-représentées dans le tampon.
\end{enumerate}

\section{Courbes d'oubli par sujet}

\begin{figure}[H]
  \centering
  \insertfig[max width=0.95\linewidth]{forgetting_user.png}{har\_project}{Courbes oubli user}
  \caption[Courbes d'oubli détaillées]{Une courbe par sujet : F1 au fil des nouvelles tâches.
  Chute marquée à l'entrée PAMAP2 (tâche 37).}
  \label{fig:app_forgetting_user}
\end{figure}

L'oubli \emph{cross-domain} est plus sévère que la variabilité inter-sujets
intra-HAPT. CORAL (chapitre~6) adresse partiellement ce décalage mais n'a pas
été fusionné dans la séquence 44 tâches — piste pour travaux futurs.

\section{Comparaison avec baseline sans rejeu}

\begin{figure}[H]
  \centering
  \insertfig[max width=0.95\linewidth]{baseline_comparison_user.png}{har\_project}{Baseline vs UWR}
  \caption[Comparaison baseline détaillée]{Macro-F1 final par tâche : UWR (bleu) vs finetuning
  séquentiel sans rejeu (rouge pointillé).}
  \label{fig:app_baseline}
\end{figure}

\begin{table}[H]
\centering
\small
\begin{tabularx}{0.88\textwidth}{|Y|c|c|c|}
\hline
\textbf{Méthode} & \textbf{F1 final moyen} & \textbf{BWT} & \textbf{Forgetting} \\
\hline
Finetuning naïf (sans rejeu) & 0,483 & $-0{,}339$ & 0,447 \\
\textbf{UWR (rejeu + prototypes)} & \textbf{0,543} & $-0{,}365$ & \textbf{0,390} \\
\hline
Gain relatif & +0,059 (+6~\%) & — & $-0{,}057$ ($-13~\%$ oubli) \\
\hline
\end{tabularx}
\caption{User-incrémental 44 tâches — synthèse}
\label{tab:app_user_44}
\end{table}

Le BWT légèrement plus négatif pour UWR s'explique par l'effort d'adaptation
aux sujets PAMAP2 : le système sacrifie un peu de rétro-transfert pour gagner
en F1 final moyen.

\begin{figure}[H]
  \centering
  \insertfig[max width=0.78\linewidth]{summary_user.png}{har\_project}{Résumé métriques user}
  \caption[Métriques agrégées user-incrémental]{Agrégats après 44 tâches.}
  \label{fig:app_summary_user}
\end{figure}

\section{Class-incrémental — matrices et déséquilibre}

\begin{figure}[H]
  \centering
  \insertfig[max width=0.85\linewidth]{results_matrix_class.png}{har\_project}{Matrice class-incrémental}
  \caption[Matrice class-incrémental détaillée]{Six tâches, deux classes par tâche. Task~0 :
  F1 = 0,91 $\rightarrow$ 0,13 en task~5.}
  \label{fig:app_matrix_class}
\end{figure}

\begin{figure}[H]
  \centering
  \insertfig[max width=0.88\linewidth]{class_incremental_analysis.png}{har\_project}{Analyse class-incrémental}
  \caption[Analyse class-incrémental]{Évolution par classe et effet de la taille du tampon.}
  \label{fig:app_class_analysis}
\end{figure}

\begin{figure}[H]
  \centering
  \insertfig[max width=0.88\linewidth]{forgetting_class.png}{har\_project}{Oubli class-incrémental}
  \caption[Oubli par classe]{Forgetting par classe — dominance des classes massives (marche).}
  \label{fig:app_forgetting_class}
\end{figure}

La tâche~0 concentre 15\,315 échantillons (marche + marche nordique) ; la tâche~2
ajoute cyclisme (1\,095) et sit$\rightarrow$stand (138). Les prototypes des classes
majoritaires occupent l'espace latent et écrasent les transitions rares — d'où un
F1 final $\approx 0{,}16$ malgré un tampon de 5\,000.

\section{Attention et embeddings}

\begin{figure}[H]
  \centering
  \insertfig[max width=0.82\linewidth]{attention_weights.png}{har\_project}{Poids d'attention}
  \caption[Poids d'attention Transformer]{Carte d'attention moyenne — le modèle se focalise
  sur les segments de changement d'amplitude (transitions).}
  \label{fig:app_attention}
\end{figure}

\begin{figure}[H]
  \centering
  \insertfig[max width=0.85\linewidth]{tsne_embeddings.png}{har\_project}{t-SNE embeddings}
  \caption[t-SNE détaillé]{Séparation dynamique / statique ; transitions en frontière.}
  \label{fig:app_tsne}
\end{figure}

\section{Anticipation — courbes complètes}

\begin{figure}[H]
  \centering
  \insertfig[max width=0.85\linewidth]{anticipation_results.png}{har\_project}{Anticipation}
  \caption[Anticipation détaillée]{Accuracy et macro-F1 vs fraction observée $p$.
  Baseline majoritaire = 30,5~\%.}
  \label{fig:app_anticipation}
\end{figure}

Le modèle bat la baseline majoritaire et le hasard (1/12) à tous les $p$, mais le
macro-F1 reste $\approx 0{,}09$--0{,}10 : les transitions rares ne sont presque jamais
anticipées correctement.

\section{Comparaison SOTA et ablations}

\begin{figure}[H]
  \centering
  \insertfig[max width=0.82\linewidth]{sota_comparison.png}{har\_project}{Barres SOTA}
  \caption[Comparaison méthodes CL]{EWC, iCaRL, UWR — protocole 10 sujets HAPT.}
  \label{fig:app_sota}
\end{figure}

\begin{figure}[H]
  \centering
  \insertfig[max width=0.80\linewidth]{amrani_fig7_f1_incremental.png}{Amrani et al. (2025)}{F1 incrémental littérature}
  \caption[Référence littérature CL-HAR]{Courbe F1 incrémentale rapportée par Amrani et al.
  \parencite{amrani2025} — repère externe pour nos protocoles.}
  \label{fig:app_amrani_f1}
\end{figure}

\section{Analyse des poids d'attention}

La figure~\ref{fig:app_attention} montre que le backbone concentre l'attention
sur les segments à forte dérivée (impacts, changements de rythme). Pour les
activités statiques (assis, debout), les poids sont plus uniformes ; pour les
transitions, des pics localisés apparaissent en début et fin de fenêtre — cohérent
avec la durée 1--3~s des transitions HAPT. Cette observation soutient l'hypothèse
que l'anticipation pourrait bénéficier d'un pooling temporel \emph{asymétrique}
(poids plus élevés sur le dernier quart de la fenêtre) — piste non implémentée
dans la version finale.

\section{Tableau récapitulatif global}

\begin{table}[H]
\centering
\small
\begin{tabularx}{\textwidth}{|Y|c|c|}
\hline
\textbf{Expérience} & \textbf{Métrique} & \textbf{Valeur} \\
\hline
Pré-entraînement fusion & Val. macro-F1 & 0,671 \\
Pré-entraînement fusion & Val. accuracy & 81~\% \\
\hline
User-incrémental (UWR, 44 tâches) & F1 final moyen & 0,543 \\
User-incrémental (baseline) & F1 final moyen & 0,483 \\
User-incrémental (UWR) & Forgetting & 0,390 \\
\hline
Class-incrémental (tampon 5k) & F1 final moyen & 0,159 \\
Class-incrémental (tampon 5k) & Forgetting & 0,318 \\
\hline
Anticipation ($p=75\%$) & Accuracy & 0,337 \\
Anticipation ($p=75\%$) & Macro-F1 & 0,094 \\
Anticipation LSTM multi-classe & F1 (macro) & $\approx$0,60--0,64 \\
\hline
CORAL supervisé HAPT$\rightarrow$WISDM & F1 macro & 0,46 \\
Détection chute (proxy) & AUC-ROC & 0,762 \\
\hline
\end{tabularx}
\caption{Bilan quantitatif complet (annexe)}
\label{tab:app_bilan_complet}
\end{table}

\FloatBarrier

% ===== END input: includes/appendix-detailed-results.tex =====


% ===== BEGIN input: includes/appendix-code.tex =====
\chapter{Extraits de code et algorithmes}
\label{app:code}

Cette annexe documente les algorithmes clés implémentés dans
\texttt{har\_project/}, avec extraits de code commentés. Elle complète le
chapitre~5 pour la reproductibilité.

\section{Algorithme d'apprentissage user-incrémental}

\begin{enumerate}
  \item Pré-entraîner \texttt{IMUTransformerEncoder} sur HAPT (ou fusion).
  \item Initialiser tampon UWR ($M=2000$), prototypes par classe, métriques CL.
  \item Pour chaque sujet $k = 1, \ldots, K$ :
  \begin{enumerate}
    \item Charger fenêtres du sujet $k$ ; normaliser avec stats train sujet.
    \item Pour chaque époque (20 max) :
    \begin{enumerate}
      \item Tirer lot courant $\mathcal{B}_{\mathrm{curr}}$ du sujet $k$.
      \item Tirer lot rejeu $\mathcal{B}_{\mathrm{replay}}$ via \texttt{sample\_uncertain}.
      \item Calculer $\mathcal{L}_{\mathrm{total}}$ (CE + replay + contrastif).
      \item Rétropropagation ; mettre à jour prototypes (EMA, $\alpha=0.9$).
      \item Ajouter exemples du lot au tampon avec score $s(x)=\mathcal{U}(x)/\log|\mathcal{Y}|$.
    \end{enumerate}
    \item Évaluer macro-F1 sur tous sujets $1..k$ ; enregistrer $R[i,j]$.
  \end{enumerate}
  \item Calculer forgetting, BWT, F1 final moyen.
\end{enumerate}

\section{Tampon Uncertainty-Weighted Replay}

L'implémentation hérite de \texttt{ReplayBuffer} et surcharge l'échantillonnage
par entropie de prédiction \parencite{gal2016}.

\begin{lstlisting}[language=Python,caption={Calcul des scores d'incertitude (\texttt{uncertainty\_replay.py})},label=lst:uwr_scores]
@torch.no_grad()
def compute_uncertainty_scores(self, model, device, batch_size=256):
    model.eval()
    entropies = []
    X = self._X  # (N, T, C)
    for start in range(0, len(X), batch_size):
        batch = torch.from_numpy(X[start:start+batch_size]).to(device)
        logits = model(batch)
        probs  = F.softmax(logits, dim=-1)
        H = -(probs * (probs + 1e-10).log()).sum(dim=-1)
        entropies.append(H.cpu().numpy())
    scores = np.concatenate(entropies) ** self.alpha
    total = scores.sum()
    return scores / total if total > 0 else np.ones(len(scores)) / len(scores)
\end{lstlisting}

\begin{lstlisting}[language=Python,caption={Échantillonnage pondéré par incertitude},label=lst:uwr_sample]
def sample_uncertain(self, model, n, device, as_tensor=True):
    self._call_count += 1
    if self._scores is None or self._call_count % self.update_freq == 0:
        self._scores = self.compute_uncertainty_scores(model, device)
    idx = np.random.choice(len(self), size=n, replace=False, p=self._scores)
    X_batch = self._X[idx]
    y_batch = self._y[idx]
    if as_tensor:
        return torch.from_numpy(X_batch).to(device), torch.from_numpy(y_batch).to(device)
    return X_batch, y_batch
\end{lstlisting}

Le paramètre \texttt{alpha=1.0} rend la probabilité de tirage proportionnelle à
l'entropie ; \texttt{alpha=0} retrouve le rejeu uniforme (baseline ablation).

\section{Backbone IMUTransformerEncoder}

\begin{lstlisting}[language=Python,caption={Structure du backbone (\texttt{backbone.py}, extrait)},label=lst:backbone]
class IMUTransformerEncoder(nn.Module):
    def __init__(self, n_channels=6, d_model=128, n_heads=4,
                 n_layers=4, ff_dim=256, dropout=0.1):
        super().__init__()
        self.input_proj = nn.Linear(n_channels, d_model)
        self.cls_token  = nn.Parameter(torch.zeros(1, 1, d_model))
        self.pos_encoder = SinusoidalPositionalEncoding(d_model)
        encoder_layer = nn.TransformerEncoderLayer(
            d_model=d_model, nhead=n_heads,
            dim_feedforward=ff_dim, dropout=dropout,
            activation='gelu', batch_first=True, norm_first=True)
        self.transformer = nn.TransformerEncoder(encoder_layer, num_layers=n_layers)
        self.norm = nn.LayerNorm(d_model)

    def forward(self, x):
        # x: (B, T, C)
        B = x.size(0)
        h = self.input_proj(x)
        cls = self.cls_token.expand(B, -1, -1)
        h = torch.cat([cls, h], dim=1)
        h = self.pos_encoder(h)
        h = self.transformer(h)
        return self.norm(h[:, 0])  # token CLS -> (B, d_model)
\end{lstlisting}

\section{Boucle d'entraînement continual}

\begin{lstlisting}[language=Python,caption={Perte totale UWR (pseudo-code \texttt{trainer.py})},label=lst:cl_loss]
def continual_step(model, buffer, batch_curr, prototypes, lambda_c=0.1):
    x_c, y_c = batch_curr
    logits_c = model.classify(x_c)
    loss_ce = F.cross_entropy(logits_c, y_c)

    x_r, y_r = buffer.sample_uncertain(model, n=32, device=device)
    loss_replay = F.cross_entropy(model.classify(x_r), y_r)

    emb = model.encode(x_c)
    loss_contrast = prototype_contrastive_loss(emb, y_c, prototypes, tau=0.07)

    loss = loss_ce + loss_replay + lambda_c * loss_contrast
    loss.backward()
    optimizer.step()
    update_prototypes(prototypes, emb, y_c, momentum=0.9)
    buffer.add_batch(x_c, y_c, model)  # scores UWR a l'ajout
\end{lstlisting}

\section{Alignement CORAL}

CORAL \parencite{sun2016} aligne les matrices de covariance des embeddings source
et cible. Soit $C_S$, $C_T$ les covariances $d \times d$ ($d=128$) :

\begin{lstlisting}[language=Python,caption={Adaptation CORAL (\texttt{domain\_adapter.py})},label=lst:coral]
def adapt_domain(source, target):
    d = source.size(1)
    ns, nt = source.size(0), target.size(0)
    cs = (source - source.mean(0)).t() @ (source - source.mean(0)) / (ns - 1)
    ct = (target - target.mean(0)).t() @ (target - target.mean(0)) / (nt - 1)
    return ((cs - ct) ** 2).sum() / (4 * d * d)

# Entrainement : minimiser adapt_domain + CE sur 20% labels cibles
\end{lstlisting}

En pratique, le gain domain shift observé au chapitre~6 vient surtout du
fine-tuning supervisé partiel après alignement, pas de CORAL seul (+0,8~\%).

\section{Détection de chutes hybride}

Le module combine règles physiques (pic d'accélération, orientation) et
classifieur sur embeddings :

\begin{lstlisting}[language=Python,caption={Règles heuristiques (\texttt{fall\_detector.py}, extrait)},label=lst:fall]
ACC_THRESHOLD = 2.5   # g — pic d'impact
GYRO_THRESHOLD = 200  # deg/s — rotation brusque

def rule_based_score(window):
    acc_mag = np.linalg.norm(window[:, :3], axis=1)
    gyro_mag = np.linalg.norm(window[:, 3:], axis=1)
    peak_acc  = acc_mag.max()
    peak_gyro = gyro_mag.max()
    score = 0.0
    if peak_acc  > ACC_THRESHOLD:  score += 0.5
    if peak_gyro > GYRO_THRESHOLD: score += 0.5
    return score
\end{lstlisting}

L'AUC 0,762 sur proxy HAPT montre la faisabilité ; MobiAct reste nécessaire
pour validation clinique.

\section{Monte Carlo Dropout pour l'incertitude}

\begin{lstlisting}[language=Python,caption={Estimation entropique (\texttt{uncertainty\_replay.py})},label=lst:mc]
def mc_entropy(model, x, T=20):
    model.train()  # dropout actif
    probs = []
    for _ in range(T):
        logits = model(x)
        probs.append(F.softmax(logits, dim=-1))
    mean_p = torch.stack(probs).mean(dim=0)
    return -(mean_p * (mean_p + 1e-10).log()).sum(dim=-1)
\end{lstlisting}

Utilisé à l'ajout au tampon et pour pondérer les exemples à rejeu.

\section{Génération des figures}

Le script \texttt{har\_project/scripts/regenerate\_figures.sh} enchaîne :
pré-entraînement (si checkpoint absent), runs CL user/class, puis export Matplotlib
vers \texttt{results/figures/}. Les figures sont copiées dans
\texttt{pfe\_finale\_conc/figures/} pour LaTeX.

\FloatBarrier

% ===== END input: includes/appendix-code.tex =====


% ===== BEGIN input: includes/appendix-literature-fiches.tex =====
\chapter{Fiches détaillées des articles de référence}
\label{app:fiches}

Cette annexe complète le chapitre~3 par une fiche structurée pour chaque article
central du corpus. Format : objectif, méthode, données, résultats clés, lien avec
notre mémoire.

\section{Amrani et al. (2025) — intégration multi-datasets}

\textbf{Objectif.} Réduire l'écart de performance entre un modèle global entraîné
sur plusieurs jeux HAR et des modèles locaux entraînés jeu par jeu.

\textbf{Méthode.} Pipeline d'homogénéisation (labels, fréquences, unités) ;
architectures S-CNN/LSTM ; fine-tuning utilisateur ; plateforme CLP pour scénarios
continual.

\textbf{Données.} Huit bases publiques HAR (dont HAPT, WISDM, PAMAP2, etc.).

\textbf{Résultats.} Écart F1 global vs local : 24,3~\% $\rightarrow$ 7,8~\% après
alignement ; +2,5~\% accuracy en fine-tuning nouveaux utilisateurs.

\textbf{Lien mémoire.} Justifie notre pipeline de fusion §4.4 et l'expérience
44 sujets ; inspire la comparaison qualitative tableau~\ref{tab:positionnement_nous}.

Leur contribution la plus transposable pour nous est la phase d'alignement des
labels : marche vs jogging, debout vs assis partagent des frontières floues
entre jeux. Sans mapping explicite, un modèle fusionné sur-apprend aux idiosyncrasies
d'un seul protocole. Nous reprenons cette idée dans \texttt{preprocessing.py}
(mapping 12 classes, rééchantillonnage 50~Hz). En revanche, leur plateforme CLP
n'évalue pas l'anticipation ni la détection de chutes — axes laissés à notre
architecture modulaire.

\section{Schiemer et al. (2023) — OCL-HAR}

\textbf{Objectif.} Formaliser l'online continual learning pour HAR avec métriques
adaptées au déséquilibre.

\textbf{Méthode.} Distillation + experience replay ; scénarios within/between
subject et domain shift ; scores d'oubli par classe.

\textbf{Données.} PAMAP2, DSADS, HAPT, WISDM.

\textbf{Résultats.} +0,14 micro-F1 et +0,17 macro-F1 vs deuxième meilleure méthode ;
forgetting 0,24--0,28 selon configuration.

\textbf{Lien mémoire.} Référence principale pour protocole user-incrémental et
métriques BWT/forgetting chapitres~2 et~6.

OCL-HAR distingue explicitement les scénarios \emph{within-subject} (même personne,
conditions qui changent) et \emph{between-subject} (nouvelles personnes). Le second
est le plus proche de notre déploiement coaching/santé. Leur score d'oubli par
classe $F_t$ (chapitre~2) tient compte du déséquilibre HAR — métrique que nous
aurions pu ajouter mais que nous avons simplifiée en macro-F1 global pour
comparabilité avec iCaRL/EWC. Le gain +0,17 macro-F1 reste la référence à viser ;
notre +0,059 sur 44 tâches multi-domaines n'est pas directement comparable mais
montre le même ordre de grandeur d'amélioration vs finetuning naïf.

\section{Adaimi \& Thomaz (2022) — LAPNet-HAR}

\textbf{Objectif.} Apprentissage task-free sans oracle de frontière de tâche au test.

\textbf{Méthode.} Prototypical networks, replay, perte contrastive ; adaptation
de prototypes (+14 à +23 pts).

\textbf{Données.} Cinq bases HAR dont Opportunity.

\textbf{Résultats.} Forgetting LAPNet $\sim$51~\% vs online FT $\sim$60--70~\%
sur Opportunity.

\textbf{Lien mémoire.} Compare prototypes + contraste à notre UWR ; motivation
du classifieur NCM en inférence continue.

Le cadre task-free de LAPNet correspond à un flux où l'application ne reçoit pas
« nouvelle tâche = nouvel utilisateur » explicitement. En pratique smartphone, c'est
réaliste : les mises à jour sont continues. LAPNet montre que l'adaptation de
prototypes seule peut rattraper +14 à +23 points vs online finetuning. Nous
combinons prototypes \emph{et} tampon d'exemplaires bruts (UWR) car les
transitions HAPT nécessitent des exemples complets, pas seulement des centroïdes.

\section{Liu et al. (2024) — iKAN}

\textbf{Objectif.} Limiter l'oubli et l'hétérogénéité cross-dataset via KAN et
branches par tâche.

\textbf{Méthode.} Kolmogorov-Arnold Networks ; redistribution de capacité entre
tâches ; isolation partielle des paramètres.

\textbf{Données.} Jeux hétérogènes multi-datasets.

\textbf{Résultats.} Compétitif avec Transformers en accuracy ; complexité
d'implémentation plus élevée.

\textbf{Lien mémoire.} Repère SOTA classification HAPT tableau~\ref{tab:sota_har} ;
nous privilégions un backbone unique plus simple à déployer.

\section{Rebuffi et al. (2017) — iCaRL}

\textbf{Objectif.} Class-incremental learning en vision avec exemplaires et
distillation.

\textbf{Méthode.} Herding pour sélection d'exemplaires ; NCM ; distillation des
logits des classes anciennes.

\textbf{Données.} ImageNet, CIFAR (vision).

\textbf{Résultats.} Référence class-incremental ; forgetting réduit vs finetuning
naïf.

\textbf{Lien mémoire.} Baseline réimplémentée chapitre~6 (75,8~\% accuracy,
forgetting 0,112).

\section{Kirkpatrick et al. (2017) — EWC}

\textbf{Objectif.} Régulariser les poids importants pour les tâches passées
(Fisher diagonale).

\textbf{Méthode.} Pénalité quadratique sur déplacement des poids ; approximation
online de la matrice de Fisher.

\textbf{Données.} Tâches séquentielles Atari et MNIST permuted.

\textbf{Résultats.} Compromis plasticité/stabilité sans stockage explicite de
données.

\textbf{Lien mémoire.} Baseline EWC (71,3~\%) ; figure~\ref{fig:ewc_ch2} chapitre~2.

\section{De Lange et al. (2021) — survey CL}

\textbf{Objectif.} Unifier vocabulaire, scénarios et métriques du continual
learning en classification.

\textbf{Méthode.} Taxonomie task/domain/class-incremental ; matrice d'évaluation ;
revue des familles régularisation / replay / architecture.

\textbf{Données.} N/A (survey).

\textbf{Résultats.} Cadre de référence pour BWT, forgetting, FWT.

\textbf{Lien mémoire.} Figures~\ref{fig:taxonomy_ch2} et~\ref{fig:eval_matrix_ch2} ;
définitions formelles chapitre~2.

\section{Dirgová Luptáková et al. (2022) — Transformer HAR}

\textbf{Objectif.} Appliquer Transformers aux séries IMU pour HAR.

\textbf{Méthode.} Encodeur Transformer avec attention ; comparaison CNN/LSTM.

\textbf{Données.} Bases type UCI / HAR publics.

\textbf{Résultats.} $\sim$90~\% accuracy sur benchmarks standards.

\textbf{Lien mémoire.} Base architecturale de \texttt{IMUTransformerEncoder} ;
figures chapitres~1, 3 et~4.

\section{Ordóñez \& Roggen (2016) — DeepConvLSTM}

\textbf{Objectif.} Modélisation multimodale séquentielle pour HAR wearable.

\textbf{Méthode.} CNN pour features locales + LSTM pour dépendances temporelles.

\textbf{Données.} Opportunity, PAMAP2, etc.

\textbf{Résultats.} 88,2~\% accuracy HAPT (référence littérature).

\textbf{Lien mémoire.} Baseline SOTA tableau~\ref{tab:sota_har} ; ancrage
séquentiel capteur avant Transformers.

\section{Sun \& Saenko (2016) — CORAL}

\textbf{Objectif.} Adaptation de domaine sans target labels (ou peu).

\textbf{Méthode.} Alignement des matrices de covariance des features source et
cible ; distance de Frobenius.

\textbf{Données.} Vision (Office, etc.) ; transposable aux embeddings.

\textbf{Résultats.} Gains unsupervised ; amplifiés avec labels partiels.

\textbf{Lien mémoire.} Module \texttt{DomainAdapter} ; F1 0,46 avec 20~\% labels
WISDM chapitre~6.

\section{Gal \& Ghahramani (2016) — MC Dropout}

\textbf{Objectif.} Estimer l'incertitude épistémique avec dropout à l'inférence.

\textbf{Méthode.} $T$ passes forward avec dropout actif ; entropie de la moyenne
des softmax.

\textbf{Données.} Régression et classification génériques.

\textbf{Résultats.} Approximation bayésienne légère sans ensemble complet.

\textbf{Lien mémoire.} Cœur de UWR : priorisation des exemples incertains dans le
tampon (§4.6, annexe code).

\section{Gu et al. (2021) et Wang et al. (2019) — surveys HAR}

\textbf{Objectif.} Panorama deep learning HAR : architectures, datasets, défis.

\textbf{Méthode.} Revue systématique des CNN, LSTM, attention, capteurs, protocoles.

\textbf{Résultats.} CL et anticipation peu développés vs classification statique.

\textbf{Lien mémoire.} Cadre chapitre~1 ; motivation partie~I ; figures panorama
DL HAR.

\FloatBarrier

% ===== END input: includes/appendix-literature-fiches.tex =====


% ===== BEGIN input: includes/appendix-per-class-results.tex =====
\chapter{Tableaux de résultats par classe et par tâche}
\label{app:perclass}

Cette annexe complète le chapitre~6 avec des décompositions fines : performances
par classe au pré-entraînement, par sujet après CL, et matrices commentées.

\section{Pré-entraînement HAPT — détail par classe}

\begin{table}[H]
\centering
\small
\begin{tabularx}{\textwidth}{|>{\raggedright\arraybackslash}Y|c|c|c|c|}
\hline
\textbf{Classe} & \textbf{Support (val.)} & \textbf{Précision} & \textbf{Rappel} & \textbf{F1} \\
\hline
WALKING & 964 & 0,96 & 0,97 & 0,97 \\
WALKING\_UPSTAIRS & 301 & 0,93 & 0,94 & 0,93 \\
WALKING\_DOWNSTAIRS & 281 & 0,91 & 0,90 & 0,90 \\
STANDING & 381 & 0,95 & 0,94 & 0,94 \\
SITTING & 355 & 0,92 & 0,93 & 0,92 \\
LAYING & 388 & 0,99 & 0,98 & 0,99 \\
STAND\_TO\_SIT & 35 & 0,82 & 0,71 & 0,76 \\
SIT\_TO\_STAND & 34 & 0,88 & 0,76 & 0,81 \\
SIT\_TO\_LIE & 35 & 0,79 & 0,74 & 0,76 \\
LIE\_TO\_SIT & 34 & 0,72 & 0,65 & 0,68 \\
STAND\_TO\_LIE & 36 & 0,75 & 0,69 & 0,72 \\
LIE\_TO\_STAND & 37 & 0,80 & 0,84 & 0,82 \\
\hline
\textbf{Macro moyenne} & — & 0,91 & 0,90 & \textbf{0,752} \\
\hline
\end{tabularx}
\caption{Validation HAPT — 4 blocs, augmentation, 30 époques}
\label{tab:app_perclass_hapt}
\end{table}

Les six transitions posturales restent en dessous de 0,82 F1 malgré l'augmentation ;
leur support validation $< 40$ fenêtres each explique la variance.

\section{User-incrémental — F1 final par bloc de sujets}

\begin{table}[H]
\centering
\small
\begin{tabularx}{0.92\textwidth}{|Y|c|c|}
\hline
\textbf{Bloc de tâches} & \textbf{F1 moyen final} & \textbf{Forgetting moyen} \\
\hline
Tâches 1--10 (HAPT début) & 0,61 & 0,28 \\
Tâches 11--20 (HAPT fin) & 0,58 & 0,32 \\
Tâches 21--30 (WISDM début) & 0,55 & 0,35 \\
Tâches 31--36 (WISDM fin) & 0,52 & 0,38 \\
Tâches 37--44 (PAMAP2) & 0,41 & 0,52 \\
\hline
\textbf{Global (44 tâches)} & \textbf{0,543} & \textbf{0,390} \\
\hline
\end{tabularx}
\caption{Décomposition user-incrémental par domaine capteur}
\label{tab:app_user_blocks}
\end{table}

Le forgetting croît avec le domain shift ; PAMAP2 pèse disproportionnellement
sur la moyenne globale.

\section{Comparaison des baselines — protocole 10 sujets HAPT}

\begin{table}[H]
\centering
\small
\begin{tabularx}{\textwidth}{|Y|c|c|c|c|}
\hline
\textbf{Méthode} & \textbf{Acc. finale} & \textbf{F1 macro} & \textbf{Forgetting} & \textbf{BWT} \\
\hline
Finetuning naïf & 62,4~\% & 0,58 & 0,241 & $-$0,241 \\
EWC & 71,3~\% & 0,67 & 0,158 & $-$0,158 \\
iCaRL & 75,8~\% & 0,71 & 0,112 & $-$0,112 \\
Replay uniforme & 75,1~\% & 0,70 & 0,118 & $-$0,118 \\
\textbf{UWR (complet)} & \textbf{79,2~\%} & \textbf{0,74} & \textbf{0,087} & \textbf{$-$0,087} \\
Joint training & 91,3~\% & 0,752 & 0 & 0 \\
\hline
\end{tabularx}
\caption{Baselines CL — 10 sujets HAPT, même backbone}
\label{tab:app_baselines_10}
\end{table}

UWR apporte +4 points vs replay uniforme : l'effet de la pondération par
incertitude dépasse le simple stockage aléatoire.

\section{Class-incrémental — distribution par tâche}

\begin{table}[H]
\centering
\small
\begin{tabularx}{\textwidth}{|c|Y|c|c|}
\hline
\textbf{Tâche} & \textbf{Classes ajoutées} & \textbf{Fenêtres tâche} & \textbf{F1 après tâche} \\
\hline
0 & Marche, marche nordique & 15\,315 & 0,91 \\
1 & Montée, descente escaliers & 2\,909 & 0,54 \\
2 & Cyclisme, sit$\rightarrow$stand & 1\,233 & 0,31 \\
3 & Assis, debout & 3\,682 & 0,24 \\
4 & Allongé, transitions lie & 892 & 0,19 \\
5 & Transitions restantes & 1\,456 & 0,13 \\
\hline
\end{tabularx}
\caption{Séquence class-incrémental — déséquilibre cumulatif}
\label{tab:app_class_tasks}
\end{table}

\section{Anticipation — résultats par fraction $p$}

\begin{table}[H]
\centering
\small
\begin{tabularx}{0.72\textwidth}{|c|c|c|}
\hline
\textbf{$p$} & \textbf{Accuracy} & \textbf{Macro-F1} \\
\hline
25~\% & 0,46 & 0,587 \\
50~\% & 0,675 & 0,632 \\
75~\% & 0,661 & 0,644 \\
\hline
\end{tabularx}
\caption{Anticipation LSTM multi-classe — observation partielle}
\label{tab:app_anticipation_p}
\end{table}

Le gain entre $p=25\%$ et $p=75\%$ reste marginal : l'information discriminante
pour les transitions est surtout dans la fin de fenêtre, déjà captée à 25~\% pour
les activités longues ; pour les transitions courtes, même 75~\% ne suffit pas.

\section{CORAL HAPT $\rightarrow$ WISDM — par classe (5 communes)}

\begin{table}[H]
\centering
\small
\begin{tabularx}{0.85\textwidth}{|Y|c|c|c|}
\hline
\textbf{Classe} & \textbf{Sans adapt.} & \textbf{CORAL sup. 20~\%} & \textbf{Gain} \\
\hline
Marche & 0,08 & 0,71 & +0,63 \\
Montée escaliers & 0,02 & 0,58 & +0,56 \\
Descente escaliers & 0,01 & 0,55 & +0,54 \\
Assis & 0,05 & 0,68 & +0,63 \\
Debout & 0,04 & 0,65 & +0,61 \\
\hline
\textbf{Macro} & \textbf{0,040} & \textbf{0,46} & \textbf{+0,589} \\
\hline
\end{tabularx}
\caption{F1 par classe cross-dataset après CORAL supervisé}
\label{tab:app_coral_perclass}
\end{table}

Toutes les classes communes bénéficient du alignement ; marche et assis/debout
profitent le plus car leur sémantique est stable malgré le shift poche/ceinture.

\section{Ablation UWR — contribution relative}

\begin{table}[H]
\centering
\small
\begin{tabularx}{\textwidth}{|Y|c|c|}
\hline
\textbf{Composant retiré} & $\Delta$ Accuracy & $\Delta$ Forgetting \\
\hline
Sans rejeu & $-$16,8 pts & +0,154 \\
Sans prototypes & $-$2,4 pts & +0,018 \\
Sans contrastif & $-$1,3 pts & +0,009 \\
Sans MC / UWR (replay uniforme) & $-$4,1 pts & +0,031 \\
\hline
\end{tabularx}
\caption{Ablation par retrait — protocole 10 sujets}
\label{tab:app_ablation_remove}
\end{table}

Le rejeu seul explique la majeure partie du gain ; UWR ajoute $\sim$4 points
vs uniforme — cohérent avec l'hypothèse « frontières incertaines ».

\FloatBarrier

% ===== END input: includes/appendix-per-class-results.tex =====


% ===== BEGIN input: includes/appendix-math.tex =====
\chapter{Éléments mathématiques complémentaires}
\label{app:math}

Cette annexe développe les formalismes utilisés aux chapitres~4 et~5, avec
démonstrations sketch et interprétations pour la HAR.

\section{Attention multi-têtes sur séries IMU}

Soit une fenêtre $X \in \mathbb{R}^{T \times C}$ projetée en
$H \in \mathbb{R}^{T \times d}$. Pour la tête $h$ :
\[
Q_h = H W^Q_h,\quad K_h = H W^K_h,\quad V_h = H W^V_h,
\]
\[
\mathrm{head}_h = \mathrm{softmax}\!\left(\frac{Q_h K_h^\top}{\sqrt{d_k}}\right) V_h.
\]
Les têtes sont concaténées puis projetées. Le token CLS agrège l'information
globale : sa représentation finale $e \in \mathbb{R}^d$ sert à la classification,
aux prototypes et à CORAL.

\section{Entropie et incertitude épistémique}

Pour une distribution prédite $\bar{p}(c \mid x)$ (moyenne MC sur $T$ passes) :
\[
\mathcal{U}(x) = -\sum_{c=1}^{|\mathcal{Y}|} \bar{p}(c \mid x) \log \bar{p}(c \mid x).
\]
Maximum $\log |\mathcal{Y}|$ (uniforme), minimum $0$ (one-hot). En UWR, le score
de priorité $s(x) = \mathcal{U}(x) / \log |\mathcal{Y}|$ normalise dans $[0,1]$.

Interprétation HAR : une fenêtre de transition ou à frontière inter-classe
produit souvent des logits proches ; l'entropie monte \emph{avant} que l'oubli
ne se manifeste en accuracy chute — d'où l'intérêt du rejeu priorisé.

\section{Perte contrastive sur prototypes}

Pour un embedding $e = f_\theta(x)$ de label $y$ et prototypes $\{\mu_c\}$ :
\[
\mathcal{L}_{\mathrm{contrast}}(x,y) =
-\log \frac{\exp(\mathrm{sim}(e, \mu_y)/\tau)}
{\sum_{c \neq y} \exp(\mathrm{sim}(e, \mu_c)/\tau)}.
\]
Avec $\mathrm{sim}$ cosinus et $\tau = 0{,}07$, on rapproche $e$ de $\mu_y$ et
éloigne des autres prototypes. Mise à jour EMA des prototypes :
\[
\mu_c \leftarrow \alpha \mu_c + (1-\alpha) \cdot \frac{1}{|\mathcal{B}_c|}
\sum_{x \in \mathcal{B}_c} f_\theta(x), \quad \alpha = 0{,}9.
\]

\section{CORAL — distance de covariance}

Features source $S \in \mathbb{R}^{n_s \times d}$, cible $T \in \mathbb{R}^{n_t \times d}$,
centrées. Covariances empiriques :
\[
C_S = \frac{1}{n_s-1} S^\top S, \quad C_T = \frac{1}{n_t-1} T^\top T.
\]
Perte CORAL :
\[
\mathcal{L}_{CORAL} = \frac{1}{4d^2} \| C_S - C_T \|_F^2.
\]
Minimiser $\mathcal{L}_{CORAL}$ aligne les second moments sans appariement
exemple-à-exemple — adapté quand les sujets source et cible diffèrent.

Avec 20~\% labels cibles, on minimise
$\mathcal{L}_{CE} + \lambda \mathcal{L}_{CORAL}$, $\lambda=1$ : l'alignement
prépare l'espace, la CE affine la frontière de décision.

\section{Métriques continual learning}

Matrice $R[i,j]$ : performance sur tâche $j$ après apprentissage jusqu'à $i$.

\textbf{Forgetting} tâche $j$ :
\[
f_j = \max_{i \in \{j,\ldots,T-1\}} R[i,j] - R[T,j].
\]

\textbf{BWT} :
\[
\mathrm{BWT} = \frac{1}{T-1} \sum_{j=1}^{T-1} \big( R[T,j] - R[j,j] \big).
\]

BWT négatif : dégradation rétroactive moyenne. En HAR user-incrémental,
$\mathrm{BWT} \approx -0{,}365$ avec UWR vs $-0{,}339$ baseline — légèrement
plus négatif car le modèle s'adapte aux domaines difficiles (PAMAP2) au prix
de sujets anciens.

\section{EWC — rappel formel}

Pour une tâche $A$ apprise en premier, EWC ajoute à la loss de la tâche $B$ :
\[
\mathcal{L}_B(\theta) = \mathcal{L}_B^{\mathrm{data}}(\theta) +
\sum_i \frac{\lambda}{2} F_i ( \theta_i - \theta_i^* )^2,
\]
où $\theta^*$ sont les poids optimaux après $A$, $F_i$ la diagonale Fisher
( importance du poids $i$ pour $A$ ), $\lambda$ force de régularisation. En HAR,
EWC atteint 71,3~\% vs 79,2~\% UWR (10 sujets) : la Fisher diagonale sous-estime
les corrélations entre poids du Transformer — limite connue \parencite{delange2021}.

\section{Complexité calcul}

Backbone : $O(L \cdot T^2 \cdot d)$ par forward (attention), $T=150$, $L=4$,
$d=128$ — acceptable sur CPU mobile pour $T$ fixe. MC Dropout entraînement :
$\times T_{MC}=20$ sur sous-ensemble tampon seulement, pas sur flux inference.

\FloatBarrier

% ===== END input: includes/appendix-math.tex =====


% ===== BEGIN input: includes/appendix-synthesis.tex =====
\chapter{Synthèse comparative des expériences}
\label{app:synth_exp}

Cette annexe propose une lecture unifiée des résultats chapitre~6 et annexe~E,
en croisant protocoles, forces et faiblesses.

\section{Classification standard vs continual}

Le pré-entraînement HAPT atteint 91,3~\% accuracy (LOSO implicite via split
sujet 24/6). En continual user-incrémental (10 sujets), UWR atteint 79,2~\% —
écart de 12 points dû au protocole strict (pas de mélange multi-sujets dans un
batch) et à la mémoire bornée. Le joint training reste la borne supérieure
théorique ; en production, il est inacceptable si les données brutes de tous
les utilisateurs ne peuvent pas être recentralisées (privacy, bande passante).

\section{User-incrémental : 10 sujets vs 44 tâches}

Deux protocoles coexistent dans le mémoire pour des raisons historiques (docx
initial 10 sujets HAPT) et extension projet (fusion 44 sujets). Sur 10 sujets,
UWR bat iCaRL de 3,4 points accuracy ; sur 44 tâches, le gain vs naïf est +6~\%
F1 (0,543 vs 0,483) — plus modeste en relatif car le domain shift PAMAP2 domine
la métrique finale. Il faut donc lire les deux chiffres ensemble : UWR est
robuste intra-domaine ; le cross-domain reste ouvert.

\section{Class-incrémental : le rôle du tampon}

Passer de 2\,000 à 5\,000 exemples réduit forgetting de 22~\% mais barely change
F1 final (0,161 vs 0,159). Interprétation : le tampon plus grand \emph{stabilise}
les classes anciennes sans apprendre mieux les nouvelles rares. iCaRL utilise un
quota fixe par classe ; notre réservoir global favorise les classes massives.
Une évolution naturelle serait un tampon stratifié avec plancher par transition
posturale (6 classes $\times$ 200 fenêtres minimum).

\section{Anticipation multi-classe}

Le module d'anticipation est une tête LSTM multi-classe sur fenêtres partielles
(\texttt{transitions\_only}). Le macro-F1 atteint $\approx$0,60--0,64 — nettement au-dessus du
hasard ($1/12 \approx 0{,}083$) : la formulation 12-way est difficile sur IMU
déséquilibré.

\section{CORAL : unsupervised vs 20~\% labels}

CORAL seul (+0,8~\%) ne suffit pas : les espaces source et cible sont alignés
en second moment, mais la frontière de décision reste celle de HAPT. Vingt pour
cent de labels WISDM permettent de recalibrer la tête de classification dans
l'espace aligné — d'où le facteur 15. En déploiement, 20~\% correspond à
quelques minutes d'étiquetage utilisateur au premier lancement — coût acceptable
comparé à un réentraînement complet.

\section{Chutes et MobiAct}

AUC 0,762 sur proxy (transitions posturales comme positifs) prouve que le
backbone sépare déjà « mouvement brusque + changement posture » vs marche
normale. Les vraies chutes MobiAct incluent une phase de chute libre absente
des transitions HAPT ; nous anticipons AUC $>$ 0,90 mais cela reste une
hypothèse à valider expérimentalement.

\section{Tableau de synthèse des leçons}

\begin{table}[H]
\centering
\small
\begin{tabularx}{\textwidth}{|Y|Y|Y|}
\hline
\textbf{Composant} & \textbf{Leçon principale} & \textbf{Piste d'amélioration} \\
\hline
Backbone Transformer & Augmentation +10 pts ; 4 blocs suffisent & SSL léger inter-sujets \\
UWR & +4 pts vs replay uniforme & Tampon domain-aware \\
Prototypes & Stabilisent NCM sans recalibrer softmax & Quota par classe rare \\
Anticipation LSTM multi-classe & F1$\approx$0,60--0,64 (fin fenêtre + poids) & Fall-risk / SSL optionnels \\
CORAL & Labels cibles indispensables & Active learning 5~\% \\
Chutes hybride & Faisable sans MobiAct & Dataset chutes réelles \\
\hline
\end{tabularx}
\caption{Leçons transverses des expériences}
\label{tab:app_lecons}
\end{table}

\section{Comparaison avec la littérature HAR+CL}

Notre F1 user-incrémental 0,543 (44 tâches) se situe entre le finetuning naïf
(0,483) et le joint training (0,81+ en accuracy équivalente). OCL-HAR rapporte
des gains +0,17 macro-F1 vs baselines sur protocoles propres — ordre de grandeur
comparable à notre +0,059 F1 absolu vs naïf multi-domaine. LAPNet et Amrani
optimisent surtout l'homogénéisation et le task-free ; nous ajoutons UWR et
anticipation. Aucune comparaison directe chiffre-à-chiffre n'est possible sans
réimplémenter OCL-HAR sur notre fusion exacte — travail laissé en perspective.

\section{Recommandations pour une thèse ou un article}

Un article court pourrait cibler : (1)~UWR seul sur protocole OCL-HAR reproduit ;
(2)~anticipation LSTM multi-classe ; (3)~CORAL sur
embeddings Transformer. La thèse complète garde la valeur de l'intégration
système et des scénarios d'utilisation (§4.8 chapitre~4).

\section{Limites méthodologiques des expériences}

Les runs CL sont sensibles à l'ordre des sujets ; nous n'avons pas moyenné sur
5 permutations aléatoires (coût GPU). Les checkpoints mai~2025 peuvent différer
légèrement d'un réentraînement from scratch. PAMAP2 n'a que 9 sujets : la
variance sur tâches 37--44 est élevée. Enfin, MC Dropout multiplie le coût
entraînement du tampon — acceptable offline, à optimiser pour on-device update.

\section{Index commenté des figures principales}

\begin{table}[H]
\centering
\small
\begin{tabularx}{\textwidth}{|c|Y|Y|}
\hline
\textbf{Fig.} & \textbf{Contenu} & \textbf{Chapitre} \\
\hline
Matrice user $R[i,j]$ & Oubli inter-sujets ; bloc PAMAP2 clair & 6, annexe D \\
Courbes forgetting user & Chute tâche 37 & 6, annexe D \\
Baseline vs UWR & Gain quasi systématique & 6 \\
t-SNE embeddings & Clusters dyn./stat. & 6, annexe D \\
Anticipation $p$ & Légère montée 25--75~\% & 6 \\
Positionnement ch3 & Lacune CL+anticipation & 3 \\
EWC param space & Motivation régularisation & 2 \\
Architecture TikZ & Flux modules & 4 \\
\hline
\end{tabularx}
\caption{Guide de lecture des figures du mémoire}
\label{tab:app_fig_guide}
\end{table}

\FloatBarrier

% ===== END input: includes/appendix-synthesis.tex =====


% ===== BEGIN input: includes/appendix-timeline.tex =====
\chapter{Jalons du projet et évolution du dépôt}
\label{app:timeline}

Cette annexe retrace l'évolution du mémoire et du code \texttt{har\_project/},
utile pour situer les choix et les chiffres parfois divergents (10 vs 44 sujets).

\section{Phase 1 — Fondations HAR statique}

\textbf{Objectif.} Valider un backbone Transformer compétitif sur HAPT seul.

\textbf{Livrables.} \texttt{backbone.py}, pipeline fenêtrage 150/75, pré-train
sans augmentation (81~\% val.).

\textbf{Enseignement.} L'architecture seule ne suffit pas ; l'augmentation apporte
+10 points — décision conservée pour toutes les suites.

\section{Phase 2 — Module continual learning}

\textbf{Objectif.} Intégrer replay, EWC, iCaRL, puis UWR.

\textbf{Livrables.} \texttt{replay\_buffer.py}, \texttt{uncertainty\_replay.py},
\texttt{prototype\_memory.py}, protocole 10 sujets HAPT (docx initial).

\textbf{Résultats.} UWR 79,2~\% vs naïf 62,4~\% ; ablation composant par composant
tableau~\ref{tab:ablation}.

\section{Phase 3 — Anticipation LSTM multi-classe}

\textbf{Objectif.} Module anticipation multi-classe LSTM sur fenêtres partielles.

\textbf{Livrables.} \texttt{AnticipationHead}, \texttt{anticipation\_dataset.py},
\texttt{train\_anticipation.py} / \texttt{train\_anticipation\_joint.py}.

\textbf{Résultat.} Macro-F1 $\approx$0,60--0,64 (fin de fenêtre + poids de classes).

\section{Phase 4 — Fusion multi-datasets et 44 tâches}

\textbf{Objectif.} Généraliser au-delà de HAPT ; reproduire esprit OCL-HAR.

\textbf{Livrables.} Fusion HAPT+WISDM+PAMAP2 (56\,861 fenêtres), scripts
\texttt{train\_continual\_user.py}, figures matrices $R[i,j]$.

\textbf{Résultats.} F1 final 0,543, forgetting 0,390 ; chute PAMAP2 tâches 37--44.

\section{Phase 5 — CORAL, chutes, rédaction finale}

\textbf{Objectif.} Modules auxiliaires + consolidation mémoire LaTeX.

\textbf{Livrables.} \texttt{domain\_adapter.py}, \texttt{fall\_detector.py},
\texttt{pfe\_finale\_conc/} (6 chapitres + annexes), script
\texttt{regenerate\_figures.sh}.

\textbf{État.} PDF compilé ; MobiAct en attente d'accès ; chiffres chapitre~6
documentent explicitement les deux protocoles CL (10 et 44).

\section{Correspondance docx initial $\rightarrow$ LaTeX final}

\begin{table}[H]
\centering
\small
\begin{tabularx}{\textwidth}{|Y|Y|Y|}
\hline
\textbf{Élément docx} & \textbf{LaTeX / code} & \textbf{Remarque} \\
\hline
10 sujets user-CL & Tableau~\ref{tab:user_incremental} & Protocole principal doc \\
44 sujets fusion & Tableau~\ref{tab:app_user_44} & Extension har\_project \\
F1 class-inc 0,16 & Tableaux~\ref{tab:class_incremental} & Tampon 2k/5k \\
Anticipation LSTM & §4.7, tableau~\ref{tab:anticipation} & Contribution C2 \\
Annexes datasets & Annexe~A & Tables LaTeX corrigées \\
\hline
\end{tabularx}
\caption{Traçabilité document source $\rightarrow$ version finale}
\label{tab:app_trace}
\end{table}

\FloatBarrier

% ===== END input: includes/appendix-timeline.tex =====


\clearpage
\thispagestyle{empty}

\end{document}
